\documentclass[11pt]{article}
\usepackage[T1]{fontenc}
\usepackage[utf8]{inputenc}
\usepackage[english]{babel}
\usepackage{amsmath,amssymb,mathtools,array,booktabs,pdflscape,adjustbox,diagbox}
\usepackage{geometry}
\usepackage{microtype}
\geometry{a4paper,margin=1.45cm}
\setlength{\parindent}{1.25em}
\setlength{\parskip}{0pt}
\newcommand{\widebar}[1]{\overline{#1}}
\newcommand{\A}{\Delta}
\newcommand{\D}{\Delta}
\allowdisplaybreaks
\setlength{\emergencystretch}{10em}
\sloppy
\hbadness=10000
\vbadness=10000
\hfuzz=2pt
\vfuzz=10pt
\raggedbottom
\newcommand{\R}{\varrho}
\newcommand{\hata}{\widehat{a}}
\newcommand{\hatv}{\widehat{v}}
\newcommand{\modu}[1]{\;\operatorname{mod} #1}
\newcommand{\mcell}[1]{\ensuremath{\begin{gathered}#1\end{gathered}}}
\providecommand{\pair}[2]{(#1\,|\,#2)}
\providecommand{\eps}{\varepsilon}
\providecommand{\rhoidx}{\varrho}
\begin{document}

\begin{center}
{\Large\bfseries 1. On the Formation of the form System of the Ternary\\ Biquadratic form}\par
\vspace{1.5em}
Reports of the Physical-Medical Society in Erlangen 39 (1907), pp. 176--179
\end{center}

\vspace{4em}
\begin{center}
(Extract from the author's dissertation.)
\end{center}

Work by Gordan, Maisano, and Pascal\footnote{P. Gordan, ``Über the volle Formensystem of the ternären biquadratischen form'' $f=x_1^3x_2+x_2^3x_3+x_3^3x_1$ (Math. Annalen vol. XVII, pp. 217--233, 1880); G. Maisano, 1. ``Sistemi completi dei primi cinque gradi della forma ternaria biquadratica e degl' invarianti, covarianti e contravarianti di sesto grado'' (Giorn. di Battaglini XIX). 2. ``Sui covarianti indipendenti di $6^\circ$ grado nei coefficienti della forma biquadratica ternaria'' (Rend. Circ. Mat. di Palermo I, 1887); E. Pascal, ``Contributo alla teoria della forma ternaria biquadratica e delle sue varie decomposizioni in fattori'' (Memoria premiata dalla R. Accademia delle scienze fisiche e matematiche di Napoli, 1905).} deals with the form system of the ternary biquadratic form. Gordan sets up the complete form system, consisting of 54 formations, for the special automorphic form
\[
f=x_1^3x_2+x_2^3x_3+x_3^3x_1,
\]
on the basis of principles similar to those which he had given for form systems in the binary domain.

In Maisano's work, for the general biquadratic form, the forms up to and including the fifth order\footnote{By ``order'' the dimension in the coefficients is to be understood, and by ``degree'' the dimension in the variables.} are set up, together with some invariants, covariants, and contravariants of higher order, according to the method used by Gordan in volume I of the \emph{Mathematische Annalen} for the ternary cubic form. Pascal, using Maisano's results, is chiefly concerned with the question of the decomposition of the biquadratic form into factors\footnote{In his second short note Maisano attempts to prove a linear dependence among the three covariants of order 6 and degree 6. In contrast to this not quite complete proof, Pascal believes he has proved the linear independence of the three covariants of order 6, as well as that of the three invariants of order 9. That, however, a linear relation does in fact exist between the three covariants on the one hand and the three invariants on the other, emerged in the course of my computations; in Pascal's notation the explicit formulas are
\[
0=20\Omega_1+6\Omega_2-3\Omega_3+4fC_1-12fC_2-4A\cdot\Delta,
\]
\[
0=4A^3-15AB+30C-90D+9E.
\]}.

\textbf{The aim of my investigations is to set up the form system for the general ternary biquadratic form; to begin with, only the main foundations are given, and a so-called ``relatively complete system''\footnote{For the terminology, see Gordan--Kerschensteiner, \emph{Vorlesungen über Invariantentheorie}, vol. II, p. 227.} is established.}

The basic idea for forming systems of ternary forms is the same as in the binary domain. Starting from a first relatively complete system -- the system of forms taken over from the binary case -- one passes, according to a definite law, to systems with ever higher modulus, until either the system of a modulus -- taking the modulus as ground form -- becomes finite and known, or until a modulus can be reduced to forms having invariants as factors. By transvection of the relatively complete system with the system of the modulus, in the first case one obtains the absolutely complete system, while in the second case the relatively complete and absolutely complete systems coincide. By Hilbert's general proof the finiteness of form systems, this procedure must necessarily come to an end.

In our case, the modulus $(abc)$ of the first relatively complete system can be led back to the moduli
\[
\Delta=(abc)^2a_x^2b_x^2c_x^2\quad\text{and}\quad \nu=(abu)^4;
\]
from this the relatively complete system modulo $\nu$ is readily obtained. As the sequence of moduli we now choose the forms
\[
\nu;\qquad \nu^{(\nu)}=(\nu\nu_1x)^4=s_x^4,
\]
\[
\nu^{(s)}=(ss'u)^4\ldots,
\]
and adjoin, as further moduli, two quadratic forms occurring in the formation of the relatively complete system modulo $s$, namely $u_\sigma^2$ and $t_x^2$. One can then show that the modulus following $s$, namely $(ss'u)^4$, is reducible to the modulus $(\varrho,t)$. Since, however, the simultaneous system of two quadratic forms is finite and known, the termination described above has thereby been reached. As the relatively complete system modulo $(\varrho,t)$ one obtains 331 formations. ---

I add a few details concerning the method used.

The fundamental process for producing forms, the folding process, can be defined in the ternary domain as follows:

If in a symbolic product
\[
s_x^m t_x^n u_\sigma^\mu u_\tau^\nu
\]
one replaces the pairs of factors
\[
\begin{array}{c|c|c|c|c}
 & s_xt_x & u_\sigma u_\tau & s_xu_\sigma\ \text{or}\ s_xu_\tau & t_xu_\tau\ \text{or}\ t_xu_\sigma\\
\text{respectively by} & (stu) & (\sigma\tau x) & s_\sigma\ \text{or}\ s_\tau & t_\tau\ \text{or}\ t_\sigma\\
\text{folding} & \mathrm{I} & \mathrm{II} & \mathrm{III} & \mathrm{IV},
\end{array}
\]
then the forms so arising have been obtained from the original form by folding\footnote{Gordan, Math. Annalen, vol. XVII, p. 219.}.

Here one may still always put $s_\sigma=0;\ t_\tau=0$. For the relation among the individual foldings, the following theorem holds:

\textbf{The foldings I and II are fundamental foldings, from which, independently of the order of composition, the foldings III and IV can be composed. In other words: to form all forms arising from a given expression by folding, one need only apply foldings I and II}\footnote{The theorem has an exception in the case where $n$ and $\mu$, respectively $m$ and $\nu$, vanish simultaneously; the ``form sequence'' then reduces to the diagonal term.}.

By a ``form sequence'' we mean an initial form together with all forms arising from it by folding into itself. By the theorem, a form sequence $s_x^m t_x^n u_\sigma^\mu u_\tau^\nu$ can be arranged in a rectangular scheme. In proceeding, one obtains:
\begin{enumerate}
\item by moving one column to the right, all forms arising by folding I from the adjacent forms;
\item by moving one row downward, all forms arising by folding II from the forms standing above, and hence all forms arising by folding.
\end{enumerate}

Under these assumptions, the theorems on reducers can be stated in their most general form, under the definition
\begin{center}
``A reducer is a reducible form sequence,''
\end{center}
as follows:

\begin{center}
\textbf{If the initial form of a form sequence is reducible because one of its terms has arisen by folding with a reducer, and if the final form of the form sequence has arisen from this same term by folding, then the whole form sequence is reducible.}
\end{center}

Further reduction methods to be mentioned are:
\begin{enumerate}
\item the so-called ``double reduction,'' that is, the reduction of a form in two ways in order to obtain a relation between the higher forms;
\item ``folding with decomposed forms,'' which partly has the character of reduction by reducers and partly that of ``double reduction.''
\end{enumerate}
\clearpage

\begin{center}
{\Large\bfseries 2. On the Formation of the form System of the Ternary\\ Biquadratic form}\par
\vspace{1.5em}
Journal for the reine and angewandte Mathematik 134 (1908), pp. 23--90 and two tables
\end{center}

\vspace{3em}

\begin{center}
\textbf{Introduction.}
\end{center}

Work by Gordan, Maisano, and Pascal\footnote{A short extract from the Introduction and Chapter I appeared in the \emph{Sitzungsberichte of the physikalisch-medizinischen Sozietät Erlangen 1907}, pp. 176--179.}\footnote{P. Gordan, on the complete form system of the ternary biquadratic form
\[
f=x_1^3x_2+x_2^3x_3+x_3^3x_1.
\]
(\emph{Math. Annalen}, vol. XVII (1880), pp. 217--233.)
G. Maisano, 1) \emph{Sistemi completi dei primi cinque gradi della forma ternaria biquadratica e degl' invarianti, covarianti e contravarianti di sesto grado}. (Giorn. di Battaglini XIX.)
2) \emph{Sui covarianti indipendenti di $6^\circ$ grado nei coefficienti della forma biquadratica ternaria}. (Rend. Circ. Mat. di Palermo I, 1887.)
E. Pascal, \emph{Contributo alla teoria della forma ternaria biquadratica e delle sue varie decomposizioni in fattori}. (Memoria premiata dalla R. Accademia delle scienze fisiche e matematiche di Napoli. 1905.)} deals with the form system of the ternary biquadratic form. Gordan sets up the complete form system, consisting of 54 formations, of the special automorphic form
\[
f=x_1^3x_2+x_2^3x_3+x_3^3x_1
\]
on the basis of principles similar to those that he had given for form systems in the binary domain.

In Maisano's work, for the general biquadratic form, the forms up to and including the fifth order\footnote{By ``order'' the dimension in the coefficients is to be understood, and by ``degree'' the dimension in the variables.} are set up, together with some invariants, covariants, and contravariants of higher order, according to the method used by Gordan in volume I of the \emph{Mathematische Annalen} for the ternary cubic form. Pascal, using Maisano's results, is chiefly concerned with the question of the decomposition of the biquadratic form into factors\footnote{In his second short note, Maisano attempts to prove a linear dependence of the three covariants of order 6 and degree 6. In contrast to this not quite complete proof, Pascal believes he has proved the linear independence of the three covariants of order 6, and likewise that of the three invariants of order 9. That, however, a linear relation does in fact exist between the three covariants on the one hand and the three invariants on the other is shown by formulas (13), § 11, and (22), § 17 note, of the present work, where these relations are given explicitly. According to a communication from Pascal, the contradiction is explained by a numerical error in the expression for his contravariant $p$ (called $\varrho$ here), p. 46 of his paper.}.

\emph{The aim of the following investigations is to set up the form system for the general ternary biquadratic form; in this work, namely, only the principal foundations are given, and a so-called ``relatively complete system''\footnote{For the terminology see § 3a.} is established.}

The work is closely connected with Gordan's work; however, the principles that were stated there only quite generally first had to be worked out in detail. With the aid of a theorem on the connection among the foldings and by introducing the ``form sequence'' (§ 1), the reduction theorems are sharply formulated and completed (§ 3), while the recursive establishment of special series expansions (§ 2) gives the computational means for actually carrying out the reductions.

The basic idea for forming systems of ternary forms is the same as in the binary domain. Starting from a first relatively complete system -- the system of forms taken over from the binary case -- one passes, according to a definite law, to systems with ever higher modulus, until either the system of a modulus -- the modulus taken as ground form -- becomes finite and known, or until a modulus can be reduced to forms having invariants as factors. By transvection of the relatively complete system with the system of the modulus, in the first case the absolutely complete system arises, while in the second case the relatively complete and absolutely complete systems coincide. By Hilbert's general proof the finiteness of form systems, this procedure must necessarily come to an end.

In our case the modulus $(abc)$ of the first relatively complete system (§ 4) is reduced to the moduli
\[
\Delta=(abc)^2a_x^2b_x^2c_x^2\quad\text{and}\quad \nu=(abu)^4
\]
(§ 5), and then the relatively complete system modulo $\nu$ is found (§ 6). These results are already given in Gordan's work for the general biquadratic form; our manner of deriving them, however, is more easily transferred to forms of higher degree.

As the sequence of moduli we now choose the forms (§ 7)
\[
\nu,\qquad \nu(\nu)=(\nu\nu_1x)^4=s_x^4,\qquad \nu(s)=(ss'u)^4,\ldots .
\]

In forming the relatively complete system modulo $s$, two quadratic forms occurring in the system, $u_\varrho^2$ and $t_x^2$, are adjoined as moduli (§§ 9 and 10), and accordingly the relatively complete system modulo $(s,\varrho,t)$ is formed. It is then shown (§ 17) that the modulus $(ss'u)^4$ is reducible to the moduli $\varrho$ and $t$. The next higher system thereby passes over into a relatively complete system modulo $(\varrho,t)$. Since, however, the simultaneous system of two quadratic forms is finite and known, and in the general case consists of 20 formations\footnote{Cf. Clebsch--Lindemann, \emph{Vorlesungen über Geometrie}. (Vol. I, p. 288 ff.) System of two cogredient forms. Gordan, ``Über Büschel of Kegelschnitten.'' (\emph{Math. Ann.} vol. XIX, p. 530.) System of two contragredient forms.}, the termination described above is reached with the establishment of the relatively complete system modulo $(\varrho,t)$. The transvection of this system with the system of $(\varrho,t)$ in order to form the absolutely complete system is reserved.

As the relatively complete system modulo $(\varrho,t)$ one finds 331 formations, which in the appended table are ordered according to their degree in the variables $x$ and $u$.

Finally we give an overview of the symbols introduced:
\begin{align*}
f&=a_x^4,\\
\theta&=\theta_x^4u_\vartheta^2=(abu)^2a_x^2b_x^2,\\
K&=K_x^6u_k^3=(a\theta u)u_\vartheta^2a_x^3\theta_x^3,\\
j&=u_j^6=(a\theta u)^4u_\vartheta^2,\\
\Delta&=\Delta_x^6=a_\vartheta^2a_x^2\theta_x^4=(abc)^2a_x^2b_x^2c_x^2,\\
N&=N_x^8u_\eta=(a\Delta u)a_x^3\Delta_x^5,\\
\nu&=u_\nu^4=(abu)^4,\\
H&=H_x^2u_\eta^4=(\nu\nu_1x)^2u_\nu^2u_{\nu_1}^2,\\
L&=L_x^3u_l^6=(\nu\eta x)H_x^2u_\nu^3u_\eta^3,\\
g&=g_x^6=(\nu\eta x)^4H_x^2,\\
\sigma&=u_\sigma^6=H_\nu^2u_\nu^2u_\eta^4,\\
s&=s_x^4=(\nu\nu_1x)^4,\\
Z&=Z_x^4u_\zeta^2=(ss'u)^2s_x^2s_x'^2,\\
\varrho&=u_\varrho^2=\theta_\nu^4u_\vartheta^2,\\
t&=t_x^2=a_\eta^4H_x^2,\\
i&=a_\nu^4,\\
J&=s_\nu^4.
\end{align*}

\begin{center}
{\Large\bfseries Chapter I. \quad General Theorems on Ternary Forms.}
\end{center}

\begin{center}
\textbf{§ 1. \quad The folding process. \quad form sequences.}
\end{center}

The fundamental process for producing forms is the folding process, which in the ternary domain can be defined as follows. Let there be given a symbolic product
\[
s_x^m t_x^n u_\sigma^\mu u_\tau^\nu .
\]
If one replaces the pairs of factors
\[
s_xt_x\qquad u_\sigma u_\tau\qquad s_xu_\sigma\ \text{or}\ s_xu_\tau
\qquad t_xu_\tau\ \text{or}\ t_xu_\sigma
\]
respectively by
\[
(stu)\qquad (\sigma\tau x)\qquad s_\sigma\ \text{or}\ s_\tau
\qquad t_\tau\ \text{or}\ t_\sigma,
\]
folding
\[
\text{I}\qquad\text{II}\qquad\text{III}\qquad\text{IV},
\]
then the resulting forms have arisen from the original one by folding\footnote{Gordan, \emph{Math. Annalen} vol. XVII, p. 219.}.

Here the expression $s_x^m t_x^n u_\sigma^\mu u_\tau^\nu$ can be regarded either as an actual product of two forms $S\cdot T$, or as a single form with several rows of symbols:
\[
s_x^m t_x^n u_\sigma^\mu u_\tau^\nu=A_x^{m+n}u_x^{\mu+\nu}.
\]
In the first case we speak of the folding of the form $S$ with $T$, in the second case of the ``folding of the form $A$ in itself.'' For brevity we always assume in what follows, as indeed is the case for all forms later to be considered, that
\[
\tag{a}
s_\sigma^\lambda=0,\qquad t_\tau^\lambda=0
\]
for any $\lambda$ and $x$ distinct from 0, and in conjunction with any other foldings. This says that the form $s_x^m t_y^n u_\sigma^\mu u_\tau^\nu$ is a so-called ``normal form,'' that is, it vanishes under application of the $\Omega$-process, both with respect to the variables $x$ and $u$ and with respect to the variables $y$ and $v$.

Thus condition (a) represents no restriction, but can always be achieved\footnote{Cf. Gordan, \emph{Math. Annalen} vol. V, p. 104.}.

For the connection among the individual foldings the following theorem holds:

\emph{Theorem I. Foldings I and II are fundamental foldings from which, independently of the order of composition, foldings III and IV can be composed. In other words: in order to form all forms arising by folding from a given expression, one has to apply only foldings I and II.}

Proof: By the identity theorem, respectively the product theorem for matrices, taking (a) into account, one obtains
\[
(\widehat{st}\sigma x)=-t_\sigma,\qquad
(\widehat{\sigma\tau}su)=-s_\tau,
\]
\[
(\widehat{st}\tau x)=s_\tau,\qquad
(\widehat{\sigma\tau}tu)=t_\sigma,
\]
\[
(stu)(\sigma\tau x)=s_\tau+t_\sigma-u_x\,s_\tau t_\sigma .
\]
Thus by composing foldings I and II one obtains foldings III and IV, and this holds both when applying folding II to folding I, i.e. to the factors $(stu)u_\sigma u_\tau$, and when applying folding I to folding II, to the factors $(\sigma\tau x)s_xt_x$.

By a \emph{form sequence}\footnote{Cf. Clebsch, ``Über eine Fundamentalaufgabe of the Invariantentheorie'' (Abh. of the Gött. Ges. d. Wiss. vol. XVII): § 17 and end of § 18. The ``form sequence'' differs from the ``properly reduced equivalent system'' introduced there by the principle of ordering according to higher forms.} we mean an initial form together with all forms that have arisen from it by folding in itself, and we denote the form sequence by the initial form. A ``higher form'' is to mean here any form with more folding in itself, while among the equally entitled foldings $s_\tau$ and $t_\sigma$ one has to be normalized as higher. According to the law of formation of the form sequence it follows that:

1) Each form of the form sequence is linear in the symbols of the initial form.

2) The form sequence is uniquely determined for initial forms with two rows each of contragredient symbols; for initial forms with more rows of symbols it is to be uniquely normalized by distinguishing special foldings among those connected by the identity theorem.

By Theorem I, a form sequence $s_x^m t_x^n u_\sigma^\mu u_\tau^\nu$ ($m>n;\ \mu>\nu$) can be arranged according to the following rectangular scheme:
\begingroup\small
\[
\begin{array}{c|c|c|c}
s_x^m t_x^n u_\sigma^\mu u_\tau^\nu
& (stu)
& (stu)^2
& \cdots\\[0.3em]
(\sigma\tau x)
& t_\sigma,\ s_\tau
& t_\sigma(stu),\ s_\tau(stu)
& \cdots\\[0.3em]
(\sigma\tau x)^2
& t_\sigma(\sigma\tau x),\ s_\tau(\sigma\tau x)
& t_\sigma^2,\ t_\sigma s_\tau,\ s_\tau^2
& \cdots\\[0.3em]
\vdots&\vdots&\vdots&\ddots\\[0.3em]
(\sigma\tau x)^\nu
& t_\sigma(\sigma\tau x)^{\nu-1},\ s_\tau(\sigma\tau x)^{\nu-1}
& t_\sigma^2(\sigma\tau x)^{\nu-2},\ t_\sigma s_\tau(\sigma\tau x)^{\nu-2},\ s_\tau^2(\sigma\tau x)^{\nu-2}
& \cdots\\[0.3em]
\vdots
& t_\sigma(\sigma\tau x)^\nu
& t_\sigma^2(\sigma\tau x)^{\nu-1},\ t_\sigma s_\tau(\sigma\tau x)^{\nu-1}
& \cdots
\end{array}
\]
\endgroup
and the correspondingly continued lower part\footnote{Here, and similarly in what follows, the lower part marked with an asterisk is to be set at the correspondingly marked place at the upper right of the scheme.}
\begingroup\small
\[
\begin{array}{c|c|c}
(stu)^n & \cdots & \cdots\\[0.3em]
t_\sigma(stu)^{n-1},\ s_\tau(stu)^{n-1}
& s_\tau(stu)^n & \cdots\\[0.3em]
t_\sigma^2(stu)^{n-2},\ t_\sigma s_\tau(stu)^{n-2},\ s_\tau^2(stu)^{n-2}
& t_\sigma s_\tau(stu)^{n-1},\ s_\tau^2(stu)^{n-1}
& s_\tau^2(stu)^n .
\end{array}
\]
\endgroup

Here one obtains, when one moves

1) one column to the right, all forms arising by folding I from the adjacent forms,

2) one row downward, all forms arising by folding II from the forms above,

and hence all forms arising by folding.

An arbitrary form of the total form sequence, $t_\sigma^\kappa s_\tau^\lambda(stu)^\mu$ or $t_\sigma^\kappa s_\tau^\lambda(\sigma\tau x)^\nu$, forms the beginning of a new form sequence, consisting of all those forms of the rectangular scheme with $t_\sigma^\kappa s_\tau^\lambda(stu)^\mu$ or $t_\sigma^\kappa s_\tau^\lambda(\sigma\tau x)^\nu$ as left upper corner, which have the factor $t_\sigma^\kappa s_\tau^\lambda$.

Supplementary remark. Theorem I has an exception in the case where the numbers $n$ and $\mu$ (respectively $m$ and $\nu$) become equal to zero, so that foldings I, II, and IV no longer exist, while folding III (respectively IV) still does. In this case we designate as the form sequence the diagonal member of the scheme:
\[
\begin{array}{ccccc}
s_x^m u_\tau^\nu &&&& t_x^n u_\sigma^\mu\\
& s_\tau && t_\sigma &\\
& s_\tau^2 && t_\sigma^2 &\\
&&\ddots&&\\
&&s_\tau^\nu && t_\sigma^\mu .
\end{array}
\]
In accordance with the absence of foldings I and II, the series expansion according to polars of the form sequence in this case always reduces to the initial member; however, the theorems on reducers remain valid.

\begin{center}
\textbf{§ 2. \quad Series expansions according to polars of the form sequence.}
\end{center}

For forms with $n$ variables two kinds of series expansions have been set up explicitly\footnote{Cf. Gordan, ``Über Kombinanten,'' §§ 2 and 5 (\emph{Math. Ann.} vol. V).}:

1) for forms with one row each of contragredient variables, $s_x^m u_\tau^\nu$, a series progressing by powers of $u_x$, whose coefficients are ``normal forms'';

2) for forms with two rows of cogredient variables, $s_x^m t_x^n$, an expansion according to polars of the elementary covariants (polars of the form sequence $s_x^m t_x^n$), which in the ternary case reads as follows\footnote{Cf. Study, \emph{Methoden to the Theorie of the ternären Formen} (Teubner 1889) II. § 7. In distinction to the derivation there, ours gives the method for the rapid computational determination of the numerical coefficients $C_{pr\varrho}$. The Clebsch--Study expansion, upon specialization a), a'), would read
\[
(stu)^\mu u_\tau^{\nu-\kappa}v_\tau^\kappa s_x^m t_x^{\,n-\lambda}(tuv)^\lambda
=\sum_{p,r,\varrho} C_{pr\varrho}
\bigl[s_\tau^\varrho(stu)^{\mu+r-\varrho}\bigr]_{\widehat{uv}^{\lambda-r+\varrho}v^{\nu-\varrho},\,u^p,\,(vw=\widehat{xw})},
\]
and therefore does not permit the simplification that occurs already in the course of our calculation, once one knows that all terms with the factor $u_x$ are to be omitted.}:
\[
\tag{I}
s_x^m t_x^{n-\lambda}t_y^\lambda
=\sum_{\varrho}
\frac{\binom{m}{\varrho}\binom{\lambda}{\varrho}}
{\binom{m+n-\varrho+1}{\varrho}}
\bigl[(st\widehat{xy})^\varrho s_x^{m-\varrho}t_x^{n-\varrho}\bigr]_{y^{\lambda-\varrho}}.
\]

We combine the two expansions by setting up, for forms with two rows each of contragredient variables,
\[
s_x^m t_x^{n-\lambda}t_y^\lambda  u_\sigma^\mu u_\tau^{\nu-\kappa}v_\tau^\kappa,
\]
expansions according to powers of $u_x$, whose coefficients are composite polars of the form sequence $s_x^m t_x^n u_\sigma^\mu u_\tau^\nu$, taken with respect to the variables $x$ and $u$.

It is enough to set up the series for two contragredient rows of symbols (respectively variables), since, by combining each two rows of symbols (variables) into a single new row, forms with more rows of symbols (variables) can be reduced to this case.

In order to ensure that all forms of the form sequence contain only two rows of symbols, respectively variables, we specialize:
\[
\begin{array}{ll}
\text{a) } \sigma=\widehat{st}, & \text{a') } y=\widehat{uv},\\[0.3em]
\text{b) } s=\widehat{\sigma\tau}, & \text{b') } v=\widehat{xy},
\end{array}
\]
and carry out the expansion for the specialization a), a').

By direct transfer of expansion I one obtains composite polars for forms $(stu)^\mu s_x^m t_x^{n-\lambda}t_y^\lambda$, whereas for forms $(stu)^\mu u_\tau^{\nu-\kappa}v_\tau^\kappa s_x^m t_x^{n-\lambda}t_y^\lambda$, instead of composite polars, terms of such a kind arise whose evaluation makes necessary the introduction of ever new auxiliary variables that are to be set equal only at the end, thereby complicating the calculation unnecessarily.

We therefore take an indirect route, by representing the composite polars of all forms of the form sequence, that is, expressions
\[
\bigl[s_\tau^\varrho(stu)^{\mu-\varrho+r}\bigr]_{y^\lambda}^{v^\kappa},
\]
by the initial members of these polars, and by inversion of the resulting recursive system of equations obtaining the desired expansion according to polars.

By the transfer principle, for the $\lambda$-th polar of a form $s_x^m t_x^n:[s_x^m t_x^n]_{y^\lambda}$ ($\lambda\le n$) the following expansion holds (the case $\lambda>n$ is reduced to the first by the expansion of $[s_y^m t_y^n]_{x^{m+n-\lambda}}$)\footnote{Gordan, \emph{The Resultante binärer Formen} (Chap. I, § 5). Rend. Circ. Mat. di Palermo XXII. 1906.}:
\[
\bigl[s_x^m t_x^n\bigr]_{y^\lambda}
=\sum_r (-1)^r\,
\frac{\binom{m}{r}\binom{\lambda}{r}}{\binom{m+n}{r}}\,
(st\widehat{xy})^r s_x^{m-r}t_x^{n-\lambda}t_y^{\lambda-r},
\]
and correspondingly for the $\kappa$-th polar of $u_\sigma^\mu u_\tau^\nu:[u_\sigma^\mu u_\tau^\nu]_{v^\kappa}$
\[
\bigl[u_\sigma^\mu u_\tau^\nu\bigr]_{v^\kappa}
=\sum_\varrho (-1)^\varrho\,
\frac{\binom{\mu}{\varrho}\binom{\kappa}{\varrho}}{\binom{\mu+\nu}{\varrho}}\,
(\sigma\tau\widehat{uv})^\varrho u_\sigma^{\mu-\varrho}u_\tau^{\nu-\kappa}v_\tau^{\kappa-\varrho}.
\]

By multiplication it follows, with introduction of the specialization a) and a'),
\[
\begin{aligned}
\bigl[s_x^m t_x^n(stu)^\mu u_\tau^\nu\bigr]_{\widehat{uv}^{\lambda}}^{v^\kappa}
={}&
\sum_{r,\varrho}c_{r\varrho}\,
(st\widehat{uv})^r(st\widehat{uv})^\varrho
s_x^{m-r}t_x^{n-\lambda}(tuv)^{\lambda-r}(stu)^{\mu-\varrho}
u_\tau^{\nu-\kappa}v_\tau^{\kappa-\varrho},
\end{aligned}
\]
and from this, by expansion according to powers of $u_x$ with the first member distinguished ($\sum'_{r,\varrho}$ means that the member $r=0$, $\varrho=0$ is to be omitted) and taking account of (a) on p. 26,
\[
\tag{II}
\begin{aligned}
&s_x^m t_x^{n-\lambda}(tuv)^\lambda(stu)^\mu u_\tau^{\nu-\kappa}v_\tau^\kappa\\
&\quad =
\bigl[s_x^m t_x^n(stu)^\mu u_\tau^\nu\bigr]_{\widehat{uv}^{\lambda}}^{v^\kappa}\\
&\qquad
-\sum_{r,\varrho}' c_{r\varrho}\,s_\tau^\varrho(stu)^{\mu-\varrho+r}
u_\tau^{\nu-\kappa}v_\tau^{\kappa-\varrho}
s_x^{m-r}t_x^{n-\lambda}(tuv)^{\lambda-r+\varrho}v_x^r\\
&\qquad
+u_x\sum_{\varrho}\sum_{r=1}^{\lambda} c_{r\varrho}\,
s_\tau^\varrho(stu)^{\mu-\varrho+r-1}(stv)
u_\tau^{\nu-\kappa}v_\tau^{\kappa-\varrho}
s_x^{m-r}t_x^{n-\lambda}(tuv)^{\lambda-r+\varrho}v_x^{r-1}\\
&\qquad
\vdots\\
&\qquad
-(-1)^\lambda  u_x^\lambda\sum_\varrho c_{\lambda\varrho}\,
s_\tau^\varrho(stu)^{\mu-\varrho}(stv)^\lambda
u_\tau^{\nu-\kappa}v_\tau^{\kappa-\varrho}
s_x^{m-\lambda}t_x^{n-\lambda}(tuv)^\varrho,
\end{aligned}
\]
where
\[
c_{r\varrho}=(-1)^{r+\varrho}\cdot
\frac{\binom{m}{r}\binom{\lambda}{r}}{\binom{m+n}{r}}\cdot
\frac{\binom{\mu}{\varrho}\binom{\kappa}{\varrho}}{\binom{\mu+\nu}{\varrho}}
\qquad
\begin{array}{l}
\text{for }\lambda\le n,\\
\kappa\le\nu.
\end{array}
\]
and correspondingly for the remaining combinations of the numbers $\lambda,n;\kappa,\nu$. Thus the expression $s_x^m t_x^{n-\lambda}(tuv)^\lambda(stu)^\mu u_\tau^{\nu-\kappa}v_\tau^\kappa$ is reduced to a composite polar, and, by applying the identity
\[
(stv)u_\tau=(stu)v_\tau-s_\tau(tuv),
\]
to a sum of analogously formed expressions corresponding to higher forms of the form sequence.

By iteration one arrives at the expansion
\[
\tag{III}
s_x^m t_x^{n-\lambda}(tuv)^\lambda(stu)^\mu u_\tau^{\nu-\kappa}v_\tau^\kappa
=
\sum_{p,r,\varrho}u_x^p\,C_{pr\varrho}\cdot
\bigl[s_\tau^\varrho(stu)^{\mu-\varrho+r}\bigr]_{\widehat{uv}^{\lambda-r+\varrho}v^{\kappa-\varrho+p},\,v_x^{r-p}}.
\]
The numerical coefficients $C_{pr\varrho}$ are computed uniquely and recursively from the coefficients $c_{r\varrho}$, $c'_{r\varrho}$ of the system of equations arising by iteration of (II) (cf. the example on p. 42). Analogous expansions arise for the other specializations.

\begin{center}
\textbf{§ 3. \quad Reduction theorems.}
\end{center}

The known theorems on the reduction of forms and form systems shall now be brought into connection with §§ 1 and 2, for which some definitions taken over from the theory of binary forms are necessary.

a) We call a system of forms a \emph{relatively complete system} modulo a prescribed sequence of forms if it has the property that all formations arising by folding an arbitrary product of these forms can be expressed as integral rational functions of the forms of the system, up to an additive term consisting of expressions that have arisen by folding the system forms with the system of the modulus\footnote{Gordan--Kerschensteiner, \emph{Vorlesungen über Invariantentheorie}. Vol. II. p. 227.}.

b) We call a form \emph{reducible} when it can be expressed by forms having invariants as factors or by ``higher forms,'' that is, by higher forms of the total form sequence to which the form belongs, or by forms containing the symbols of the modulus in higher order.

The theorems on reducers\footnote{Gordan, \emph{Math. Annalen} vol. XVII, p. 222.} can now be stated in their most general form as follows:

Definition: \emph{A reducer is a reducible form sequence.}

\emph{Theorem II. If the initial form of a form sequence is reducible because one of its members has arisen by folding with a reducer (has a reducer as factor), and if the final form of the form sequence has arisen from exactly this member by folding, then the total form sequence is reducible}\footnote{The simplification that results from Theorem II can be seen in the following example. For the special form $f=x_1^3x_2+x_2^3x_3+x_3^3x_1$, $a_y^2a_x^2u_y^2$ is a reducer. It follows by Theorem II that the form sequence
\[
\theta_\nu(\vartheta\nu x)\theta_x^3u_\vartheta u_\nu^2
=
a_\nu^2(abu)-a_\nu b_\nu(aba),
\]
is reduced, since $\theta_\nu^4=a_\nu^2b_\nu^2(abu)^2$ has arisen from the member $a_\nu^2(abu)$ by folding. In Gordan, loc. cit., p. 231, by contrast, the reduction of the forms
\[
\theta_\nu(\vartheta\nu x),\quad \theta_\nu(\vartheta\nu x)^2,\quad
\theta_\nu^2,\quad \theta_\nu^2(\vartheta\nu x),\quad
\theta_\nu^2(\vartheta\nu x)^2
\]
is carried out individually.}.

Proof: The form sequence arises from the initial member by folding in itself, hence by higher folding with the reducer on the one hand, and by folding the reducer in itself on the other hand. In both cases, by § 2, all forms of the form sequence can be represented as polars (transvections) over the form sequence of the reducer.

The inference from the initial form to the total form sequence no longer holds for the remaining reduction methods. These are:

1) so-called ``double reduction.''

By this we mean the reduction, in two different ways, of an expression that has two mutually independent reducers as factors, whereby relations arise between the higher forms to which one has reduced.

By systematic application of ``double reduction'' one must, on the one hand, obtain all relations\footnote{Thus, for example, the following were foand by ``double reduction'': the reduction of the form $a_\eta^4$ (§ 10), the only form included superfluously in Maisano's system; and also the relations mentioned in the note to the Introduction between the three covariants and the three invariants.}; on the other hand, if ``double reduction'' applied to different expressions supplies no new relations, one has both a criterion for the independence of the higher forms and a check on the calculation.

2) \emph{Folding with decomposable forms.}

By ``decomposable forms'' are to be understood -- with a slight generalization of the usual concept -- forms that can be expressed by products of forms of lower degree and by ``higher'' forms. Products of forms pass into products under one folding; likewise products of nonlinear forms under twofold folding in the specializations a') and b') (§ 2), according to the product theorem:
\[
a_yb_y a_xb_x=\frac12a_y^2b_x^2+\frac12b_y^2a_x^2-\frac12(ab\widehat{yx})^2.
\]
First polars (transvections) and certain second polars of decomposable forms are therefore again decomposable forms. It follows that:

A member of a onefold (respectively twofold) transvection over a decomposable form, arising by onefold (respectively twofold) folding with a decomposable form, itself becomes a decomposable form:

a) if the next-higher forms in the form sequence of the decomposable form decompose or become reducible, or also if, under onefold folding, the specializations $y=\widehat{uv}$ or $v=\widehat{xy}$ occur;

b) if the remaining onefold foldings lead to higher forms (cf. § 12, B. $H_k$ and $H_k(k\eta x)$).

Such a member of a transvection can also decompose:

c) if the remaining onefold foldings lead to lower forms that decompose independently of the folding with the decomposable form, that is, that can be expressed by products and by the form to be reduced, although in each case one must first verify by calculation that the numerical coefficient of the member concerned does not vanish identically. This is nothing other than ``double reduction'' of the lower form (cf. § 23, C. $H_q(\theta H u)(\theta s u)$).

Decomposable forms arise by all onefold foldings with functional determinants\footnote{cf. Gordan, loc. cit. § 3.}, and moreover by certain higher foldings. One has
\[
(st\widehat{yx})s_y=\frac12\,t\,s_y^2-\frac12\,s\,t_y^2+\frac12(st\widehat{yx})^2,
\]
\[
(st\widehat{yx})^2s_y^2=\frac13\,t\,s_y^3-\frac13\,s\,t_y^3+s_y(st\widehat{yx})^2-\frac13(st\widehat{yx})^3.
\]
Analogous formulas hold for the dualistic forms
\[
(\sigma\tau\widehat{\nu u})v_\sigma
\quad\text{and}\quad
(\sigma\tau\widehat{\nu u})v_\sigma^2.
\]

From the theorems of this section the following rule follows for setting up all irreducible forms that arise from the folding of $S$ with $T$:

Beginning with the lowest foldings, proceed along any previously fixed path (for example row by row, or symmetrically with respect to the diagonal member) in the formation of the form sequence $ST$ until one reaches a form that has a reducer as factor. The form sequence defined by this form is to be omitted as soon as the final form satisfies the condition stated in Theorem II. After all reducible form sequences have been eliminated, all remaining reducible forms, and likewise all decomposable forms, are to be removed by applying double reduction. Places in the total form sequence that are covered twice by independently reducible form sequences give rise to the reduction of higher forms (cf. § 11, D).


\clearpage
\noindent\textbf{Theorem III.}
\emph{The forms taken over from the binary domain, that is, those forms which arise from the system forms of the corresponding binary form by bordering according to Clebsch's transfer principle, form a relatively complete system modulo $(abc)$.}

\noindent\textbf{Proof.}
To a binary relation asserting that the forms $Q'_1,\ldots,Q'_n$ form a complete system of a binary ground form,
\[
Q'-F(Q')=0=\sum_\lambda\bigl\{(ab)c_x+(bc)a_x+(ca)b_x\bigr\}\varphi_\lambda^{*},
\]
there corresponds, by bordering, the ternary relation
\[
Q-F(Q_i)=\sum_\lambda  u_x\cdot(abc)\varphi_\lambda  .
\]
This is precisely the defining equation for a relatively complete system modulo $(abc)$. For the biquadratic form, the system of the forms taken over consists of the following forms:
\[
\begin{gathered}
f=a_x^4,\qquad
\theta=\theta_x^4u_\vartheta^2=(abu)^2a_x^2b_x^2,\qquad
K=K_x^6u_k^3=(a\theta u)u_\vartheta^2a_x^3\theta_x^3,\\
j=u_j^6=(a\theta u)^4u_\vartheta^2,\qquad
\nu=u_\nu^4=(abu)^4,
\end{gathered}
\]
among which the first four form a relatively complete system modulo $((abc),\nu)$.

\begin{center}
\textbf{§ 5. \quad Reduction of the modulus $(abc)$ to the moduli $\A$ and $\nu$.}
\end{center}

The modulus $(abc)$, given only by a bracket factor, is to be reduced, in agreement with the definition (§ 3, a), to forms of the system. In other words, the form sequence $(abc)$ is to be normalized uniquely (cf. § 1).

\[
a_s^2a_xu_x=\frac13 a_s^3u_x=\frac13 S\cdot u_x. \tag{2}
\]
\[
\left.
\begin{aligned}
(a\A u)^2&=a_s-\frac{S}{6}u_x^2,\\
(a\A u)^3&=0.
\end{aligned}
\right\} \tag{3}
\]
\[
\left.
\begin{aligned}
\theta(s)=(ss,x)^2&=2a_t+\frac13 S\cdot\theta-\frac13 Tu_x^2,\\
\theta(t)=(tt,x)^2&=-\frac13 S\cdot a_t+\frac23 T\cdot a_s+\frac1{18}S^2\cdot\theta-\frac1{18}u_x^2\cdot S\cdot T.
\end{aligned}
\right\} \tag{4}
\]

Sequence of moduli: $(abc)$; $\A$; $s$; $t$ (the system of the modulus $t$ consists of the single form $t$).\footnote{Gordan--Kerschensteiner, loc. cit., p. 134.}

It will be shown that the forms
\[
\A=(abc)^2a_x^2b_x^2c_x^2,\qquad \nu=(abu)^4
\]
suffice as moduli, that is, that the forms
\[
\begin{gathered}
\A=(abc)^2a_x^2b_x^2c_x^2,\qquad
a_\nu=(abc)(bcu)^3a_x^3,\\
a_\nu^2=(abc)^2(bcu)^2u_x^2,\qquad
i=a_\nu^4=(abc)^4
\end{gathered}
\]
define the form sequence $(abc)$. In the form sequence $a_\nu$, the form $a_\nu^3$ is missing according to the identity
\[
a_\nu^3=(abc)^3a_x(bcu)=\frac13 u_x\cdot(abc)^4=\frac13 i\cdot u_x.
\]

For reducing a modulus to a higher one, according to the definition, we have to reduce the most general foldings, or all special foldings coming into consideration, with the modulus to foldings with the higher modulus.

In the present case the most general foldings arise from the expressions
\[
(abc)a_x^3b_y^3c_z^3,\qquad
(abc)a_x^2b_y^2c_z^2a_zb_xc_x
\]
(taken over from the cubic forms),
\[
(abc)a_xb_yc_za_x^2b_x^2c_x^2
\]
(taken over from the quadratic forms), and from the polarisation of these expressions.

For the calculation of these expressions by the polars of $\A$ and of the form sequence $a_\nu$, respectively of their initial terms, we take the indirect route, as in § 2. We expand the polar terms
\[
(abc)^2a_x^2b_y^2c_z^2,\qquad
a_\nu(\nu xy)(\nu yz)(\nu zx)a_xa_ya_z
=(abc)(bc\widehat{x}u)(bc\widehat{y}z)(bc\widehat{z}x)a_xa_ya_z\quad\text{etc.}
\]
as linear aggregates of expressions with the symbolic factor $(abc)$ and obtain, by inverting the system of equations, the desired relations, as well as a relation between the polars taken according to different combinations of the variables. The identity theorem and the product theorem for determinants are used in the evaluation.

The system of equations is, for $(abc)a_x^3b_y^3c_z^3$:
\begingroup\small
\[
\begin{aligned}
1)\quad
\sum_3 a_\nu(\nu yz)^3a_x^3
&=6(abc)a_x^3b_y^3c_z^3
-6\sum_3(abc)a_x^3b_y^2b_zc_z^2c_y^{*},\\[0.4em]
2)\quad
3(abc)^2a_x^2b_y^2c_z^2\cdot(xyz)
-\frac12\sum_3 a_\nu^2(\nu yz)^2\nu_x^2(xyz)
&=2\sum_3(abc)a_x^3b_y^2b_zc_z^2c_y\\
&\quad-4\sum_3(abc)a_xa_ya_zb_y^2b_zc_z^2c_x,\\[0.4em]
3)\quad
a_\nu(\nu xy)(\nu yz)(\nu zx)a_xa_ya_z
&=2\sum_3(abc)a_xa_ya_zb_y^2b_zc_z^2c_x,\\[0.4em]
4)\quad
\sum_3 a_\nu(\nu yz)^3u_x^3
-\frac32\sum_3 a_\nu^2(\nu yz)^2\nu_x^2(xyz)
+\frac16 i\cdot(xyz)^3
&=6\sum_3(abc)a_xa_ya_zb_y^2b_zc_z^2c_x.
\end{aligned}
\]
\endgroup

It follows for
\[
(abc)a_x^3b_y^3c_z^3,
\]
and by analogous calculation for
\[
(abc)a_x^2b_y^2c_z^2a_zb_xc_x,\qquad
(abc)a_xb_yc_za_z^2b_x^2c_x^2:
\]
\[
\begin{aligned}
(abc)a_x^3b_y^3c_z^3
&=\frac32(abc)^2a_x^2b_y^2c_z^2(xyz)
+\frac32 a_\nu(\nu xy)(\nu yz)(\nu zx)a_xa_ya_z\\
&\quad-\frac1{12}i\cdot(xyz)^3,\\[0.4em]
(\text{IV.})\quad
(abc)a_x^2b_y^2c_z^2a_zb_xc_x
&=\frac23(abc)^2a_x^2b_y^2c_z^2c_x\cdot(xyz)\\
&\quad+\frac13a_\nu(\nu xy)(\nu yz)(\nu zx)a_x^3
-\frac16 a_\nu^2(\nu xy)(\nu zx)a_x^2\cdot(xyz),\\[0.4em]
(abc)a_xb_yc_za_z^2b_x^2c_x^2
&=\frac16(abc)^2a_x^2b_z^2c_z^2\cdot(xyz).
\end{aligned}
\]

We have therefore obtained a relatively complete system modulo $(\A,\nu)$, consisting of the four forms $f,\theta,K,j$. It follows from this that all transvections of $f$ over $\theta$ not taken over from the binary domain, as well as the transvection $(a\theta u)^2$ reducible in the binary case to $(ab)^4$, can be expressed by the symbols $\A$ and $\nu$.\footnote{$\sum_3$ refers to the sum of three terms obtained from the initial term by cyclic interchange of the variables. The equations also hold individually for each term of the sum.}

We obtain the reduction formulas fundamental for what follows by specializing expansion IV, or more shortly by direct calculation (interchanging the symbols), for $a_\vartheta(a\theta u)^2$, $a_\vartheta^2((a\theta u)^2$, $(a\theta u)^2$, according to the following ansatz:
\[
\begin{aligned}
a_\vartheta(a\theta u)^2
&=(abc)(bcu)(abu)^2-\frac13(abc)(bcu)\{(bcu)a_x-u_x(abc)\}^2,\\
a_\vartheta^2(a\theta u)^2
&=(abc)^2(abu)^2-\frac13(abc)^2\{(bcu)a_x-u_x(abc)\}^2,
\end{aligned}
\]
for $((a\theta u)^2)^{*}$:
\[
\begin{aligned}
(abu)a_y^3b_x^3&=\frac32u_\vartheta(\vartheta yx)\theta_y^2+\frac14(\nu yx)^3,\\
((abu)(acu))b_x^3&=\frac32u_\vartheta(\vartheta\widehat{a}ux)(\theta au)^2+\frac14(\nu\widehat{a}ux)^3.
\end{aligned}
\]
\[
\left.
\begin{aligned}
(a\theta u)^2&=\frac16\nu\cdot f-\frac23u_x\cdot a_\nu
+\frac23u_x^2\cdot a_\nu^2-\frac{i}{18}u_x^4,\\
a_\vartheta&=\frac13u_x\cdot\A,\\
a_\vartheta(a\theta u)&=0,\\
a_\vartheta^2&=\A,\\
(a\theta u)^3&=0,\\
a_\vartheta(a\theta u)^2&=-\frac56a_\nu+\frac76u_x\cdot a_\nu^2-\frac{i}{9}u_x^3,\\
a_\vartheta^2(a\theta u)&=0,\\
(a\theta u)^4&=j,\\
a_\vartheta(a\theta u)^3&=0,\\
a_\vartheta^2(a\theta u)^2&=\frac23a_\nu^2-\frac{i}{9}u_x^2.
\end{aligned}
\right\} \tag{1}
\]

We add the formula, likewise obtained from the reduction of the modulus $(abc)$:
\[
a_\nu^3a_xu_\nu=\frac13 i\cdot u_x. \tag{2}
\]

\emph{From formulas (1.) and (2.) and from the formulas formed analogously for the ground form $\nu$, all later reduction formulas are derived by polarisation}, except for the relations for the decomposed forms. From formulas (1.) we see:

The form sequence leads:
\[
\begin{array}{rcl}
a_\vartheta(a\theta u)^2 &\text{to}& \text{symbols }\nu\text{ alone},\\[0.2em]
a_\vartheta &\text{to}& \text{symbols }\A\text{ and }\nu,\\[0.2em]
(a\theta u)^2 &\text{to}& \text{symbols }j,\A\text{ and }\nu,\\[0.2em]
(a\theta u)^2\text{ under folding I} &\text{to}& \text{symbols }j\text{ and }\nu,\\[0.2em]
(a\theta u)^2\text{ under folding II} &\text{to}& \text{symbols }\A\text{ and }\nu.
\end{array}
\]

We can therefore replace:
\begin{enumerate}
\item modulo $(\A,\nu)$: foldings with $j$ are replaced by corresponding foldings with $(a\theta u)^2$ or $(a\theta u)^3$;
\item modulo $(\nu)$: foldings with $\A$ are replaced by corresponding foldings with $a_\vartheta$ or $a_\vartheta(a\theta u)$, or also by special foldings with $(a\theta u)^2$ (which lead only to a single folding I).
\end{enumerate}
In particular:
\[
\begin{aligned}
(a\theta u)^2\theta_y^2\theta_x^2&=\frac13(jyx)^2\pmod{\nu},&
(a\theta u)^2u_y\theta_y&=-\frac16(jyx)^2\pmod{\nu},\\
(a\theta u)^2\nu_y^2&=\frac13(\A u\nu)^2\pmod{\nu},&
(a\theta u)(a\theta\nu)u_\vartheta\nu_\vartheta&=-\frac16(\A u\nu)^2\pmod{\nu},\\
(a\theta u)^2u_\nu^2\vartheta_\nu^2&=\frac13(\A u\nu)^2\A_\nu\pmod{\nu}.
\end{aligned}
\]

\begin{center}
\textbf{§ 6. \quad Reduction of the modulus $(\A,\nu)$ to the modulus $(\nu)$.}
\end{center}

The reduction formulas of the preceding paragraph give us the means to set up the relatively complete system modulo $\nu$ directly; we shall see that to the system of the transferred forms one has only to add the forms $\A$ and $(a\A u)=N$ (analogously as for the cubic form, and indeed valid in general).

For the reduction of the modulus $\A$ we consider all special foldings coming into consideration, that is, the foldings of the relatively complete system modulo $(\A,\nu)$ with the system of $\A$, and begin with the forms lowest in the coefficients, the foldings of $f$ with $\A$.

For the reduction of the forms $(a\A u)^2$, $(a\A u)^3$, $(a\A u)^4$, we replace the symbols $\A$ by lower forms of the form sequence according to § 5; that is, we develop the expressions
\[
a_\vartheta b_\vartheta(b\theta u)^2,\qquad
a_\vartheta(a\theta u)b_\vartheta(b\theta u)^2,\qquad
a_\vartheta(a\theta u)b_\vartheta(b\theta u)^3
\]
1) by polars of the form sequence $a_\vartheta$ (symbols $\A$ and $\nu$),

2) by polars of the form sequence $b_\vartheta(b\theta u)^2$ (symbols $\nu$),

and obtain the reduction formulas (calculation below):
\begin{align}
\text{a)}\quad (a\A u)^2
&=\frac12(\vartheta\nu x)^2-\frac25 f\cdot a_\nu^2
-\frac45u_x\cdot\theta_\nu(\vartheta\nu x)^2 \notag\\
&\quad+u_x^2\left\{\frac16 i\cdot f-\frac1{40}S\right\},\notag\\
\text{b)}\quad (a\A u)^3
&=-\frac3{10}\theta_\nu(\vartheta\nu x),\tag{3}\\
\text{c)}\quad (a\A u)^4
&=\frac{21}{10}\theta_\nu^2-\frac3{10}\Pi
-\frac65u_x\cdot\theta_\nu^3+\frac1{10}u_x^2\cdot\theta .\notag
\end{align}

(The same results are also reached by the double reduction of the expressions $a_\vartheta^2(b\theta u)^2$, $a_\vartheta^2(b\theta u)^3$, $a_\vartheta^2((a\theta u)^2(b\theta u)^2$ and $a_\vartheta^2((a\theta u)(b\theta u))^3$ after elimination of $a_\vartheta^2$.)

From formulas (3.) it follows that $(a\A u)^2$ is a reducer with respect to the modulus $\nu$. Hence:

1) The reduction of the system of $\A$: the form sequence $(\A\A' u)^2=a_\vartheta^2((a\A u)^2$ has the reducer $(a\A u)^2$ as a factor.

2) The reduction of the system arising by folding with the modulus $\A$:

The form sequence $(\theta\A u)^2$ has the reducer $(a\A u)^2$ as a factor.

For the forms $(\theta\A u)$ and $\A_\vartheta$ one obtains:
\[
(a\A u)^2(abu)=(\theta\A u)-u_x\cdot\A_\vartheta(\theta\A u),
\]
\[
(a\A u)(a\A b)=-\frac12\A_\vartheta+\frac12u_x\cdot\A_\vartheta^2.
\]

From the reducer $(\theta\A u)$, however, follows the reduction of the system arising by folding with the modulus $\A$.

\bigskip


The relatively complete system modulo $\nu$ therefore consists of the six forms:
\[
f,\ \theta,\ K,\ j,\ \A,\ N.
\]

As an example of the sequence expansion in \S\ 2 we give the derivation of formula (3.)a in full; in later calculations the final formula III will be set down directly. (Polars of vanishing forms have already been omitted during the calculation; the first line gives the values inserted according to Expansion II, and the second line gives the final values found recursively, beginning with the last equation of the system.) It follows, according to Expansion II:


\begin{align*}
\text{1)}\quad&m=3,\quad n=4,\quad \lambda=2,\quad \mu=0,\quad \nu=\chi=1,\\
a_\vartheta b_\vartheta(b\theta u)^2
&=[a_\vartheta]_{b^2}b+\frac67\{a_\vartheta b_\vartheta(a\theta u)(b\theta u)-u_xa_\vartheta b_\vartheta(a\theta b)(b\theta u)\}\\
&\quad-\frac17\{a_\vartheta(a\theta u)^2b_\vartheta-2u_xa_\vartheta(a\theta u)(a\theta b)b_\vartheta+u_x^2a_\vartheta(a\theta b)^2b_\vartheta\}\\
&=\frac23(a\Delta u)^2+\frac15[a_\vartheta(a\theta u)^2]b+\frac1{30}f[a_\vartheta^2(a\theta u)^2]-\frac25u_x[a_\vartheta(a\theta u)^2]b^2\\
&\quad-\frac1{15}u_x[a_\vartheta^2(a\theta u)^2]b+\frac15u_x^2[a_\vartheta(a\theta u)^2]b^3+\frac1{30}u_x^2[a_\vartheta^2(a\theta u)^2]b^2;\\[.6em]
\text{2)}\quad&m=2,\quad n=3,\quad \lambda=1,\quad \mu=1,\quad \nu=\chi=1,\\
a_\vartheta(a\theta u)b_\vartheta(b\theta u)
&=\frac25\{a_\vartheta(a\theta u)^2b_\vartheta-u_xa_\vartheta(a\theta u)(a\theta b)b_\vartheta\}+\frac12a_\vartheta^2(b\theta u)^2\\
&\quad-\frac15\{a_\vartheta^2(b\theta u)(a\theta u)-u_xa_\vartheta^2(a\theta b)(b\theta u)\}\\
&=\frac12(a\Delta u)^2+\frac25[a_\vartheta(a\theta u)^2]b+\frac1{15}[a_\vartheta^2(a\theta u)^2]f\\
&\quad-\frac25u_x[a_\vartheta(a\theta u)^2]b^2-\frac1{10}u_x[a_\vartheta^2(a\theta u)^2]b+\frac1{30}u_x^2[a_\vartheta^2(a\theta u)^2]b^2;\\[.6em]
\text{3)}\quad&m=2,\quad n=3,\quad \lambda=1,\quad \mu=1,\quad \nu=1,\quad \chi=2,\\
a_\vartheta(a\theta b)b_\vartheta(b\theta u)
&=\frac25\{a_\vartheta(a\theta b)(a\theta u)b_\vartheta-u_xa_\vartheta(a\theta b)^2b_\vartheta\}\\
&=\frac25[a_\vartheta(a\theta u)^2]b^2+\frac1{30}[a_\vartheta^2(a\theta u)^2]b-\frac25u_x[a_\vartheta(a\theta u)^2]b^3-\frac1{30}u_x[a_\vartheta^2(a\theta u)^2]b^2;\\[.6em]
\text{4)}\quad&m=1,\quad n=2,\quad \lambda=0,\quad \mu=2,\quad \nu=\chi=1,\\
a_\vartheta(a\theta u)^2b_\vartheta
&=[a_\vartheta(a\theta u)^2]b+\frac23a_\vartheta^2(a\theta u)(b\theta u)\\
&=[a_\vartheta(a\theta u)^2]b+\frac16[a_\vartheta^2(a\theta u)^2]f-\frac16u_x[a_\vartheta^2(a\theta u)^2]b;\\[.6em]
\text{5)}\quad&m=1,\quad n=2,\quad \lambda=0,\quad \mu=2,\quad \nu=1,\quad \chi=2,\\
a_\vartheta(a\theta u)(a\theta b)b_\vartheta
&=[a_\vartheta(a\theta u)^2]b^2+\frac13a_\vartheta^2(a\theta b)(b\theta u)\\
&=[a_\vartheta(a\theta u)^2]b^2+\frac1{12}[a_\vartheta^2(a\theta u)^2]b-\frac1{12}u_x[a_\vartheta^2(a\theta u)^2]b^2;\\[.6em]
\text{6)}\quad&m=1,\quad n=2,\quad \lambda=0,\quad \mu=2,\quad \nu=1,\quad \chi=3,\\
a_\vartheta(a\theta b)^2b_\vartheta&=[a_\vartheta(a\theta u)^2]b^3;\\[.6em]
\text{7)}\quad&m=2,\quad n=4,\quad \lambda=2,\quad \mu=\nu=\chi=0,\quad (\text{according to Expansion I}),\\
a_\vartheta^2(b\theta u)^2&=[a_\vartheta^2]_{ub^2}+\frac1{10}\{[a_\vartheta^2(a\theta u)^2]f-2u_x[a_\vartheta^2(a\theta u)^2]b+u_x^2[a_\vartheta^2(a\theta u)^2]b^2\};\\[.6em]
\text{8)}\quad&m=1,\quad n=3,\quad \lambda=1,\quad \mu=1,\quad \nu=\chi=0,\quad (\text{Expansion I}),\\
a_\vartheta^2(a\theta u)(b\theta u)&=\frac14[a_\vartheta^2(a\theta u)^2]f-\frac14u_x[a_\vartheta^2(a\theta u)^2]b;\\[.6em]
\text{9)}\quad&m=1,\quad n=3,\quad \lambda=1,\quad \mu=\chi=1,\quad \nu=0,\quad (\text{Expansion I}),\\
a_\vartheta^2(a\theta b)(b\theta u)&=\frac14[a_\vartheta^2(a\theta u)^2]b-\frac14u_x[a_\vartheta^2(a\theta u)^2]b^2.
\end{align*}

From 1) and 4) it follows:
\[
\begin{aligned}
(a\Delta u)^2
&=\frac65[a_\vartheta(a\theta u)^2]b+\frac15f[a_\vartheta^2(a\theta u)^2]+\frac35u_x[a_\vartheta(a\theta u)^2]b^2-\frac3{20}u_x[a_\vartheta^2(a\theta u)^2]b\\
&\quad-\frac3{10}u_x^2[a_\vartheta(a\theta u)^2]b^3-\frac1{20}u_x^2[a_\vartheta^2(a\theta u)^2]b^2,\\
(a\Delta u)^2
&=-a_\nu b_\nu+\frac35fa_\nu^2+\frac45u_xa_\nu^2b_\nu-\frac3{20}u_x^2a_\nu^2b_\nu^2-\frac1{12}u_x^2if.
\end{aligned}
\]

(For conversion into the symbols $\theta$ cf. \S\ 8 and \S\ 9.) Formula (3.)a can be computed still more briefly by direct transfer of Expansion I, introducing the auxiliary variable $w=x\widehat{u}b$ whereas already for formulas (3.)b and (3.)c the path taken here becomes the simpler one; for (3.)b one knows in advance, by \S\ 9, that all terms with factor $u_x$ are to be omitted. For the calculation of (3.)b and (3.)c one obtains:
\[
\begin{aligned}
(3.)\mathrm{b)}\quad 1)&\quad a_\vartheta(a\theta u)b_\vartheta(b\theta u)^2=\frac12(a\Delta u)^3+\frac45[a_\vartheta(a\theta u)^2]_{ub}b+\frac7{360}[a_\vartheta^2(a\theta u)^2]_{ub^2},\\
2)&\quad a_\vartheta(a\theta u)b_\vartheta(b\theta u)^2=[a_\vartheta(a\theta u)^2]_{ub}b+\frac{13}{36}[a_\vartheta^2(a\theta u)^2]_{ub^2};\\[.5em]
(3.)\mathrm{c)}\quad 1)&\quad a_\vartheta(a\theta u)b_\vartheta(b\theta u)^3=\frac12(a\Delta u)^4+\frac65[a_\vartheta(a\theta u)^2]_{ub^2}b+\frac{53}{120}[a_\vartheta^2(a\theta u)^2]_{ub^3}\\
&\qquad-\frac65u_x[a_\vartheta(a\theta u)^2]_{ub^3}b^2-\frac35u_x[a_\vartheta^2(a\theta u)^2]_{ub^2}b+\frac{19}{120}u_x^2[a_\vartheta^2(a\theta u)^2]_{ub^3}b^2,\\
2)&\quad a_\vartheta(a\theta u)b_\vartheta(b\theta u)^3=\frac34[a_\vartheta^2(a\theta u)^2]_{ub^2}.
\end{aligned}
\]

Supplementary remark: the transvectants of $f$ over $j$ that are reducible according to \S\ 5 are computed analogously. One obtains
\[
\tag{4}
\begin{aligned}
g_\nu&=-\theta_\nu+u_x\left(\frac92\theta_\nu^2-\frac12H\right)-3u_x^2\theta_\nu^3+\frac12u_x^3\theta,\\
g_\nu^2&=\frac{21}{10}\theta_\nu^2-\frac3{10}H-2u_x\theta_\nu^3+\frac12u_x^2\theta,\\
g_\nu^3&=-\frac7{10}\theta_\nu^3+\frac12u_x^2\theta,\qquad g_\nu^4=\frac35\theta.
\end{aligned}
\]




\[
\tag{4}
g_\nu^3=-\frac7{10}\theta_\nu^3+\frac12u_x^2\theta,
\qquad
g_\nu^4=\frac35\theta.
\]

From (3.) and (4.) it follows, by transfer from the binary domain, that
\[
(\theta\theta'u)^2=\frac13\,j\cdot f\pmod{\nu}.\footnotemark
\]
The remaining forms of the form sequence $(\theta\theta'u)^2$ and the form $\theta'_\nu$ are reducible to symbols $\nu$. We can therefore replace:
\[
\bmod\ \nu:\quad \text{foldings with } f\cdot j
\text{ by the corresponding foldings with }(\theta\theta'u)^2.
\]
Furthermore,
\[
(\vartheta\vartheta'x)^2=\frac43 f\cdot\A,
\qquad
\theta_{g\nu}(\vartheta\vartheta'x)=-N.
\]
\footnotetext{Gordan--Kerschensteiner, loc. cit., p. 181.}

\begin{center}
{\Large\bfseries Chapter III.\quad The relatively complete system modulo $(s,\varrho,t)$.}
\end{center}

\begin{center}
\textbf{\S\ 7.\quad Overview of the formation of the system.}
\end{center}

In order to pass from the relatively complete system modulo $\nu$ to the relatively complete system with the next higher modulus, one has, by Definition \S\ 3, to form the system of $\nu$ with respect to this higher modulus and to fold it with the forms of the relatively complete system modulo $\nu$. But $\nu=u_\nu^4$, as a contravariant of the fourth degree, is dual to the form $f$. If therefore we regard $\nu$ as the ground form, we obtain a relatively complete system of six forms modulo $\nu(\nu)=(\nu\nu_1x)^4=s_x^4$.

Thus, for the formation of the relatively complete system modulo $s$ (System III), we have to transvect
\[
\begin{array}{ll}
\text{System I:} & f,\ \theta,\ K,\ j,\ \A,\ N\quad \text{over}\vphantom{\dfrac11}\\[0.3em]
\text{System II:} & \nu,\ H=(\nu\nu_1x)^2u_\nu^2u_{\nu_1}^2,\ L=(\nu\eta x)H_x^2u_\nu^3u_\eta^3,\\
& g=(\nu\eta x)^4H_x^2,\ \sigma=H_\eta^2u_\nu^2u_\eta^4,\ (\nu\sigma x).
\end{array}
\]
It will be shown (formula (13.)) that the form $g$ is reducible, and hence that System II consists of only five forms.

In the course of the calculation the quadratic forms
\[
\varrho=u_\varrho^2=\theta_\nu^4u_\vartheta^2,
\qquad
 t=t_x^2=a_\eta^4H_x^2
\]
are adjoined as moduli, so that one obtains a relatively complete system modulo $(s,\varrho,t)$; and the modulus $(\varrho,t)$ is to be regarded as a higher modulus than the modulus $s$.

The arrangement of the individual forms is to follow the order in the coefficients. We normalize the folding
\[
\theta_\nu(K_\nu,\theta_\nu,\ldots)
\]
as higher than the folding
\[
H_y(H_y,L_y,\ldots).
\]
Consequently, by \S\ 1 and \S\ 3b, the order of the individual forms of the same order is uniquely determined.

The formation of the forms of the same order is carried out in tabular form:

I. Indication of which forms of Systems I and II are to be folded with each other in order to obtain the prescribed order in the coefficients.

II. Listing of the irreducible forms according to the scheme of the form sequence.

III. Reduction of the reducible or decomposable form sequences or forms, that is, of those places where the total form sequence breaks off, thereby proving that all irreducible constructions have been exhausted by those indicated.

IV. Consequences: 1) indication of newly arising reducers; 2) indication of those forms which do not enter, in products with forms of the same system, into folding with forms of the other system.

It will be shown that only the following occur:

1) foldings of one form of System I with one form of System II;

2) foldings of powers of $\theta$, respectively $H$, or of two forms of one system with one form of the second system.

We shall obtain the following scheme for folding:
\[
\begin{array}{r@{\ }c@{\ }l@{\qquad}c@{\qquad}r@{\ }c@{\ }l}
\multicolumn{3}{c}{\text{System I}} & \text{folded with} & \multicolumn{3}{c}{\text{System II}}\\[0.25em]
1. & \text{order} & f & & 2. & \text{order} & \nu\\
2. & ,, & \theta & & 4. & ,, & H\\
3. & ,, & K,j,\A & & 6. & ,, & L,\sigma\\
4. & ,, & N,\theta^2,f\cdot j & & 8. & ,, & (\nu\sigma x),H^2\\
5. & ,, & \theta\cdot K & & 10. & ,, & H\cdot L\\
6. & ,, & \theta^3,\A\cdot j & & 12. & ,, & H^3.
\end{array}
\]

For checking the calculation, besides the ``double reduction'', the two special forms serve:


\[
\begin{array}{ll}
\text{I.}&\begin{array}{rlrlrl}
f&=\displaystyle\sum_3x_1^3x_2,& u_\nu^2&=\frac{i}{6}u_x^2,& s&=\frac{i}{3}f,\\
a_\eta^2&=-\frac19 i\theta,& H_\nu^2&=\frac{i}{6}\theta,& \sigma&=-\frac{i}{6}j,\\
g&=-\frac23i\Delta,& \varrho&=0,& t&=0,
\end{array}\\[1.2em]
\text{II.}&\begin{array}{rlrlrl}
f&=\displaystyle\sum_3x_1^4,& s&=\frac43if,& a_\eta^2&=\frac29i\theta,\\
\Delta_\nu^2&=\frac1{15}i\theta,& \sigma&=\frac43ij,& g&=\frac43i\Delta,\\
\varrho&=0,& t&=0.
\end{array}
\end{array}
\]

Supplementary remark: for the ground form $\nu$, formulas (1.), (2.), (3.), (4.) correspond to the following:
\[
\tag{5}
\begin{array}{c|c|c}
(\nu\eta x)^2=A & H_\nu(\nu\eta x)=0 & H_\nu^2=\sigma\\[0.4em]
(\nu\eta x)^3=0 & H_\nu(\nu\eta x)^2=B & H_\nu^2(\nu\eta x)=0\\[0.4em]
(\nu\eta x)^4=g & H_\nu(\nu\eta x)^3=0 & H_\nu^2(\nu\eta x)^2=C,
\end{array}
\]
\[
A=\frac16s\nu-\frac23u_xs_\nu+\frac23u_x^2s_\nu^2-\frac{J}{18}u_x^4,
\quad B=-\frac56s_\nu+\frac76u_xs_\nu^2-\frac{J}{9}u_x^3,
\quad C=\frac23s_\nu^2-\frac{J}{9}u_x^2.
\]
\[
\tag{6}
s_\nu^3u_xs_\nu=\frac13u_xs_\nu^4=\frac13Ju_x.
\]


\[
\tag{7}
\begin{aligned}
\text{a)}\quad (\nu\sigma x)^2
&=\frac12(Hsu)^2-\frac25\nu\cdot s_\nu^2-
\frac45u_x\,s_\eta(Hsu)^2
+u_x^2\left\{\frac16J\cdot\nu-\frac1{40}(ss'u)^4\right\},\\
\text{b)}\quad (\nu\sigma x)^3&=-\frac3{10}s_\eta(Hsu),\\
\text{c)}\quad (\nu\sigma x)^4&=\frac{21}{10}s_\eta^2-\frac3{10}Z-
\frac65u_xs_\eta^3+\frac1{10}u_x^2s_\eta^4.
\end{aligned}
\]
\[
\tag{8}
\begin{aligned}
\text{a)}\quad g_\nu&=-s_\eta+u_x\left(\frac92s_\eta^2-\frac12Z\right)-3u_x^2s_\eta^3+\frac12u_x^3s_\eta^4,\\
\text{b)}\quad g_\nu^2&=\frac{21}{10}s_\eta^2-\frac3{10}Z-2u_xs_\eta^3+\frac12u_x^2s_\eta^4,\\
\text{c)}\quad g_\nu^3&=-\frac7{10}s_\eta^3+\frac12u_xs_\eta^4,\\
\text{d)}\quad g_\nu^4&=\frac35s_\eta^4.
\end{aligned}
\]

\begin{center}
\textbf{\S\ 8.\quad Forms of the third order (System III).}
\end{center}

\noindent\textbf{Folding of $f$ with $\nu$.}

\noindent\emph{Irreducible forms:}
\[
\begin{array}{cccc}
a_\nu&\cdot&\cdot&\cdot\\
\cdot&a_\nu^2&\cdot&\cdot\\
\cdot&\cdot&\cdot&\cdot\\
\cdot&\cdot&\cdot&a_\nu^4=i.
\end{array}
\]

\noindent\emph{Reduction:}
\[
a_\nu^3=\frac13 i u_x\qquad\text{(formula (2.)).}
\]
\noindent\emph{Consequences:} 1) $a_\nu^3$ is a reducer. 2) $f$ enters into folding with $\nu$ only in products with contragredient forms $(j)$. The same holds for $\nu$ with respect to $f$.

\begin{center}
\textbf{\S\ 9.\quad Forms of the fourth order (System III).}
\end{center}

\noindent\textbf{Folding of $\theta$ with $\nu$.}

\noindent\emph{Irreducible forms:}
\[
\begin{array}{cccc}
(\theta\nu x)&\theta_\nu&\cdot&\cdot\\
(\theta\nu x)^2&\theta_\nu(\theta\nu x)&\theta_\nu^2&\cdot\\
\cdot&\theta_\nu(\theta\nu x)^2&\cdot&\theta_\nu^3.
\end{array}
\]

\noindent\emph{reductions:} From the reducer $a_\nu^3$, by double reduction of the expressions $a_\nu^3(abu)$, $a_\nu^3b_\nu$, $a_\nu^3b_\nu(abu)$ according to the ansatz:


\[
\begin{aligned}
\theta_\nu^2(\vartheta\nu x)
&=a_\nu^2(\widehat{a}b\nu x)(abu)-\frac13(\nu\nu_1x)^3=a_\nu b_\nu(\widehat{a}b\nu x)(abu)+\frac16(\nu\nu_1x)^3\\
&=a_\nu^3(abu)-a_\nu^2b_\nu(abu)=2a_\nu^2b_\nu(abu)=\frac23a_\nu^3(abu),\\[.4em]
\theta_\nu^2(\vartheta\nu x)^2
&=a_\nu^2(\widehat{a}b\nu x)^2-\frac13(\nu\nu_1x)^4=a_\nu b_\nu(\widehat{a}b\nu x)^2+\frac16(\nu\nu_1x)^4\\
&=if-2a_\nu^3b_\nu+a_\nu^2b_\nu^2-\frac13s=2a_\nu^3b_\nu-2a_\nu^2b_\nu^2+\frac16s=\frac23if-\frac23a_\nu^3b_\nu-\frac16s,\\[.4em]
\theta_\nu^3(\vartheta\nu x)&=a_\nu^2b_\nu(\widehat{a}b\nu x)(abu)=a_\nu^3b_\nu(abu).
\end{aligned}
\]


The formulas obtained are:
\[
\tag{9}
\begin{array}{c|c|c}
\theta_\nu^2(\theta\nu x)=0
& \theta_\nu^3(\theta\nu x)=0
& \theta_\nu^4u_x^2=u_\varrho^2=\varrho\quad(\text{adjoined modulus, \S\ 7}),\\[0.6em]
\theta_\nu^2(\theta\nu x)^2=\dfrac49 i f-\dfrac16s
& &
\end{array}
\]

\noindent\emph{Consequences:} 1) $\theta_\nu^2(\theta\nu x)^2$ is a reducer. 2) $\nu$ enters into folding with $\theta$ only in products with contragredient forms $(g)$.



\begin{center}
\textbf{\S\ 10.\quad Forms of the fifth order (System III).}
\end{center}

\[
\text{Folding of }\left\{\begin{array}{l}
K,j,\A \text{ with }\nu,\\
f\text{ with }H.
\end{array}\right.
\]

\noindent\emph{Irreducible forms:}
\[
\begin{array}{c@{\qquad}c@{\qquad}c@{\qquad}c}
\text{A)} & \begin{array}{cccc}
\cdot&K_\nu&\cdot&\cdot\\
(k\nu x)^2&\cdot&K_\nu^2&\cdot\\
(k\nu x)^3&K_\nu(k\nu x)^2&\cdot&\cdot\\
\cdot&K_\nu(k\nu x)^3&\cdot&\cdot
\end{array}
&
\text{B)}\ \begin{array}{c}(j\nu x)\\(j\nu x)^2\\(j\nu x)^3\\ \cdot\end{array}
&
\text{C)}\ \begin{array}{c}\A_\nu\\\cdot\\\cdot\\\cdot\end{array}
\end{array}
\]
\[
\text{D)}\qquad
\begin{array}{cccc}
(aHu)&(aHu)^2&\cdot&\cdot\\
a_\eta&a_\eta(aHu)&a_\eta(aHu)^2&\cdot\\
\cdot&a_\eta^2&\cdot&\cdot
\end{array}
\]

\noindent\emph{reductions:}

\noindent A) form sequence $K_\nu^3$: reducer $a_\nu^3$.

form $K_\nu^2(k\nu x)^2$: reducer $\theta_\nu^2(\vartheta\nu x)^2$.
\[
(k\nu x),\quad K_\nu(k\nu x),\quad K_\nu^2(k\nu x)
\]
are decomposable forms according to \S\ 3.

\noindent B) form
\[
\begin{aligned}
(j\nu x)^4&=(\hata\theta\nu x)^4
=a_\nu^4\theta+\theta_\nu^4 f-4a_\nu^3\theta_\nu
-4a_\nu\{a_\nu\theta_x-(\hata\theta\nu x)\}^3\\
&\quad+6a_\nu^2\{a_\nu\theta_x-(\hata\theta\nu x)\}^2
=\R f+\frac53 i\theta-6a_\nu^2(\hata\theta\nu x)^2
+4a_\nu(\hata\theta\nu x)^3,
\end{aligned}
\]
and therefore
\[
(j\nu x)^4=\R f+\frac53 i\theta\modu{(\A,\nu)},\qquad (\text{cf. \S\ 5 end}).
\]

\noindent C) form $\A_\nu^4$: reducer $\theta_\nu^4$.

forms $\A_\nu^2$ and $\A_\nu^3$: reducer $a_\nu^3$ or $\theta_\nu^4$ by replacing the form $\A$ by lower forms of the form sequence (\S\ 5 end), i.e. by double reduction of the expressions
\[
a_\vartheta a_\nu^3,
\qquad a_\vartheta a_\nu^3\theta_\nu,
\qquad a_\vartheta\theta_\nu^4.
\]
The double calculation of $\A_\nu^3$ gives rise to the reduction of the higher form $a_\eta^3$.

\noindent\emph{Reduction formulas:}
\[
\tag{10}
\begin{aligned}
\text{a)}\quad \A_\nu^2
&=\frac35a_\eta^2-\frac3{10}(asu)^2+\frac13 i\theta
+\frac15u_x^2(2a_\R^2-t),\\
\text{b)}\quad \A_\nu^3
&=-\frac3{10}a_\R+\frac35u_x\left(\frac{13}{10}a_\R^2-\frac25t\right),\\
\text{c)}\quad \A_\nu^4&=a_\R^2-\frac25t.
\end{aligned}
\]

\noindent\emph{Derivation of the formulas:}
\[
\begin{aligned}
\text{a)}\quad
 a_\vartheta a_\nu^3
&=\frac13 i\theta
=[a_\vartheta]_{\nu^3}+\frac67[a_\vartheta^2]_{\nu^2}
+\frac65[a_\vartheta(a\theta u)^2]_{\nu^2}\,
u x^2
+\frac{11}{30}[a_\vartheta^2(a\theta u)^2]_\nu\,\nu x^2\\
&\quad-\frac67u_x[a_\vartheta^2]_{\nu^3}
-\frac16u_x[a_\vartheta^2(a\theta u)^2]_{\nu^2}\,\nu x^2,\\
\frac13 i\theta
&=\A_\nu^2-\frac23u_x\A_\nu^3-a_\nu a_{\nu_1}(\nu\nu_1x)^2
+\frac25a_\nu^2(\nu\nu_1x)^2
+\frac15u_x^2a_\nu^2a_{\nu_1}(\nu\nu_1x)^2;
\end{aligned}
\]
\[
\begin{aligned}
\text{b)}\quad
 a_\vartheta a_\nu^3\theta_\nu
&=0=[a_\vartheta]_{\nu^4}+\frac9{14}[a_\vartheta^2]_{\nu^3}
+\frac35[a_\vartheta(a\theta u)^2]_{\nu^2}\,\nu x^2
+\frac{17}{120}[a_\vartheta^2(a\theta u)^2]_{\nu}\,\nu x^2\\
&\quad-\frac9{14}u_x[a_\vartheta^2]_{\nu^4}
-\frac1{24}u_x[a_\vartheta^2(a\theta u)^2]_{\nu^2}\,\nu x^2,\\
0&=\frac56\A_\nu^3-\frac12u_x\A_\nu^4-\frac14a_\eta^3+\frac1{20}u_xa_\eta^4;
\end{aligned}
\]
\[
\begin{aligned}
\text{b')}\quad
 a_\vartheta\theta_\nu^4
&=a_\R=a_\vartheta\theta_\nu\{a_\nu-(a\theta\nu x)\}^3\\
&=-\frac32[a_\vartheta]_{\nu^3}+\frac35[a_\vartheta(a\theta u)^2]_{\nu^2}\,\nu x^2
-\frac{37}{120}[a_\vartheta^2(a\theta u)^2]_{\nu}\,\nu x^2\\
&\quad+\frac32u_x[a_\vartheta^2]_{\nu^4}
+\frac{49}{120}u_x[a_\vartheta^2(a\theta u)^2]_{\nu^2}\,\nu x^2,\\
a_\R&=-\frac32\A_\nu^3+\frac32u_x\A_\nu^4-\frac{11}{20}a_\eta^3+\frac7{20}u_xa_\eta^4;
\end{aligned}
\]
\[
\begin{aligned}
\text{c)}\quad
 a_\vartheta^2\theta_\nu^4
&=a_\R^2=[a_\vartheta^2]_{\nu^4}+\frac35[a_\vartheta^2(a\theta u)^2]_{\nu^2}\,\nu x^2,\\
a_\R^2&=\A_\nu^4+\frac25a_\eta^4.
\end{aligned}
\]

\noindent D) reduction of $a_\eta^3$ by double reduction of $\A_\nu^3$ (C).

The remaining reductions are obtained by a computation dualistically opposite to that of \S\ 9.

\noindent\emph{Reduction formulas:}
\[
\tag{11}
\begin{array}{rl|rl}
a_\eta^2(aHu)&=-\dfrac16(asu)^3,&
 a_\eta^2(aHu)^2&=\dfrac49 i\nu-\dfrac16(asu)^4,\\[0.7em]
a_\eta^3&=-a_\R+\dfrac35u_xa_\R^2+\dfrac15u_xt,
& a_\eta^3(aHu)&=0,\\[0.7em]
&&a_\eta^4&=t\quad(\text{adjoined as a modulus}).
\end{array}
\]

\noindent\emph{Consequences:}

1) $(j\nu x)^4$, $\A_\nu^2$, $a_\eta^2(aHu)$ are reducers.

2) $\nu$ enters only in products with contragredient forms in foldings with $K,j,\A$; the same holds for $f$ in relation to $H$, and for $j$ in relation to $\nu$, since $\theta_\nu(\theta\cdot j,\nu,x^3)$ becomes reducible according to \S\ 11 B.

\begin{center}
\textbf{\S\ 11.\quad Forms of the sixth order (System III).}
\end{center}
\[
\text{Folding of }\left\{\begin{array}{l}
\theta^2,
 f\cdot j,
 N \text{ with }\nu,\\
\theta\text{ with }H.
\end{array}\right.
\]

\noindent\emph{Irreducible forms:}
\[
\begin{array}{c@{\quad}c@{\quad}c}
\text{A)}\ (\vartheta^2\nu x)^3,\quad \theta_\nu(\vartheta^2\nu x)^3
&
\text{B)}\ a_\nu(j\nu x),\quad a_\nu^2(j\nu x)
&
\text{C)}\ N_\nu
\end{array}
\]
\[
\text{D)}\quad
\begin{array}{c|c|c|c}
 & (\theta Hu)&(\theta Hu)^2&\cdot\\
(\vartheta\eta x)&\theta_\eta&H_\vartheta(\theta Hu),\ \theta_\eta(\theta Hu)&\theta_\eta(\theta Hu)^2\\
(\vartheta\eta x)^2&H_\vartheta(\vartheta\eta x),\ \theta_\eta(\vartheta\eta x)&\theta_\eta^2&\theta_\eta^2(\theta Hu)\\
\cdot&\theta_\eta(\vartheta\eta x)^2&\cdot&\cdot
\end{array}
\]

\noindent\emph{reductions:}

\noindent A) $(\vartheta^2\nu x)^4$ decomposes according to formula (3.)a:
\[
(a\A\vartheta x)^4=\frac12(\vartheta^2\nu x)^4-\frac35 f\cdot a_\nu^2(\nu\vartheta x)^2
=\frac13\A^2+f\A_\vartheta^2.
\]

\noindent B) $a_\nu(j\nu x)^2$ decomposes according to formula (9.)a, indem wir according to \S\ 6 end $f\cdot j$ by $(\theta\theta'u)^2$ ersetzen, and the reducibility of $\theta_\vartheta'$ take account of:
\[
\frac13a_\nu(j\nu x)^2=(\theta\theta'\nu x)^2\theta_\nu
=\theta_\nu^3\theta-\theta_\nu^2\theta_\nu'
=\theta_\nu^3\theta+\theta_\nu^2(\widehat{j}\nu\theta'u)
=\theta_\nu^3\theta-\frac23\theta_\nu^3\theta.
\]
forms $a_\nu(j\nu x)^3$ and $a_\nu^2(j\nu x)^2$: reducer $a_j$ and $(j\nu x)^4$:
\[
a_\nu(j\nu x)^3=a_j(j\nu x)^3-(j\nu x)^3(j\nu\widehat{a}u),
\]
\[
a_\nu^2(j\nu x)^2=a_j^2(j\nu x)^2-2a_j(j\nu x)^2(j\nu\widehat{a}u)+(j\nu x)^2(j\nu\widehat{a}u)^2.
\]

\noindent C) form sequence $N_\nu^2$; reducer $\A_\nu^2$. $(n\nu x)$ and $N_\nu(n\nu x)$ decompose as a transvection over functional determinants according to \S\ 3.

\noindent D) \emph{Overview of the reductions:}

a) form sequence $\theta_\eta^2H_\vartheta$: reducer $a_\eta^2(aHu)$. (Leads to forms with higher symbols.)

b) form sequence $\theta_\eta H_\vartheta$: reducer $a_\eta^2(aHu)$ by double reduction. (Leads to forms with higher symbols and to higher forms of the total form sequence, i.e. to the form sequence $\theta_\eta^2$.) Instead of $\theta_\eta H_\vartheta$ we consider the form sequence
\[
\theta_\eta(\vartheta\eta x)(\theta Hu)=\theta_\eta H_\vartheta+\theta_\eta^2,
\]
ersetzen in it the symbols $\theta$ by the lower form $(abu)$, gelangen so to the double reduction of the form sequence
\[
a_\eta^2(aHu)(abu)=\theta_\eta H_\vartheta+\frac32\theta_\eta^2\modu{s}
\]
up to the terminal form
\[
a_\eta^3b_\eta(bHu)(abu)=\theta_\eta^3H_\vartheta+\frac12\theta_\eta^4.
\]

c) form sequence $\theta_\eta^3$: By double reduction of those forms which simultaneously belong to the form sequences $\theta_\eta H_\vartheta$ and $\theta_\eta^2H_\vartheta$ - i.e. the forms of the form sequence $\theta_\eta^2H_\vartheta$ - to forms with higher symbols on the one hand, and to forms with higher symbols and to the higher form sequence $\theta_\eta^3$ on the other hand.

d) forms $\theta_\eta^2(\vartheta\eta x)$, $\theta_\eta^2(\vartheta\eta x)^2$, $\theta_\eta^3(\vartheta\eta x)$: By double reduction by means of the reducers $a$ analogous to the computation of \S\ 9.

e) forms $\theta_\eta^2(\theta Hu)^2$ and $\theta_\eta^2(\vartheta\eta x)^2$: By double reduction of the expressions $\A_\nu^2(a\A u)^4$ and $(a\A\nu x)^4$ by means of the reducers $\A_\nu^2$ and $(a\A u)^4$.

f) forms $H_\vartheta$ and $H_\vartheta^2$: Decomposition of forms. Is obtained for $H_\vartheta^2$ by double reduction of the expression $\A_\nu^2(a\A u)^2$ or also by direct calculation of the products $(a_\nu^2)^2$, $a_\nu^2b_{\nu_1}$ by forms of the system. (The analogous calculation of $(a_\nu)^2$ gives a relation between decomposable forms.)

g) Reduction of higher forms by the fact that forms of the form sequence $\theta\cdot H$ admit two reduction possibilities (belong to two reducible form sequences), namely:
\begin{quote}\small
$g$: by double reduction of $\theta_\eta^2(\vartheta\eta x)^2$; allows reduction d) and e) zu.\par
$s_\vartheta^2(\theta su)$: by double reduction of $\theta_\eta^3(\vartheta\eta x)$; allows reduction c) and d) zu.\par
$s_\vartheta^2(\theta su)^2$: by double reduction of $\theta_\eta^2H_\vartheta^2$; allows reduction a) and has the reducers $\theta_\nu^2(\vartheta\nu x)^2$ as a factor.\par
$s_\nu^2$: by double reduction of $\theta_\eta^3H_\vartheta$; allows reduction a) and c) zu.
\end{quote}
The proof of the independence of the remaining forms is obtained by double reduction of the form sequence $\A_\nu^2(a\A u)^2=\theta_\eta^2$.

\noindent\emph{Execution of the reductions:}

The existence of reductions a), b), c), d) is clear a priori, since according to the stated law of formation the relations arising from double reduction are independent of one another. For reductions e), f), g), it must be shown by computation that no identities arise.

Proceeding symmetrically with the diagonal term in the form sequence $\theta\cdot H$ up to the form $\theta_\eta$, and partly indicating the calculation, we also give the formulas obtained by reductions a), b), c), d), since they are used partly in reductions e), f), g), and partly later.

\noindent 1) $H_\vartheta$ and $H_\vartheta^2$ according to f). According to the ansatz\footnote{The sign $\sim$ in place of the equality sign means that products arising from $u_x$ with higher forms, other than those arising by foldings III and IV, are omitted.}
\[
\begin{aligned}
a_\nu b_{\nu_1}^2
&=\nu a_\nu b_\nu^2-(aHu)(bHu)b_\eta
\sim \nu a_\nu b_\nu^2-f a_\eta(aHu)^2-(abu)(aHu)b_\eta+(abu)^2b_\eta,\\
a_\nu^2b_{\nu_1}^2
&=b_\nu^2\{a_\nu u_\nu-(a\nu\nu_1u)\}^2
\sim \nu\{a_\nu^2b_\nu^2-2a_\nu b_\nu(aHu)(bHu)
-\frac16(asu)^2(bsu)^2\}\\
&\sim \nu a_\nu^2b_\nu^2-2a^2(aHu)(abu)-\frac16(asu)^2(bsu)^2+(\vartheta\eta x)^2(\theta Hu)^2
\end{aligned}
\]
after substitution of the value of $\theta_\eta H_\vartheta$, the formulas arise:
\[
\tag{12}
\begin{aligned}
\text{a)}\quad H_\vartheta&\sim 2a_\nu a_\R^2+2\nu b_\nu(\vartheta\eta x)^2+2f a_\eta(aHu)^2-2\theta_\eta,\\
\text{b)}\quad H_\vartheta^2&\sim (a^2)^2+2\theta_\eta^2+\frac23(\theta su)^2
+\frac1{12}s\nu-\frac19if\nu-\frac16f(asu)^4.
\end{aligned}
\]

\noindent 2) $\theta_\eta H_\vartheta$ according to b):
\[
\begin{aligned}
a_\eta^2(aHu)(abu)&\sim [a_\eta^2(aHu)]_{bu}+\frac12 f[a_\eta^2(aHu)^2]
=a_\eta b_\eta(aHu)(abu)\\
&\quad+a_\eta(a\widehat{b}\eta x)(abu)(aHu)
\sim \theta_\eta(\vartheta\eta x)(\theta Hu)+\frac12\theta_\eta^2+\frac14(\nu\eta x)^2.
\end{aligned}
\]
According to the formulas (5.) and (11.) follows:
\[
\theta_\eta H_\vartheta\sim -\frac32\theta_\eta^2-\frac14(\theta su)^2-\frac1{12}s\nu+\frac29if\nu.
\]

\noindent 3) $\theta_\eta H_\vartheta(\theta Hu)$ is computed according to b), analogously to $\theta_\eta H_\vartheta$:
\[
\theta_\eta H_\vartheta(\theta Hu)\sim-\theta_\eta^2(\theta Hu)-\frac16(\theta su)^3.
\]

\noindent 4) $\theta_\eta H_\vartheta(\vartheta\eta x)$ according to b); $\theta_\eta^2(\vartheta\eta x)$ according to d) (cf. \S\ 9):
\[
\theta_\eta H_\vartheta(\vartheta\eta x)\sim-\theta_\eta^2(\vartheta\eta x)-\frac16s_\vartheta(\theta su)
\sim-\frac13(\vartheta\R x)-\frac16s_\vartheta(\theta su),
\]
\[
\theta_\eta^2(\vartheta\eta x)\sim\frac23a_\eta^3(abu)\sim-\frac13(\vartheta\R x).
\]

\noindent 5)
\[
\begin{array}{@{}c@{\quad}c@{\quad}c@{}}
\theta_\eta H_\vartheta^2\text{ according to b)}&
\theta_\eta^2H_\vartheta\text{ according to a)}&
\begin{array}{c}\theta_\eta^2H_\vartheta\text{ according to b)}\\[-0.2em]\text{and thereby }\theta_\eta^3\text{ according to c)}\end{array}\\
\text{from }a_\eta^2(aHb)(bHu)&
\text{from }a_\eta^3(Hab)(abu)&
\text{from }a_\eta^3(bHu)(abu)
\end{array}
\]
\[
\theta_\eta H_\vartheta^2\sim-\frac32\theta_\eta^2H_\vartheta-\frac14s_\vartheta(\theta su)^2+\frac16s_\nu-\frac49ia_\nu
\sim-\theta_\R-\frac16s_\nu-\frac19ia_\nu,
\]
\[
\theta_\eta^2H_\vartheta\sim\frac23\theta_\R-\frac16s_\vartheta(\theta su)^2+\frac29s_\nu-\frac29ia_\nu,
\]
\[
\theta_\eta^3H_\vartheta\sim\frac23\theta_\R-\frac23\theta_\eta^3+\frac5{36}s_\nu,
\qquad
\theta_\eta^3\sim\frac14s_\vartheta(\theta su)^2-\frac18s_\nu+\frac13ia_\nu.
\]

\noindent 6) $\theta_\eta^2(\theta Hu)^2$ according to e):
\[
\A_\nu^2(a\A u)^2=[(a\A u)^4]_{\nu^2}
=-\frac{12}{5}\theta_\eta^2(\theta Hu)^2-\frac3{10}\sigma-\frac1{20}(\theta su)^4+\nu\R
\quad(\text{according to formula (3.)c}),
\]
\[
\A_\nu^2(a\A u)^4=[\A_\nu^2]_{aa}^4=-\frac35\theta_\eta^2(\theta Hu)^2+\frac15\sigma-\frac3{10}(\theta su)^4+\frac13ij
\quad(\text{according to formula (10.)a}),
\]
\[
\theta_\eta^2(\theta Hu)^2=-\frac16\sigma+\frac1{12}(\theta su)^4-\frac19ij+\frac13\nu\R.
\]

\noindent 7) $\theta_\eta^2(\vartheta\eta x)^2$ according to d) and e). From this reduction of the higher form $g$.

By a computation analogous to that in \S\ 9, it follows:
\[
\begin{aligned}
\theta_\eta^2(\vartheta\eta x)^2
&=\frac12 i f-\frac12a_\eta^2b_\eta-\frac1{12}g
=\frac23tf-\frac23a_\eta^3b_\eta-\frac16g\\
&=-\frac16g-\frac13(\vartheta\R x)^2+\frac4{15}fa_\R^2+\frac8{15}ft
\quad(\text{according to formula (11.)}).
\end{aligned}
\]
By the corresponding computation as above follows:
\[
((a\A\nu x)^4)=[(a\A u)^4]_{\nu x^4}
=\frac{21}{10}\theta_\eta^2(\vartheta\eta x)^2+\frac7{10}s_\vartheta^2-\frac3{10}g
\quad(\text{according to formula (3.)c}),
\]
\[
\begin{aligned}
(a\A\nu x)^4&=-\frac13i\A+6[\A_\nu^2]a^2-4[\A_\nu^3]a+\A_\nu^4f\\
&= -\frac{36}{5}\theta_\eta^2(\vartheta\eta x)^2-\frac95s_\vartheta^2-\frac35g-\frac35(\vartheta\R x)^2\\
&\quad+\frac53i\A+\frac{37}{25}fa_\R^2+\frac{74}{25}ft\quad(\text{according to formula (10.)}).
\end{aligned}
\]
Thus
\[
\begin{aligned}
\frac{93}{10}\theta_\eta^2(\vartheta\eta x)^2
&=-\frac3{10}g-\frac52s_\vartheta^2-\frac35(\vartheta\R x)^2\\
&\quad+\frac53i\A+\frac{37}{25}fa_\R^2+\frac{74}{25}ft,
\end{aligned}
\]
and by comparing the two values of $\theta_\eta^2(\vartheta\eta x)^2$ according to g):
\[
\tag{13}
0=g-2s_\vartheta^2+2(\vartheta\R x)^2+\frac43i\A-\frac45fa_\R^2-\frac85ft.
\]

\noindent 8) $\theta_\eta^2H_\vartheta(\theta Hu)$ according to a) and b), from this $\theta_\eta^3(\theta Hu)$ according to c) (forms with the Faktor $u_x$ vanish):
\[
\theta_\eta^2H_\vartheta(\theta Hu)=-\frac16s_\vartheta(\theta su)^3-\frac13(\nu\R x),
\]
\[
\theta_\eta^2H_\vartheta(\theta Hu)=-\frac13\theta_\eta^3(\theta Hu)-\frac13(\nu\R x),
\]
\[
\theta_\eta^3(\theta Hu)=\frac12s_\vartheta(\theta su)^3.
\]

\noindent 9) $\theta_\eta^2H_\vartheta(\vartheta\eta x)$ according to a) and b), from this $\theta_\eta^3(\vartheta\eta x)$ according to c). $\theta_\eta^3(\vartheta\eta x)$ according to d), from this reduction of the higher form $s_\vartheta^2(\theta su)$ according to g) (forms with the Faktor $u_x$ vanish):
\[
\theta_\eta^2H_\vartheta(\vartheta\eta x)=\frac13(atu),
\]
\[
\theta_\eta^2H_\vartheta(\vartheta\eta x)=\frac13(atu)+\frac13\theta_\eta^3(\vartheta\eta x)-\frac16s_\vartheta^2(\theta su),
\]
\[
\theta_\eta^3(\vartheta\eta x)=\frac12s_\vartheta^2(\theta su),
\]
\[
\theta_\eta^3(\vartheta\eta x)=\frac25\theta_\R(\vartheta\R x)+\frac15(atu),
\]
\[
\tag{14}
s_\vartheta^2(\theta su)=\frac45\theta_\R(\vartheta\R x)+\frac25(atu).
\]

\noindent 10) $\theta_\eta^2H_\vartheta^2$ according to a) and by reducer $\theta_\nu^2(\vartheta\nu x)^2$, $\theta_\eta^3H_\vartheta$ according to a) and c), $\theta_\eta^4$ according to c), from this reduction of the higher forms $s_\vartheta^2(\theta su)^2$ and $s_\nu^2$ according to g):
\[
\begin{aligned}
\text{According to a)}\quad
\theta_\eta^2H_\vartheta^2&=[a_\eta^2(aHu)^2]b^2+\frac13[a_\eta^4]_{bu^2}-\frac13H_\nu^2(\nu\eta x)^2\\
&=\frac49ia_\nu^2-\frac16s_\vartheta^2(\theta su)^2-\frac5{18}s_\nu^2+\frac13(atu)^2+\frac1{54}Ju_x^2.
\end{aligned}
\]
\[
\begin{aligned}
\text{By }\theta_\nu^2(\vartheta\nu x)^2:
\quad\theta_\eta^2H_\vartheta^2
&=[\theta_\nu^2(\vartheta\nu x)^2]_{\nu_1}+\frac13[\theta_\nu^4]_{\nu_1}\hat\nu_1x^2-\frac13s_\vartheta^2(\theta su)^2\\
&=\frac49ia_\nu^2-\frac16s_\nu^2-\frac13s_\vartheta^2(\theta su)^2+\frac13(\nu\R x)^2.
\end{aligned}
\]
\[
\begin{aligned}
\text{According to a)}\quad
\theta_\eta^3H_\vartheta&=-[a_\eta^2(aHu)]_{bu}b^2-\frac12[a_\eta^2(aHu)^2]b^2-\frac13[a_\eta^3]b+\frac1{12}[a_\eta^4]_{bu^2}\\
&\quad+\frac12u_x[a_\eta^2(aHu)^2]b^3\\
&=-\frac29ia_\nu^2+\frac1{12}s_\nu^2+\frac4{15}\theta_\R^2-\frac7{90}(\nu\R x)^2+\frac1{30}(atu)^2+\frac2{27}i^2u_x^2-\frac1{36}Ju_x^2.
\end{aligned}
\]
\[
\begin{aligned}
\text{According to c)}\quad
\theta_\eta^3H_\vartheta:
 a_\eta^3(Hab)(bHu)&=\frac12(atu)^2
=\theta_\eta^2H_\vartheta(\theta Hu)(\vartheta\eta x)+\frac14H_\nu^2(\nu\eta x)^2\\
&\quad+\frac12\theta_\eta^2H_\vartheta^2
=\frac32\theta_\eta^2H_\vartheta^2+\theta_\eta^3H_\vartheta-u_x[a_\eta^2(aHu)^2]b^3\\
&\quad+\frac14H_\nu^2(\nu\eta x)^2,
\end{aligned}
\]
and therefore
\[
\begin{aligned}
\theta_\eta^3H_\vartheta&=-\frac23ia_\nu^2+\frac5{12}s_\nu^2+\frac12(\nu\R x)^2-\frac12(atu)^2\\
&\quad+\frac4{27}i^2u_x^2-\frac1{12}Ju_x^2.
\end{aligned}
\]
According to c)
\[
\begin{aligned}
\theta_\eta^4&=\frac49ia_\nu^2-\frac16s_\nu^2+\frac{16}{15}\theta_\R^2\\
&\quad-\frac{14}{45}(\nu\R x)^2+\frac2{15}(atu)^2.
\end{aligned}
\]
According to g), from the two values for $\theta_\eta^2H_\vartheta^2$:
\[
\tag{15}
\begin{aligned}
s_\vartheta^2(\theta su)^2
&=\frac23s_\nu^2-2(atu)^2+2(\nu\R x)^2-\frac19Ju_x^2\\
&=\frac89ia_\nu^2+\frac8{15}\theta_\R^2-\frac{14}{15}(atu)^2+\frac{38}{45}(\nu\R x)^2-\frac4{27}i^2u_x^2.
\end{aligned}
\]
According to g), from the two values for $\theta_\eta^3H_\vartheta$:
\[
\tag{16}
s_\nu^2=\frac43ia_\nu^2+\frac45\theta_\R^2+\frac85(atu)^2-\frac{26}{15}(\nu\R x)^2+\frac16Ju_x^2-\frac29i^2u_x^2.
\]
Formula (16.) gives the basis for the reduction of the modulus $(ss'u)^4$, and formula (13.) gives the basis for the reduction of the system of $s$. (\S\ 17.)

\noindent\emph{Consequences:}

1) $a_\nu(j\nu x)^3$, $\theta_\eta H_\vartheta$, $\theta_\eta^2(\vartheta\eta x)$, $\theta_\eta^2(\theta Hu)^2$ are reducers, $a_\nu(j\nu x)^2$, $H_\vartheta$ and $H_\vartheta^2$ decomposable forms.

2) The form of the system II: $j(\nu)=g=(\nu\eta x)^4H_x^2$ is reducible.

3) Since the only contragredient form $g$ becomes reducible, $\nu$ does not enter at all in powers or in products with forms of the same system in folding with System I.

$\theta$ enters into foldings only as a contravariant in products or powers; correspondingly, $H$ enters only as a covariant (cf. \S\ 13 B).

\begin{center}
\textbf{\S\ 12.\quad Forms of the seventh order.}
\end{center}
\[
\text{Folding of }\left\{\begin{array}{l}
\theta\cdot K\text{ with }\nu,\\
K,j,\A\text{ with }H,\\
f\text{ with }L,\sigma.
\end{array}\right.
\]

\noindent\emph{Irreducible forms:}
\[
\text{B)}\quad
\begin{array}{c|c|c|c|c}
\cdot&\cdot&(KHu)^2&\cdot&\cdot\\
\cdot&K_\eta&\cdot&K_\eta(KHu)^2&\cdot\\
(k\eta x)^2&\cdot&K_\eta^2&\cdot&\cdot\\
(k\eta x)^3&H_k(k\eta x)^2,\ K_\eta(k\eta x)^2&\cdot&\cdot&\cdot\\
\cdot&K_\eta(k\eta x)^3&\cdot&\cdot&\cdot
\end{array}
\]
\[
\text{C)}\quad
\begin{array}{cc}
(j\eta x)&H_j\\
\cdot&H_j(j\eta x)\\
\cdot&H_j^2
\end{array}
\qquad
\text{D)}\quad
\begin{array}{c}
(\A Hu)\\\A_\eta
\end{array}
\]
\[
\text{E)}\quad
\begin{array}{cccc}
a_i&(aLu)^2&(aLu)^3&\cdot\\
\cdot&\cdot&a_i(aLu)^2&a_i(aLu)^3\\
\cdot&a_i^2&\cdot&\cdot
\end{array}
\]

\noindent\emph{reductions:}

\noindent A) $(\vartheta\cdot k,\nu,x)^4$ decomposes as single Überschiebung over the decomposable form $(\vartheta^2\nu x)^4$.

\noindent B) form sequence $H_kK_\eta$: reducer $H_\vartheta\theta_\eta$.

form sequence $K_\eta^2(KHu)$: reducer $a_\eta^2(aHu)$.

form sequence $K_\eta^2(k\eta x)$: reducer $\theta_\eta^2(\vartheta\eta x)$.

The from this arising double reduction of the form sequence $K_\eta^3$ gives rise to the reduction of the higher form sequence $(j\eta x)^3$.

$(KHu)$, $(k\eta x)$, $K_\eta(k\eta x)$ decomposable forms according to \S\ 3.

$H_k(KHu)$, $K_\eta(KHu)$ decompose as arising by single folding with the decomposable form $K_\nu(k\nu x)$ and the functional determinant
\[
K_\nu^2=\theta_\nu^2(a\theta u)\equiv a_\nu^2(a\theta u)\modu{\nu^2}.
\]
A relation between decomposable forms corresponds to the double representation of $K_\nu^2$. One obtains:
\[
H_k(KHu)=2K_\nu(\nu\nu_1k)=2\theta_{\nu_1}(\nu\nu_1a\theta)
=2\theta_\nu^2a_\nu-a_\nu^2\theta_\nu+f\theta_\eta(\theta Hu)^2-f\nu\theta_\vartheta^3
\modu{\nu^3},
\]
\[
\begin{aligned}
K_\eta(KHu)&=-K_\nu^2(\nu\nu_1x)
=-a_\nu^2(a\theta u)(\nu\nu_1x)\\
&=a_\eta(aHu)^2\theta+a_\nu^2\theta_\nu
=-\theta_\nu(\theta Hu)^2f-\theta_\nu^2a_\nu.
\end{aligned}
\]
Therefore
\[
\tag{17}
0=a_\nu^2\theta_\nu+\theta_\nu^2a_{\nu_1}+f\theta_\eta(\theta Hu)^2+\theta a_\eta(aHu)^2\modu{\nu^3}.
\]
The forms $H_k$, $H_k(k\eta x)$, $H_k^2$, $H_k^2(k\eta x)$ decompose as arising by single folding with the decomposable forms $H_\vartheta$ and $H_\vartheta^2$ (cf. \S\ 3. 2), a) and b)).

Es is
\[
H_k=H_\vartheta(a\theta u)\sim [H_\vartheta]_{ua}-fH_\vartheta(\theta Hu)
\]
(specialization $y=uv$ of 2)a).
\[
H_k(k\eta x)=H_\vartheta a_\eta-H_\vartheta\theta_\eta f,
\qquad
H_\vartheta a_\eta=\frac54[H_\vartheta]a-\frac14H_\vartheta a_\vartheta
\quad(\S\ 3. 2)b),
\]
\[
H_\vartheta^2(a\theta u)=[H_\vartheta^2]_{ua},
\qquad
H_\vartheta^2a_\eta=[H_\vartheta^2]a=H_k^2(k\eta x)+H_\vartheta^2\theta_\eta f.
\]

\noindent C) form sequence $(j\eta x)^3$: Reducible by double reduction of the form sequence $K_\eta^3$ according to B); according to the ansatz:
\[
0=a_\eta^3(a\theta u)=\theta_\eta^3(a\theta u)
=\theta_\eta+(a\theta\eta x)^3(a\theta u)-\theta_\eta^3(a\theta u)
=\frac12(j\eta x)^3\modu{(\nu^3,s,\R,t,i,J)}.
\]
The analogous ansatz holds up to the terminal form of the form sequence:
\[
0=H_\vartheta^2(a\theta\eta x)a_\eta^3
=H_\vartheta^2(a\theta\eta x)\theta_\eta^3
-\frac12H_\vartheta^2(j\eta x)^4\modu{(\nu^3,s,\R,t,i,J)}.
\]
form $(j\eta x)^2$: decomposes 1) by twofold polarization of the relation for the decomposable form $H_\vartheta^2$ ((12.) b);

2) by direct calculation of the products $a_\nu\theta_\nu^3$; thereby Decomposition of the higher form $(\A Hu)^2$. One obtains:
\[
\tag{18}
\left.\begin{array}{ll}
\text{a)}&\dfrac13(\A Hu)^2+\dfrac13(j\eta x)^2=\theta_\nu^2a_{\nu_1}^2,\\[0.6em]
\text{b)}&(j\eta x)^2=\dfrac43\theta_\nu^3a_{\nu_1},
\end{array}\right\}\modu{(\nu^3,s,\R,t,i,J)}.
\]
according to the ansatz:
\[
H_\vartheta^2(a\theta u)^2-\frac13(\A Hu)^2
=2\theta_\eta^2(a\theta u)(aHu)+\theta_\nu^2a_{\nu_1}^2
=(a\theta u)(aHu)\{a_\eta-(a\theta\eta x)^2\}+\theta_\nu^2a_{\nu_1}^2,
\]
\[
\begin{aligned}
a_{\nu_1}\theta_\nu^3
&=-(a\widehat{\nu}\nu_1u)\theta_\nu^2(\theta\widehat{\nu}\nu_1u)
-\frac12(u\nu\nu_1u)\theta_\nu\theta_{\nu_1}(\vartheta\nu\nu_1u)\\
&=-\frac32\theta_\eta^2(\theta Hu)(aHu)
=-\frac32\theta_\eta^2(\theta Hu)(a\theta u)\\
&=-\frac32\{a_\eta-(a\theta\eta x)^2\}\{(aHu)-(a\theta u)\}(a\theta u).
\end{aligned}
\]

\noindent D) form sequence $\A_\eta(\A Hu)$: reducer $\A_\nu^2$.

$(\A Hu)^2$ decomposable form according to C).

\noindent E) form sequence $a_i^2(aLu)$: reducer $a_\eta^2(aHu)$. forms $(aLu)$ and $a_i(aLu)$: Decomposition as a transvectionen over functional determinants according to \S\ 3.

\noindent F) form sequence $a_\sigma^3$: reducer $a_\R^3$.

forms $a_\sigma^2$ and $a_\sigma^3$: reducer $a_\nu^3$ and $a_\eta^3$ by double reduction of the corresponding expressions
\[
H_\nu a_\nu^3\quad\text{and}\quad H_\nu a_\eta^3,
\qquad
H_\nu a_\nu^3a_\nu\quad\text{and}\quad H_\nu a_\eta^4
\]
(cf. \S\ 10. C) reduction of $\A_\nu^2$ and $\A_\nu^3$).

From this reduction for the higher forms $a_\nu s_\nu(asu)^2$, $a_\nu^2s_\nu(asu)^2$ and for forms in symbolsn $\R,t(a_\nu,a_\eta^2)$ taking account of (16.):
\[
\tag{19}
\left.\begin{array}{rcl}
a_\nu s_\nu(asu)^2&=&\dfrac13 i\theta_\nu^2-\dfrac1{18}iH,\\[0.5em]
a_\nu^2s_\nu(asu)^2&=&0,
\end{array}\right\}\modu{(\R,t)}.
\]
form $a_\sigma$: decomposes by polarization of (12.)b or, more briefly, by direct expansion of the products $a_\nu^2\theta_\nu^3$ according to the ansatz:
\[
a_\nu^2\theta_{\nu_1}^3=(a_\nu^2\theta_{\nu_1})a_{\nu_1}-(a\theta\nu_1x)^2
=-2a_\eta(aHu)(a\theta H)(a\theta\eta x)+(a\theta\nu_1x)^2a_\nu^2\theta_{\nu_1}
\]
and taking account of the reduction of $H_j(j\eta x)^2$ (C):
\[
\tag{20}a_\sigma=2a_\nu^2\theta_{\nu_1}^3\modu{(s,\R,t,i)}.
\]

\noindent\emph{Consequences:}

1) $(j\eta x)^3$, $\A_\eta(\A Hu)$, $a_\sigma^2$ are reducers, $(j\eta x)^2$, $(\A Hu)^2$, $a_\sigma$ decomposable forms.

2) $f$ enters into folding with $L$ only in products with contragredient forms; therefore $f$, and likewise $K$ and $N$, do not enter at all into folding with $\sigma$ and consequently with $(\nu\sigma x)$. The form $j$ enters only in products with contragredient forms; $\A$ does not enter at all in products in folding with $H$ and $L$ (this follows from $\A_\eta(\A Hu)$ and $(\A Hu)^2$). Likewise $\sigma$, and consequently $(\nu\sigma x)$, do not enter in products in folding with any form of System I. Thus, among products in System II, taking account of the consequences of \S\ 11, there remain only powers of $H$ and products of $H$ with $L$.

\begin{center}
\textbf{\S\ 13.\quad Forms of the eighth order (System III).}
\end{center}
\[
\text{Folding of }\left\{\begin{array}{l}
\A\cdot j\text{ with }\nu,\\
\theta^2,
 f\cdot j,
 N\text{ with }H,\\
\theta\text{ with }L,\sigma.
\end{array}\right.
\]

\noindent\emph{Irreducible forms:}
\[
\begin{array}{c@{\quad}l@{\qquad}c@{\quad}l}
\text{A)}&\A_\nu(j\nu x)&
\text{B)}&\begin{array}{c}(\vartheta^2\eta x)^3\\(\vartheta^2\eta x)^4\end{array}
\quad H_\vartheta(\vartheta^2\eta x)^3,
\ \theta_\eta(\vartheta^2\eta x)^3,\\[1em]
\text{C)}&(aHu)H_j,
\quad a_\eta(aHu)H_j&
\text{D)}&H_n,
\ N_\eta,
\end{array}
\]
\[
\text{E)}\quad
\begin{array}{c|c|c|c}
\cdot&(\theta Lu)^2&(\theta Lu)^3&\cdot\\
\theta_i&\cdot&L_\vartheta(\theta Lu)^2,
\ \theta_i(\theta Lu)^2&\theta_i(\theta Lu)^3\\
(\vartheta lx)^2&\theta_i^2&\cdot&\cdot\\
\cdot&\theta_i(\vartheta lx)^2&\cdot&\cdot
\end{array}
\]

\noindent\emph{reductions:}

\noindent A) $\A_\nu(j\nu x)^3$: reducer $a_\nu(j\nu x)^3$.

$\A_\nu(j\nu x)^2$ is reducible by the decomposable form $a_\nu(j\nu x)^2$ according to \S\ 3. 2)b, taking account of the reducer $\theta_{\vartheta j}$. Indeed, according to \S\ 11 B), one obtains:
\[
\frac3{10}\A_\nu(j\nu x)^2+\frac35a_\nu(j\nu x)a_\vartheta(j\nu x)
+\frac1{10}a_\nu(j\nu x)^2
=\frac35\theta_\nu^3\theta_{\vartheta j}+\frac25\theta_\nu^3\theta_{\vartheta j}\theta_{\vartheta j},
\]
(reducer $a_j$ and $\theta_{\vartheta j}$).

\noindent B) The forms $H_\vartheta(\vartheta'\eta x)^2$, $H_\vartheta^2(\vartheta'\eta x)$, $H_\vartheta^2(\vartheta'\eta x)^2$ decompose as arising by successive single folding with the decomposable forms $H_\vartheta$ and $H_\vartheta^2$, respectively $H_\vartheta a_\eta$ and $H_\vartheta^2a_\eta$ according to \S\ 3. 2)b) and according to the relation:
\[
H_\vartheta(\vartheta'\eta x)^2=2H_\vartheta a_\eta^2f-2H_\vartheta a_\eta b_\eta.
\]

\noindent C) form sequence $a_\eta(j\eta x)^2$: reducers $a_j$ and $(j\eta x)^3$.

form $a_\eta(j\eta x)$ decomposes: reducer $a_j$ and single folding with the decomposable form $(j\eta x)^2$ according to \S\ 3, 2)a (specialization $y=uv$).

form $a_\eta H_j$ decomposes: 1) by single folding with the decomposable form $a_\nu(j\nu x)^2$;

2) by a special twofold folding with the decomposable form $(j\eta x)^2$; from this one obtains a relation between decomposable forms.

From $a_\nu(j\nu x)^2$:
\[
0=\theta_\nu'\theta_\nu+2a_\eta H_j\modu{(s,\R,t,i,J)}.
\]
From $(j\eta x)^2$:
\[
0=\theta_\nu'\theta_\nu-\frac32a_\eta H_j\modu{(s,\R,t,i,J)}
\]
(products with contravariants are omitted), according to the ansatz:
\[
0=\frac12a_\nu(abu)^2\theta_\nu^3
+\frac12u_x(ubu)\theta_{\nu_1}^3(\theta bu)
-\frac34(\widehat{j}\eta bu)^2+\frac34H_j(\widehat{j}\eta bu)
-\frac18fH_j^2
\]
and by interchanging the symbols in $a_\nu(ubu)(\theta bu)\theta_{\nu_1}^3$.

\noindent D) form sequence $N_\eta(NHu)$: reducer $\A_\nu^2$.

$(NHu)$, $(n\eta x)$, $N_\eta(n\eta x)$, $H_\eta(NHu)$ decomposable forms (cf. \S\ 12 B), $(NHu)^2$ decomposes by single folding with the form $(\A Hu)^2$.

\noindent E) form sequences $\theta_iL_\vartheta$, $\theta_i^2(\theta Lu)^2$, $\theta_i^2(\vartheta lx)$:

reducers: $\theta_\eta H_\vartheta$, $\theta_\eta^2(\theta Hu)^2$, $\theta_\eta^2(\vartheta\eta x)$.

forms $(\theta Lu)$, $(\vartheta lx)$, $L_\vartheta$, $L_\vartheta(\theta Lu)$, $\theta_i(\theta Lu)$, $L_\vartheta(\vartheta lx)$, $\theta_i(\vartheta lx)$, $L_\vartheta^2$, $L_\vartheta^2(\theta Lu)$, $\theta_i^2(\theta Lu)$ decompose. (Dualistically opposite to the decomposable forms of \S\ 12 B).

\noindent F) form sequence $\theta_\sigma(\vartheta\sigma x)$: reducer $a_\sigma^2$.

forms $(\vartheta\sigma x)$ and $(\vartheta\sigma x)^2$ decompose, as arising by single folding with the decomposable form $a_\sigma$; $\theta_\sigma$, by twofold folding arising, decomposes, since the form sequence
\[
a_{\nu_2}^2(a\theta u)\theta_{\nu_1}^3\equiv H_j^2(j\eta x)\equiv0\modu{u_x(s,\R,t,i,J)}:
\]
\[
\theta_\sigma=a_\sigma(abu)^2=a_\nu^2(abu)^2\theta_\nu^3.
\]

\noindent\emph{Consequences:}

$\theta^3$ or $\theta^2K$ does not enter into folding with $H$; $\theta$ enters into folding with $L$ only as a contravariant in powers or products, and not at all into folding with $\sigma$ and $(\nu\sigma x)$.

$N$ does not enter in products in folding with irgend einer form of System II enters, ebensowenig wie with powers or products of forms of System II. There remain an products in System I, taking account of the consequences of the last paragraphs, only products $f\cdot j$, $\A\cdot j$, powers of $\theta$ and products of $\theta\cdot K$.

\begin{center}
\textbf{\S\ 14.\quad Forms of the ninth order (System III).}
\end{center}
\[
\text{Folding of }\left\{\begin{array}{l}
\theta\cdot K\text{ with }H,\\
K,j,\A\text{ with }L,\\
j,\A\text{ with }\sigma,\\
f\text{ with }H^2.
\end{array}\right.
\]

\noindent\emph{Irreducible forms:}
\[
\begin{array}{c@{\quad}l@{\qquad}c@{\quad}l}
\text{A)}&(\vartheta\cdot k,\eta,x)^4,
\ H_k(\vartheta\cdot k,\eta,x)^4&
\text{B)}&(KLu)^3,
\ K_i(KLu)^3,
\ K_i^2,
\ (klx)^3,
\ K_i(Klx)^3,\\[0.7em]
\text{C)}&L_j,
\ L_j^2&
\text{D)}&\A_i,\\[0.7em]
\text{E)}&(j\sigma x)&
\text{G)}&(aH^2u)^3,
\ (aH^2u)^4,
\ a_\eta(aH^2u)^3.
\end{array}
\]

\noindent\emph{reductions:}

A) the forms $H_\vartheta(k\eta x)^3$, $H_\vartheta^2(k\eta x)^2$, $H_\vartheta^2(k\eta x)^3$ decompose as arising by successive single folding with decomposable forms.

B) form sequences $K_iL_k$, $K_i^2(KLu)$, $K_i^2(klx)$:

reducers: $\theta_\eta H_\vartheta$, $a_\eta^2(aHu)$, $\theta_\eta^2(\vartheta\eta x)$.

forms $L_i$, $K_i(KLu)$, $K_i(klx)$, $(KLu)^2$, $(klx)^2$ decompose as a transvectionen over functional determinants (\S\ 3).

forms $L_k$, $L_k(KLu)$, $L_k(KLu)^2$, $L_k(klx)$, $L_k(klx)^2$, $L_k^2$, $L_k^2(KLu)$, $L_k^2(klx)$, $L_k^3$, $K_i(KLu)^2$, $K_i(klx)^2$ decompose as arising by single folding with the decomposable forms:
\[
L_\vartheta,
\ H_k(KHu),
\ L_\vartheta(\vartheta lx),
\ H_k^2,
\ L_\vartheta^2,
\ L_\vartheta^2(\theta Lu),
\ K_\eta(KHu),
\ \theta_i(\vartheta lx)
\]
according to \S\ 3. 2), a) and b).

C) form sequence $(jlx)^3$: reducer $(j\eta x)^3$.

forms $(jlx)^2$ and $L_j(jlx)$: arising by single folding with the decomposable form $(j\eta x)^2$.

D) form sequence $\A_i(\A Lu)$: reducer $\A_\nu^2$.

forms $(\A Lu)^2$, $(\A Lu)^3$: Single folding with the decomposable form $(\A Hu)^2$.

E) form sequence $(j\sigma x)^2$: reducer $a_\sigma^2$ by replacing $j$ by lower forms of the form sequence (\S\ 5):
\[
a_\sigma^2(a\theta u)^2=\frac13(j\sigma x)^2,
\qquad
a_\sigma^2(a\theta\sigma x)^2=\frac13(j\sigma x)^4.
\]
F) form sequence $\A_\sigma^2$: reducer $a_\sigma^2$.

form $\A_\sigma$: reducer $a_\sigma^2$ by replacing $\A$ by lower forms of the form sequence:
\[
a_\sigma a_\R^2=\frac23\A_\sigma\modu{\nu}.
\]

\noindent\emph{Consequences:} Only the form $j$ enters in folding with $\sigma$ and $(\nu\sigma x)$ enters.

$N$ does not enter in folding with $L$ enters, since the $\A_i$ correspondinglye form $N_i$ decomposes.

\begin{center}
\textbf{\S\ 15.\quad Forms of the tenth order (System III).}
\end{center}
\[
\text{Folding of }\left\{\begin{array}{l}
\theta^2,
\ f\cdot j\text{ with }L,\\
\theta\text{ with }H^2.
\end{array}\right.
\]

\noindent\emph{Irreducible forms:}
\[
\begin{array}{c@{\quad}l@{\qquad}c@{\quad}l}
\text{A)}&(\vartheta^2lx)^4,
\ \theta_i(\vartheta^2lx)^4&
\text{B)}&(aLu)^2L_j,
\ a_i(aLu)^2L_j,\\[0.8em]
\text{C)}&(\theta H^2u)^3,
\ (\theta H^2u)^4,
\ H_\vartheta(\theta H^2u),
\ \theta_\eta(\theta H^2u)^3.
\end{array}
\]

\noindent\emph{reductions:}

A) and B): Decomposition of the remaining forms by single folding with decomposable forms.

C) Decomposition of the remaining forms according to \S\ 13 B) by the dualistically opposite consideration.

\noindent\emph{Consequences:}

$\theta^2K$ does not enter into folding with $L$; correspondingly, $H^2L$ does not enter into folding with $K$. The forms $H^3$ and $H^2L$ do not enter into folding with $\theta$. The folding scheme of \S\ 7 is therefore proved.

\begin{center}
\textbf{\S\ 16.\quad Forms of the 11th, 12th, 13th, and 15th orders (System III).}
\end{center}

\noindent\emph{11th order.}
\[
\text{Folding of }\left\{\begin{array}{l}
\theta\cdot K\text{ with }L,\\
K,j\text{ with }H^2,\\
j\text{ with }(\nu\sigma x),\\
f\text{ with }H\cdot L.
\end{array}\right.
\]
\noindent\emph{Irreducible forms:}
\[
\begin{array}{c@{\quad}l@{\qquad}c@{\quad}l}
\text{A)}&(\vartheta\cdot k,l,x)^5,
\ \theta_i(\vartheta\cdot k,l,x)^5&
\text{B)}&(KH^2u)^4,
\ K_\eta(KH^2u)^4,\\[0.8em]
\text{C)}&(\nu\sigma j),&
\text{D)}&(a_\eta H\cdot L,u)^4.
\end{array}
\]

\noindent\emph{12th order.}
\[
\text{Folding of }\left\{\begin{array}{l}
\theta^3\text{ with }L,\\
\theta\text{ with }H\cdot L.
\end{array}\right.
\]
\noindent\emph{Irreducible forms:}
\[
\text{E)}\ (\vartheta^3lx)^6,
\qquad
\text{F)}\ (\theta,H\cdot L,u)^4,
\qquad L_\vartheta(\theta,H\cdot L,u^4).
\]

\noindent\emph{13th order.}
\[
\text{Folding of }K\text{ with }H\cdot L.
\]
\noindent\emph{Irreducible forms:}
\[
\text{G)}\ (K,H\cdot L,u)^5,
\qquad K_\eta(K,H\cdot L,u)^5.
\]

\noindent\emph{15th order.}
\[
\text{Folding of }K\text{ with }H^3.
\]
\noindent\emph{Irreducible form:}
\[
\text{H)}\ (KH^3u)^6.
\]

\noindent\emph{reductions:}

The forms $H_j^2$, $H_j^2$ have reducer $L_j^3$, from the relation:
\[
(\nu lx)L_x^3=-\frac12H^2\modu{(\sigma,s)}.
\]
Decomposition or reducibility of the remaining forms follows from the considerations of the earlier paragraphs.

\noindent\emph{Consequences:}

According to the folding scheme of \S\ 7 and the consequences of \S\S\ 8--15, the irreducible forms of System III (117 forms) are thereby exhausted.



\clearpage

\begin{center}
{\Large\bfseries Chapter IV.\quad The relatively complete system mod $(\varrho,t)$.}
\end{center}

\begin{center}
\textbf{\S\ 17.\quad Reduction of the modulus $(ss'u)^4$ and of the system of $s$.}
\end{center}

For the formation of the next higher relatively complete system, one must, in analogy with \S\ 7, form the system of $s$ with respect to the next modulus
\[
\nu(s)=(ss'u)^4
\]
and fold it with the forms of the relatively complete system mod $s$. It will be shown that, after introduction of the modulus $(\varrho,t)$ as the higher modulus,
\[
\begin{array}{ll}
\text{A)}&\text{the modulus }(ss'u)^4\text{ becomes reducible,}\vphantom{\dfrac11}\\
\text{B)}&\text{the system of }s\text{ consists of the single form }s.
\end{array}
\]

A) The reduction of the modulus $(ss'u)^4$ follows at once from the reducer $s_\nu^2$ found by formula (16.), \S\ 11. From
\[
 s_\nu^2=\frac43 i a_\nu^2+\frac45\theta_\varrho^2+\frac85(atu)^2-
 \frac{26}{15}(\nu\varrho x)^2+\frac16J\cdot u_x^2-
 \frac29i^2\cdot u_x^2
\]
there follows, by polarization from $x$ to $\nu_1$,
\[
\tag{21}
\begin{aligned}
(ss'u)^4&=-\frac23J\cdot\nu+6s_\nu^2s_{\nu_1}^2\\
&=\frac{24}{5}\theta_\nu^2\theta_\varrho^2+
\frac{48}{5}a_\nu^2(atu)^2-
\frac{52}{5}H_\varrho^2+
\frac43i\cdot(asu)^4+\frac13J\cdot\nu-
\frac49i^2\cdot\nu,
\end{aligned}
\]
and with this the reduction of the modulus $(ss'u)^4$.

\footnotetext{From formula (21.) there follows the relation, mentioned in the note to the introduction, between the three invariants of order 9, with use of formula (10.)c.}

B) The reduction of
\[
Z=(ss'u)^2=\theta(s)
\]
and hence the reduction of the system of $s$ is obtained by replacing the form $Z$ by the lower form $g_\nu^2$, according to formula (8.)b, \S\ 7, taking into account the relation
\[
 s_\eta^2=s_\nu^2(\nu\nu_1x)^2-\frac13Z.
\]
The form $g_\nu^2$ has, by formula (13.), \S\ 11, the reducer $s_\nu^2$ as a factor. By formula (8.) and (16.) one obtains
\[
 g_\nu^2=-Z+\frac{14}{5}ia_\eta^2+\frac{14}{15}i(asu)^2
 \modu{(\varrho,t)},
\]
and by formulas (13.), (14.), (15.), (16.),
\[
\begin{aligned}
g_\nu^2&=2[s_\vartheta^2]_{\nu^2}-\frac43 i\A_\nu^2
      =2s_\vartheta^2s_\nu^2-\frac65[s_\vartheta^2(\theta su)^2]\widehat{\nu x}^{\,2}
        -\frac43i\A_\nu^2\\
&=\frac45i\cdot a_\eta^2-\frac25i\cdot(asu)^2+\frac13J\cdot\theta
\modu{(\varrho,t)}.
\end{aligned}
\]
Thus
\[
\tag{23}
 Z=2i\cdot a_\eta^2+\frac43i\cdot(asu)^2-\frac13J\cdot\theta
 \modu{(\varrho,t)}.
\]

The relatively complete system mod $((ss'u)^4,\varrho,t)$ therefore passes into a relatively complete system mod $(\varrho,t)$, and from this, by transvection over the known system of two quadratic forms, the absolutely complete system arises. In order to form the relatively complete system mod $(\varrho,t)$, since by formula (23.) the system of $s$ reduces to the form $s$, we must transvect the relatively complete system mod $(s,\varrho,t)$ over $s$ and over powers of $s$. It will be shown that powers of $s$ enter into folding only with System I.

The announced relation is
\[
\tag{22}
\begin{aligned}
(ass')^4&=\frac{24}{5}\theta_\nu^2\theta_\varrho^2a_\nu^2a_\varrho^2+
\frac{48}{5}a_\nu^2b_\nu^2(tab)^2-
\frac{52}{5}H_\varrho^2a_\nu^4+
\frac53J\cdot i-\frac49i^3,\\
(ass')^4&=\frac{24}{5}a_\varrho^2a_{\varrho'}^2-
\frac{12}{5}t_\varrho^2+\frac53J\cdot i-\frac49i^3.
\end{aligned}
\]

Among forms of the same order and the same degree, we define as ``higher forms'' those which have arisen from higher forms of the relatively complete system mod $(s,\varrho,t)$ by folding with $s$, regarding the folding with $s$ merely as a polarization of the original forms. Thus the order of the individual forms is uniquely determined (cf. \S\ 1, form sequences 2). For instance,
\[
(\theta_\eta(\theta Hu),s,u)^2\text{ is higher than }s_\eta((\theta Hu)^2,s,u),
\]
\[
(\theta_\eta(\theta Hu)^2,s,u)\text{ is higher than }(\theta_\eta(\theta Hu),s,u)^2.
\]

We divide the resulting system into the four subsystems
\[
\begin{array}{ll}
\text{System }s_I:&\text{obtained by folding }s\text{ and powers of }s\text{ with System I;}\vphantom{\dfrac11}\\
\text{System }s_{II}:&\text{obtained by folding }s\text{ with System II;}\vphantom{\dfrac11}\\
\text{System }s_{III}:&\text{obtained by folding }s\text{ with the individual forms}\\
&\text{(without products) of System III and with products of one}\\
&\text{form each from Systems I and II;}\vphantom{\dfrac11}\\
\text{System }s_{IV}:&\text{obtained by folding }s\text{ with products of one form each}\\
&\text{from Systems I and III, II and III, III and III.}
\end{array}
\]
It should also be noted that from now on all formulas are to be understood mod $(\varrho,t)$; from formula (27.) on, they are to be understood mod $(\varrho,t,i,J,\text{ higher forms})$, without this being expressly stated each time.

For the reduction we collect the formulas obtained earlier:
\[
\tag{24}
\begin{aligned}
s_\nu^2&=\frac43ia_\nu^2+\frac16J\cdot u_x^2-\frac29i^2\cdot u_x^2,
& s_\nu^3&=\frac13J\cdot u_x,
& s_\nu^4&=J,\\[0.3em]
g&=2s_\vartheta^2-\frac43i\A,
& s_\vartheta^2(\theta su)&=0,\\[0.3em]
s_\vartheta^2(\theta su)^2&=\frac89i\cdot a_\nu^3-\frac4{27}i^2\cdot u_x^2,\\[0.3em]
Z&=2i\cdot a_\eta^2+\frac43i(asu)^2-\frac13J\cdot\theta,\\[0.3em]
g_\nu&=-s_\eta+u_x\left\{2ia_\eta^2-\frac23i(asu)^2+\frac23J\cdot\theta\right\},\\[0.3em]
g_\nu^2&=\frac45ia_\eta^2-\frac25i(asu)^2+\frac13J\cdot\theta,
&g_\nu^3&=0,
&g_\nu^4&=0.
\end{aligned}
\]

\noindent\emph{Consequences:}
\[
 s_\nu^3,
\quad s_\vartheta^2(\theta su),
\quad s_\vartheta^2s_\nu,
\quad s_\vartheta^2\theta_\nu
\]
are reducers.

\begin{center}
\textbf{\S\ 18.\quad System $s_I$.}
\end{center}

\[
\text{Folding of}\quad
\left\{\begin{array}{l}
s\text{ with } f;\ \theta;\ K,j,\A;\ f\cdot j,N;\ \A\cdot j,\\
s^3\text{ with }\theta,K.
\end{array}\right.
\]

\noindent\emph{Irreducible forms.}

\noindent Order 5:
\[
(asu),\qquad (asu)^2,\qquad (asu)^3,\qquad (asu)^4.
\]

\noindent Order 6:
\[
\begin{array}{cccc}
(\theta su)&(\theta su)^2&(\theta su)^3&(\theta su)^4\\
s_\vartheta&s_\vartheta(\theta su)&s_\vartheta(\theta su)^2&s_\vartheta(\theta su)^3\\
\cdot&s_\vartheta^2&\cdot&\cdot
\end{array}
\]

\noindent Order 7:
\[
\begin{array}{c@{\qquad}c@{\qquad}c@{\qquad}c}
\cdot&(Ksu)^2&(Ksu)^3&(Ksu)^4\\
\text{A)}\ s_k&s_k(Ksu)&s_k(Ksu)^2&s_k(Ksu)^3\\
\cdot&s_k^2&\cdot&\cdot
\end{array}
\qquad
\begin{array}{l}
\text{B)}\ s_j,\\[0.4em]
\text{C)}\ (\A su),\ (\A su)^2.
\end{array}
\]

\noindent Order 8:
\[
\begin{array}{ll}
\text{A)}&s_j(asu),\ s_j(asu)^2,\ s_j(asu)^3,\\[0.4em]
\text{B)}&(Nsu)^2,\quad s_n,\quad s_n(Nsu).
\end{array}
\]

\noindent Order 10:
\[
\text{A)}\ s_j(\A su),\qquad \text{B)}\ s_\vartheta(\theta s'u)^4.
\]

\noindent Order 11:
\[
(Ks^2u)^6,\qquad s_k(Ks^2u)^5,
\qquad s_k(Ks^2u)^6.
\]

\noindent\emph{Reductions:}

The form sequences
\[
 s_\vartheta^2(\theta su),\quad s_k^2(Ksu),\quad s_j^2,\quad (\A su)^3,
 \quad (Nsu)^3
\]
have the reducer $s_\vartheta^2(\theta su)$ as a factor; $s_j^2$ and $(\A su)^3$ are obtained by going back from $j$ and $\A$ to lower forms of the form sequence:
\[
\begin{aligned}
s_\vartheta^2(\theta su)(a\theta u)^3&=-\frac12s_j^2,\\
s_\vartheta^2(\theta su)(a\theta u)^2&=\frac13(\A su)^3,\\
s_\vartheta^2(\theta su)^2(a\theta u)^2&=\frac13(\A su)^4+\frac13s_j^2.
\end{aligned}
\]
Forms $s_\vartheta^2s_{\vartheta'}$ and $s_\vartheta^2s_{\vartheta'}^2$: reducers $s_\vartheta^2(\theta su)$ and $\theta_{\vartheta'}$:
\[
\begin{aligned}
s_\vartheta^2s_{\vartheta'}&=s_\vartheta^2\theta_{\vartheta'}-s_\vartheta^2(\theta s\widehat{\vartheta'}x),\\
s_\vartheta^2s_{\vartheta'}^2&=s_\vartheta^2s_{\vartheta'}\theta_{\vartheta'}-s_\vartheta^2s_{\vartheta'}(\theta s\widehat{\vartheta'}x).
\end{aligned}
\]
Form sequence $s_j(\A su)^2$: reducers $\A_j$ and $(\A su)^3$.

\noindent\emph{Consequences:}

$s$ does not enter into powers in folding with $f,j,\A,N$. The forms $f,j,\A,N$ enter into folding only in products with contragredient forms; however, in products with $f$ those forms may still be quadratic in $x$, and in products with $\A$ they may still be linear in $x$. Thus the following products of forms are to be considered:
\[
 f\cdot(0,\lambda),\quad f\cdot(1,\lambda),\quad f\cdot(2,\lambda),\quad
 j\cdot(\mu,0),\quad N,\quad
 \A\cdot(0,\lambda),\quad \A\cdot(1,\lambda)\quad(\lambda>0),
\]
where $(\mu,\lambda)$ denotes a form of degree $\mu$ in $x$ and of degree $\lambda$ in $u$.

Neither $\theta$ nor $K$ enters into powers or products with arbitrary forms in folding with $s$ or $s^2$. For products with the functional determinant $K$, $K\cdot\varphi_x$ ($\varphi\ne f$ or $\theta$), the relation between decomposing forms is
\[
 K\cdot\varphi_x=(a\theta u)\varphi_x=f\cdot(\varphi\theta u)-\theta\cdot(\varphi au).
\]
The above scheme for folding System $s_I$ is thereby proved.

\begin{center}
\textbf{\S\ 19.\quad System $s_{II}$.}
\end{center}

Folding of $s$ with $\nu;\ H;\ L;\ H^2$.

\noindent\emph{Irreducible forms:}
\[
\begin{array}{c@{\qquad}c@{\qquad}c@{\qquad}c}
\text{Order 6} & \text{Order 8} & \text{Order 10} & \text{Order 12}\\[0.3em]
s_\nu & (Hsu),\ (Hsu)^2 & (Lsu)^2,\ (Lsu)^3 & (H^2su)^3\\
\cdot & s_\eta & s_l & \cdot\\
\cdot & s_\nu^4=J & \cdot & \cdot
\end{array}
\]

\noindent\emph{Reductions:}

Form sequences $s_\eta(Hsu)$ and $s_l(Lsu)$: reducer $s_\nu^2$.

Form sequence $s_\sigma$: reducer $s_\nu^2$ from $H$, $s_\nu^2=\dfrac23s_\sigma$.

$(H^2su)^4$ decomposes by formula (7.)a (cf. \S\ 11 A). Hence $(H\cdot L,s,u)^4$ also decomposes.

\noindent\emph{Consequences:}

$s^2$ does not enter into folding with System II. The form $\nu$ enters only in products with contragredient forms. The form $H$, regarded as a covariant, enters into folding in products with forms $(1,\lambda)$, $(2,\lambda)$, $(3,\lambda)$, with arbitrary $\lambda$. The forms $\sigma$ and $(\nu\sigma x)$ do not enter at all, and $L$, as a functional determinant, does not enter into products in folding; the full transvections over covariants of System III become reducible. The products of Systems I and II that remain are
\[
 f\cdot\nu,\qquad \A\cdot\nu.
\]

\begin{center}
\textbf{\S\ 20.\quad Forms of order 7 (System $s_{III}$).}
\end{center}

Folding of $s$ with $f\cdot\nu$, $a_\nu$, $a_\nu^2$.

\noindent\emph{Irreducible forms:}
\[
\begin{gathered}
s_\nu(asu),\quad a_\nu(asu),\quad s_\nu(asu)^2,\quad a_\nu(asu)^2,
\quad s_\nu(asu)^3,\quad a_\nu(asu)^3,\\
a_\nu s_\nu,
\quad a_\nu s_\nu(asu),
\quad a_\nu^2(asu),
\quad a_\nu^2(asu)^2.
\end{gathered}
\]

\noindent\emph{Reductions:}

The evident reductions that arise by direct folding with the reducers $s_\nu^2$ and $s_\eta(Hsu)$ will not be mentioned, nor will the decomposition of transvections with functional determinants.

For the forms $a_\nu s_\nu(asu)^2$ and $a_\nu^2s_\nu(asu)^2$, see formula (19.), \S\ 12; moreover,
\[
 a_\nu s_\nu(asu)^3=a_\nu(a\widehat{\nu}_1\nu_2u)^3(\nu_1\nu_2\nu)=0.
\]
Form sequence $a_\nu^2s_\nu$: reducer $s_\vartheta^2(\theta su)$, obtained by going back from $a_\nu^2$ to lower forms, i.e. by double reduction of the expressions
\[
 s_\vartheta^2(a\theta s)(a\theta u),
 \qquad s_\vartheta^2(a\theta s)(a\theta u)^2
\]
according to the Ansatz
\[
\begin{aligned}
s_\vartheta^2(a\theta s)(a\theta u)
&=[(a\theta u)^2]s^3+\frac15[a_\vartheta(a\theta u)^2]s^2+\frac1{60}[a_\vartheta^2(a\theta u)^2]s\\
&=\frac12a_\nu^2s_\nu-\frac23a_\nu^2s_\nu^2,\\
s_\vartheta^2(a\theta s)(a\theta u)^2
&=\frac45[a_\vartheta(a\theta u)^2]_{us}s^2+\frac{23}{180}[a_\vartheta^2(a\theta u)^2]_{us}s\\
&=-\frac12a_\nu^2s_\nu(asu).
\end{aligned}
\]
Here $a_\nu^2b_\nu(abu)=0$.

The formulas obtained are
\[
\tag{25}
\begin{array}{rcl@{\qquad}rcl}
a_\nu^2s_\nu&=&\dfrac49i\,a_\nu^2b_\nu,
& a_\nu s_\nu(asu)^2&=&\dfrac13i\theta_\nu^2-\dfrac1{18}iH,\\[0.7em]
a_\nu s_\nu(asu)^3&=&0,
& a_\nu^2s_\nu(asu)&=&0,\\[0.7em]
a_\nu^2s_\nu(asu)^2&=&0.
\end{array}
\]

\noindent\emph{Consequences:}

1) $a_\nu s_\nu(asu)^2$ and $a_\nu^2s_\nu$ are reducers.

2) The values of the forms $(\A su)^3$, $(\A su)^4$, already recognized as reducible in \S\ 19, are thereby obtained; likewise for the corresponding form sequence $(agu)^2$, which, when folded with $\nu$, gives rise to double reduction through the reducer $g_\nu$:
\[
\tag{26}
\begin{aligned}
\text{a)}\quad (\A su)^3&=-\frac65a_\nu^2(asu)+3a_\nu s_\nu(asu)+2i\theta_\nu(\theta\nu x),\\
\text{b)}\quad (\A su)^4&=-\frac25a_\nu^2(asu)^2+\frac43i\theta_\nu^2+\frac19iH.
\end{aligned}
\]
From formula (24.) and (26.), mod $(i,J)$ (cf. p. 71),
\[
\tag{27}
\begin{aligned}
\text{a)}\quad (agu)^2&=\frac23(\A su)^2-\frac43a_\nu s_\nu+\frac35s\cdot a_\nu^2,\\
\text{b)}\quad (agu)^3&=2a_\nu s_\nu(asu)-a_\nu^2(asu)^2,\\
\text{c)}\quad (agu)^4&=-a_\nu^2(asu)^2.
\end{aligned}
\]


\begin{center}
\textbf{\S\ 21.\quad Forms of 8th order (System $s_{III}$).}
\end{center}

Folding of $s$ with the forms of 4th order (System III). (\S\ 9.)

\noindent\emph{Irreducible forms:}
\[
\begin{array}{lll}
(\vartheta\nu x)s_\nu,& ((\vartheta\nu x),s,u)^2,\\[-0.1em]
((\vartheta\nu x)^2,s,u),& ((\vartheta\nu x)^2,s,u)^2,\ \theta s_\nu,\\[-0.1em]
(\vartheta\nu x)^2s_\nu,& ((\vartheta\nu x)^2,s,u)s_\nu,
\end{array}
\]
\[
\begin{array}{lll}
((\vartheta\nu x),s,u)^3,\ (\theta_\nu su)^2,&
((\vartheta\nu x),s,u)^4,\ (\theta_\nu su)^3,\\[-0.1em]
((\vartheta\nu x)^2,s,u)^3,\ (\theta_\nu(\vartheta\nu x),s,u)^2,\ (\theta_\nu^2su),&
((\vartheta\nu x),s,u)^3,\ (\theta_\nu^2su)^2,\ (\theta_\nu(\vartheta\nu x),s,u)^4.
\end{array}
\]

\noindent\emph{Reductions:}

$((\vartheta\nu x),s,u)s_\nu$ decomposes as a transvection over functional determinants, by \S\ 3.

Form sequence $((\vartheta\nu x),s,u)^2s_\nu$: reducers $s_\nu^2$ and $s_\vartheta^2\theta_\nu$:
\[
(\widehat\vartheta\nu su)s_\nu(\theta su)=s_\nu s_\vartheta(\theta su)=-\frac12(\widehat\vartheta\nu su)^2(\theta su),
\]
\[
\theta_\nu(\widehat\vartheta\nu su)s_\nu=\theta_\nu s_\nu s_\vartheta=-\frac12\theta_\nu(\widehat\vartheta\nu su)^2.
\]
Form sequence $\theta_\nu s_\nu(\theta su)$: reducer $a_\nu s_\nu(asu)^2$ by double reduction of the expressions
\[
a_\nu s_\nu(asu)^2(abu)\quad\text{and}\quad a_\nu s_\nu(asu)^2(abu)(sbu);
\]
$\theta_\nu s_\nu(\theta su)^3$ by double reduction of $(agu)^3(abu)(bgu)^3$.

Form sequence $\theta_\nu(\vartheta\nu x)s_\nu$: reducer $a_\nu^2s_\nu$ from $\theta_\nu(\vartheta\nu x)s_\nu=a_\nu^2(abu)s_\nu$.

Form sequence $(\theta_\nu(\vartheta\nu x)^2,s,u)$: double reduction of the forms of the sequence $\theta_\nu(\widehat\vartheta\nu su)s_\nu$, which at the same time belong to the form sequences $((\vartheta\nu x)su)^2s_\nu$ and $\theta_\nu(\vartheta\nu x)s_\nu$.

The form $(\theta_\nu(\vartheta\nu x),s,u)$ decomposes as a single transvection over the functional determinant $\theta_\nu(\vartheta\nu x)=a_\nu^2(abu)$.

The form $((\vartheta\nu x)^2,s,u)^4$: double reduction of the expression $(a\A u)^2(\A su)^4$, or also of $(\theta gu)^4$, from formulas (27.) and analogously to the reduction of $(j\eta x)^4$ (\S\ 12 C)).

Double reduction gives the formulas:
\[
\tag{28}
\begin{aligned}
0&=\theta_\nu s_\nu(\theta su)-\frac16(\widehat\vartheta\nu su)^2(\theta su)-\frac1{12}(Hsu),\\
0&=\theta_\nu s(\theta su)^2-\frac14(\widehat\vartheta\nu su)^2(\theta su)^2-\frac1{24}(Hsu)^2,\\
0&=(\widehat\vartheta\nu su)^2(\theta su)^2+2\theta_\nu(\widehat\vartheta\nu su)(\theta su)^2
  +3\theta_\nu^2(\theta su)^2-\frac13H(su)^2,\\
0&=\theta_\nu s_\nu(\theta su)^3+\frac12\theta_\nu(\widehat\vartheta\nu su)(\theta su)^3,\\
0&=\theta_\nu s_\nu s_\vartheta-\frac14s_\eta,
\quad 0=\theta_\nu s_\nu s_\vartheta(\theta su),
\quad 0=\theta_\nu s_\nu s_\vartheta(\theta su)^2.
\end{aligned}
\]
The last group corresponds to $((\theta_\nu(\vartheta\nu x)^2,s,u)^2,\ldots)$.

\noindent\emph{Consequences:}

1) $\theta(\vartheta\nu x)s_\nu$, $(\widehat\vartheta\nu su)(\theta su)s_\nu$, $\theta_\nu(\widehat\vartheta\nu su)^2$, $(\widehat\vartheta\nu su)^2(\theta su)^2$ are reducers.

2) The forms of 4th order $(5,\lambda)$, $(6,\lambda)$ etc. do not enter into folding with $s^2$.

3) For the form sequence $(\theta gu)^2$, formulas (27.), with the reductions of this paragraph taken into account, give the following formulas:
\[
\tag{29}
\begin{aligned}
\text{a)}\quad (\theta gu)^2&=\frac43\theta_\nu s_\nu+\frac23s_\nu s_\vartheta-\frac13s\theta_\nu^2-\frac19sH-\frac13f\,a_\nu^2(asu)^2+\frac13a_\nu^2(asu)^2,\\
\text{b)}\quad (\theta gu)^3&=-\frac13(\widehat\vartheta\nu su)^2(\theta su)-\theta_\nu(\widehat\vartheta\nu su)(\theta su)-\theta_\nu^2(\theta su),\\
\text{c)}\quad (\theta gu)^4&=2\theta_\nu^2(\theta su)^2-\frac13(Hsu)^2,\\
g\vartheta(\theta gu)&=(\widehat\vartheta\nu su)(\vartheta\nu x)s_\nu-\theta_\nu(\vartheta\nu x)(\widehat\vartheta\nu su),\\
g\vartheta(\theta gu)^2&=\frac23s\theta_\nu^3,
\quad g\vartheta(\theta gu)^3=\frac13\theta_\nu^3(\theta su),
\quad g\vartheta(\theta gu)^4=0.
\end{aligned}
\]

\begin{center}
\textbf{\S\ 22.\quad Forms of 9th order (System $s_{III}$).}
\end{center}

Folding of $s$ with the forms of 5th order (System III) and the product $\A\cdot\nu$ (\S\ 10).

\noindent\emph{Irreducible forms:}
\[
\begin{array}{ll}
\text{A)}& (k\nu x)^2s_\nu,
((k\nu x)^3,s,u),
(k\nu x)^3s_\nu,
((k\nu x)^2,s,u)^2,
K_\nu s_\nu,\\
& ((k\nu x)^3,s,u)^2,
((k\nu x)^2,s,u)s_\nu,
(K_\nu(k\nu x)^3,s,u),\\
& (K_\nu su)^2,\ (K_\nu su)^3,\ (K_\nu su)^4,
((k\nu x)^2,s,u)^3,
(K_\nu^2su)^4,\\[0.3em]
\text{B)}& (j\nu x)s_\nu,
((j\nu x)^2,s,u),
(j\nu x)^3,s,u,
((j\nu x)^2,s,u)^2,\\[0.3em]
\text{C)}& s_\nu(\A su),\quad (\A_\nu su),\\[0.3em]
\text{D)}& (aHu)s_\eta,
(a_\eta su),
((aHu),s,u)^2,
((aHu)^2,s,u),\\
& (aHu)^2s_\eta,
(a_\eta su)^2,
(a_\eta^2su),
((aHu),s,u)^3,
((aHu)^2,s,u)^2,\\
& ((aHu),s,u)^4,
(a_\eta su)^3,
(a_\eta(aHu),s,u)^2,
(a_\eta(aHu)^2,s,u),
(a_\eta(aHu),s,u)^3.
\end{array}
\]

\noindent\emph{Reductions:}

A) The reductions follow directly from the reductions of \S\ 21 by single folding. The form sequence $K_\nu s_\nu(Ksu)^2$ has the reducers $a_\nu s_\nu(asu)^2$ and $\theta_\nu s_\nu(\theta su)^2$ as factors. Hence the higher forms $K_\nu^2(Ksu)^2$, $K_\nu^2(Ksu)^3$ decompose according to the Ansatz:
\[
0=a_\nu s_\nu(asu)^2(a\theta u)=K_\nu s_\nu(Ksu)^2-\theta_\nu s_\nu(\theta su)^2(a\theta u).
\]

B) Form sequence $(j\nu x)^2s_\nu$: reducer $a_\nu^2s_\nu$ from
\[
(j\nu x)^2s_\nu=3a_\nu^2s_\nu(a\theta u)^2.
\]
$((j\nu x)^3,s,u)^2$: reducer $a_\nu^2s_\nu$ from
\[
((j\nu x)^3,s,u)^2=-6a_\nu^2s_\nu(a\theta u)^2(\theta su).
\]

C) Forms $\A_\nu(\A su)^2$ and $s_\nu(\A su)^2$: reducer $g_\nu$ from formula (27.)a.

Form $\A_\nu s_\nu$:
1) reducer $a_\nu^2s_\nu$ from $a_g a_\nu^2s_\nu=\frac23\A_\nu s_\nu$;
2) decomposes by formula (27.)a through double reduction of the expression $(\A s\widehat\nu x)^2$.

Hence the higher form $a_\eta s_\eta$ decomposes.

D) Form sequence $(aHu)^2s_\eta(asu)$: reducer $a_\nu s_\nu(asu)^2$ by double reduction of the form sequence $[a_\nu s_\nu(asu)^2]_{\nu_1x^2}$; reduction of $(aHu)^2s_\eta(asu)^2$ from $a_\nu s_\nu(asu)^2$ and $a_\nu s_\nu(asu)^3$. Hence the reduction of $(a_\eta,s,u)^4$.

Form sequence $a_\eta(aHu)s_\eta$: reducers $a_\nu^2s_\nu$, $a_\eta^2s_\eta$ by double reduction of the expression $(\A s\widehat\nu x)^3$, since $a_\eta^2s_\eta$ does not arise by folding from the reducible term $a_\nu^2s_\nu(\nu\nu_1x)$ of the form $a_\eta(aHu)s_\eta$.

Form sequence $(a_\eta^2su)^2$: reducer $a_\nu^2s_\nu$ from
\[
a_\nu^2s_\nu s_{\nu_1}=-\frac12(a_\eta^2(Hsu))^2.
\]
The form $(a_\eta(aHu),s,u)$ decomposes as a transvection over a functional determinant.

The reduction formulas obtained are:
\[
\tag{30}
\begin{aligned}
\text{a)}\quad 0&=(aHu)^2(asu)s_\eta-a_\eta(asu)(Hsu)^2-a_\eta(aHu)(asu)(Hsu),\\
\text{b)}\quad 0&=(aHu)^2(asu)^2s_\eta-a_\eta(aHu)(asu)^2(Hsu),\\
\text{c)}\quad 0&=a_\eta(asu)^2(Hsu)^2,\\
0&=a_\eta s_\eta+\frac14a_\nu^2g-\frac14s\cdot a_\eta^2,\\
0&=a_\eta(aHu)s_\eta-\frac12a_\eta^2(Hsu),
\quad 0=a_\eta^2(Hsu)^2,\ldots,
\quad 0=a_\eta^2s_\eta.
\end{aligned}
\]

\noindent\emph{Consequences:}

1) $(j\nu x)^2s_\nu$, $\A_\nu s_\nu$, $\A_\nu(\A su)^2$, and consequently $N_\nu(Nsu)^2$, $(aHu)^2(asu)s_\eta$, $a_\eta(aHu)s_\eta$, $a_\eta^2(Hsu)^2$, are reducers; $a_\eta s_\eta$ is a decomposable form that passes into decomposable forms under all single and double foldings.

2) The forms of 5th order $(5,\lambda)$, $(6,\lambda)$ etc. do not enter into folding with $s^2$.

\begin{center}\textbf{\S\ 23.\quad Forms of the tenth order (System $s_{III}$).}\end{center}
Folding of $s$ with the forms of the sixth order (System III) (\S\ 11).

\noindent\emph{Irreducible forms:}
\[
\begin{array}{rl}
\text{A)}& (\vartheta^2\nu x)^3s_\nu,
((\vartheta^2\nu x)^3,s,u),
((\vartheta^2\nu x)^3,s,u)^2,
(\vartheta^2\nu x)^3s_\nu s_\vartheta,
\theta_\nu(\vartheta^2\nu x)^3s_\vartheta,\\[0.3em]
\text{B)}& a_\nu(j\nu x)s_\nu,
(a_\nu(j\nu x),s,u)^2,
(a_\nu(j\nu x),s,u)^3,
(a_\nu(j\nu x),s,u)^4,
(a_\nu^2(j\nu x),s,u)^2,
(a_\nu^2(j\nu x),s,u)^3,\\[0.3em]
\text{C)}& ((\vartheta\eta x)^2,s,u),
((\vartheta\eta x),s,u)^2,
(\theta Hu)s_\eta,
(\theta_\eta su),
((\theta Hu),s,u)^2,
((\theta Hu)^2,s,u),\\
& ((\theta\eta x),s,u)^3,
(\theta_\eta su)^2,
(\theta_\eta^2su),
((\theta Hu),s,u)^3,
((\theta Hu)^2,s,u)^2,
(H_\vartheta(\theta Hu),s,u)^2,\\
& (\theta_\eta(\theta Hu),s,u)^2,
(\theta_\eta(\theta Hu)^2,s,u),
((\theta Hu),s,u)^4,
(\theta_\eta su)^4,
(\theta_\eta(\theta Hu),s,u)^3,
(\theta_\eta^2(\theta Hu),s,u)^2.
\end{array}
\]
\noindent\emph{Reductions:}
A) The forms $((\vartheta^2\nu x)^3,s,u)^2$, $((\vartheta^2\nu x)^3,s,u)^3$ decompose:
\[
(\vartheta\nu x)(\widehat{\vartheta\nu su})(\widehat{\vartheta\nu su})=(\vartheta\nu x)s_\vartheta s_\nu-(\vartheta\nu x)s_\nu s_\vartheta=-2\theta(\vartheta\nu x)s_\nu s_\vartheta+(\vartheta\nu x)s_\vartheta^2.
\]
The remaining forms either have reducers as factors or are single foldings with decomposed forms.

B) Reducer $a_\nu^2s_\nu$ and $(\widehat{j\nu su})s_\nu$.

C) 1. The form sequence $(\theta Hu)^2s_\eta$: a) reducer $(aHu)^2(asu)s_\eta$; b) double reduction of the form sequences $(\theta gu)^3g_\nu$ and $g_\nu(agu)^3(bgu)^3(abu)$. Hence reduction of the higher forms $(\theta,s,u)^3$, $(H_\vartheta(\theta Hu),s,u)^3$ and $(a_\nu b_{\nu_1}^2,s,u)^4$ [the latter form replacing $(H_\vartheta su)^4$].

2. $(\vartheta\eta x,s,u)^4$: reducer $(a_\eta su)^4$.

3. $((\vartheta\eta x)^2,s,u)^3$: reducer $s_\vartheta^2s_\nu$ from $(\widehat{\vartheta\eta su})^2(Hsu)$. The forms $(\vartheta\eta x)s_\eta$, $(\vartheta\eta x)^2s_\eta$, $((\vartheta\eta x)^2,s,u)^2$, $\theta_\eta s_\eta$ decompose by folding with the decomposed form $a_\eta s_\eta$.

4. Form sequences $H_\vartheta(\theta Hu)s_\eta$, $\theta_\eta(\theta Hu)s_\eta$, $H_\vartheta(\vartheta\eta x)s_\eta$, $\theta_\eta(\vartheta\eta x)s_\eta$: reducers $a_\eta(aHu)s_\eta$ and $a_\eta^2s_\eta$.

The form sequence $(\theta_\eta(\vartheta\eta x),s,u)^2$: reducer $a_\eta^2(Hsu)^2$ [except for $(\theta_\eta^2(\theta Hu),s,u)^2$, which arises from the term $a_\eta^2(\theta su)(Hsu)$ by folding].

5. $(H_\vartheta(\vartheta\eta x),s,u)$, $(H_\vartheta(\vartheta\eta x),s,u)^2$: double reduction of $a_\eta(aHu)s_\eta(abu)$ and $g_\vartheta(\theta gu)g_\nu$.

The form sequence $(H_\vartheta(\vartheta\eta x),s,u)^3$: reducer $((\vartheta\nu x)^2,s,u)^2s_\nu$.

6. $(H_\vartheta(\theta Hu),s,u)$ decomposes by single folding with the decomposed form $(\theta_\nu(\vartheta\nu x),s,u)$ according to \S\ 3. 2), c), that is, by double reduction of the lower decomposed form $(H_\vartheta su)^2$:\footnote{The sign $\equiv$ means that products of forms of lower degree have been omitted.}
\begin{align*}
0&=\theta_\nu(\vartheta\nu_1)(\theta su)+\frac23\theta_\nu^2(\vartheta\nu_1x)(\theta su)+\theta_\nu(\vartheta\nu x)(\theta su)s_{\nu_1}\\
&\equiv -\frac12H_\vartheta(\theta Hu)(\theta su)+\theta_\eta(\theta su)(Hsu)+\frac12H_\vartheta(\theta su)(Hsu)\\
&\equiv -\frac12H_\vartheta(\theta Hu)(\theta su)+\theta_\eta(\theta su)(Hsu)+\frac12[H_\vartheta]_{sx}^2-\frac12\frac35H_\vartheta(\theta Hu)(\theta su)\\
&\equiv -\frac35H_\vartheta(\theta Hu)(\theta su).
\end{align*}
By double reduction, according to 1. and 5., the following formulae arise:
\begin{equation}
\tag{31}
\begin{aligned}
0&=\theta_\eta(\theta su)(Hsu)^2+\frac14H_\vartheta(\theta Hu)(\theta su)^2\\
&\quad-\frac34\theta_\eta(\theta Hu)(\theta su)(Hsu)-\frac54\theta_\eta(\theta Hu)^2(\theta su)^3\\
&\quad-(asu)^3b_\eta(\theta Hu)^2-a_\nu(asu)^3b_\nu^2,\\
0&=H_\vartheta(\theta Hu)(\theta su)^3-7\theta_\eta(\theta Hu)(\theta su)^2(Hsu),\\
0&=a_\nu(asu)^2b_{\nu_1}^{2}(bsu)^2-\theta_\eta(\theta su)^2(Hsu)^2\\
&\quad+6\theta_\eta(\theta Hu)(\theta su)^2(Hsu),\\
0&\equiv (H_\vartheta(\vartheta\eta x),s,u),
&0&\equiv H_\vartheta(\widehat{\vartheta\eta su})(Hsu)-4\theta_\eta^2(Hsu).
\end{aligned}
\end{equation}
\noindent\emph{Consequences:}
\begin{enumerate}
\item $(\theta Hu)^2s_\eta$, $H_\vartheta(\theta Hu)s_\eta$, $\theta_\eta(\theta Hu)s_\eta$, $H_\vartheta(\vartheta\eta x)s_\eta$, $\theta_\eta(\vartheta\eta x)s_\eta$, $(H_\vartheta(\vartheta\eta x),s,u)$, $(\theta_\eta(\vartheta\eta x),s,u)^2$ are reducers.
\item $s^2$ does not enter into folding with the forms $(5,\lambda)$, etc., of order 6.
\end{enumerate}

\begin{center}\textbf{\S\ 24.\quad Forms of the 11th order (System $s_{III}$).}\end{center}
Folding of $s$ with the forms of the 7th order (System III) (\S\ 12).

\noindent\emph{Irreducible forms:}
\[
\begin{array}{rl}
\text{A)}& ((k\eta x)^3,s,u),\quad (K_\eta(k\eta x)^3,s,u),\quad (K_\eta su)^2,
\quad ((KHu)^2su)^2,\\
& ((KHu)^2,s,u)^3,
\quad ((KHu)^2,s,u)^4,
\quad (K_\eta(KHu)^2,s,u)^3,\\[0.3em]
\text{B)}& ((j\eta x),s,u)^2,
\quad (H_jsu),
\quad ((j\eta x),s,u)^3,\\[0.3em]
\text{C)}& (\D_\eta su),\\[0.3em]
\text{D)}& (aLu)^2s_\eta,
\quad (a_\eta su)^2,
\quad ((aLu)^2,s,u)^2,
\quad ((aLu)^3,s,u)^2,
\quad ((aLu)^3,s,u)^3,
\quad (a_l(aLu)^2,s,u)^2.
\end{array}
\]
\noindent\emph{Reductions:}

A) Reduction or decomposition by single folding with the corresponding forms in the symbols $\theta-H$.

$(K_\eta(KHu)^2,s,u)$ and $(K_\eta(KHu)^2,s,u)^2$: decomposition according to \S\ 3. 2), a), by folding with $K_\nu^2(Ksu)$, $K_\nu^2(Ksu)^2$.

$(K_\eta(KHu),s,u)^4\equiv(a_\nu^2\theta_{\nu_1},s,u)^4$: decomposition according to \S\ 3. 2), b), from the decomposed form $K_\nu^2(Ksu)^3$.

B) Form sequence $(j\eta x)s_\eta$: reducer $(j\nu x)^2s_\nu$.

Form $H_j(\widehat{j\eta su})(Hsu)$: reducer $(j\nu x)(\widehat{j\nu su})^2$.

Form $H_j(j\eta x)(Hsu)$: decomposes by reducing the decomposed form $((j\eta x)^2,s,u)^2$ by the reducer $(j\nu x)^2s_\nu$ (formula (18.)a):
\[
0=-2(j\nu x)^2s_\nu s_{\nu_1}=[(j\eta x)^2]_{su^2}+H_j(j\eta x)(Hsu).
\]
C) Reducers $\D_\nu s_\nu$ and $\D_\nu(\D su)^2$.

D) Reduction and decomposition by single folding with the corresponding forms in the symbols $f-H$.

E) From (30.)c one obtains the reduction of the decomposed form sequence $(a_\sigma su)^2$:
\begin{align*}
0&=a_\eta(Hsu)^2(as\widehat{\nu x})^2
=a_\eta(Hsu)^2(a\widehat{\nu\eta}u)^2+2a_\eta(a\widehat{\nu\eta}u)(\widehat{\eta\sigma u})(Hsu)^2
=\frac13a_\sigma(asu)^2,\\
0&=a_\eta a_\nu(Hsu)^2(as\widehat{\nu x})(asu)
=a_\eta(a\widehat{\nu\eta}u)^2(Hsu)^2(asu)+a_\eta(a\widehat{\nu\eta}u)(\widehat{\eta\sigma u})(Hsu)^2(asu)
=\frac13a_\sigma(asu)^3.
\end{align*}
\noindent\emph{Consequence:} $s^2$ does not enter into folding with the forms of order 7.

\begin{center}\textbf{\S\ 25.\quad Forms of the 12th, 13th, 14th, and 15th orders (System $s_{III}$).}\end{center}
Folding of $s$ with the forms of orders 8, 9, 10, 11 (System III) (\S\S\ 13--16).

\noindent\emph{Irreducible forms:}

\noindent 12th order.
\[
\begin{array}{rl}
\text{A)}& ((\vartheta^2\eta x)^4,s,u),\quad \theta_\eta(\vartheta^2\eta x)^3s_\vartheta,\\
\text{B)}& ((aHu)H_j,s,u)^2,
\quad ((aH\cdot u)H_j,s,u)^3,\\
\text{C)}& (\theta_\nu su)^2,
\quad ((\theta Lu)^2,s,u)^2,
\quad ((\theta Lu)^3,s,u),
\quad ((\theta Lu)^2,s,u)^3,
\quad (\theta_l(\theta Lu)^2,s,u)^2.
\end{array}
\]
\noindent\emph{Reductions:} The reductions arising by direct folding with reducers or decomposed forms are no longer to be mentioned.

B) Decomposition of $((aHu)H_j,s,u)$ and $(a_\eta(aHu)H_j,s,u)$ as transvections over functional determinants.

Decomposition of $(a_\eta(aHu)H_j,s,u)^2$: double reduction of the decomposed form $[\theta_\nu\theta_{\nu_1}^3]_{su^3}$ (\S\ 13. C):
\begin{align*}
0&\equiv[\theta_{\nu_1}^3\theta_\nu]_{su^3}
\equiv-\frac34\theta_\nu(\theta su)^2(\theta\theta'u)\theta_{\nu_1}^3\\
&\equiv-\frac98(\theta\theta'\eta su)(\theta\theta'u)(\theta Hu)(\theta'Hu)\theta_\nu'(\theta su)
\equiv-\frac38a_\eta(aHu)H_j(asu)^2.
\end{align*}

C) From formulae (31.) follows the reduction of the decomposed forms
\[
 [L_\vartheta(\theta Lu)]_{su^3}\quad\hbox{and}\quad [\theta_l(\theta Lu)]_{s^2u^3},
\]
that is, of the products $[\theta_\nu^2H]_{su^3}$ and $[a_\nu^2(bHu)^2]_{su^3}$:
\begin{equation}
\tag{32}
\begin{aligned}
0&\equiv\theta_\nu^2(\theta su)^2(Hsu),
0&\equiv [L_\vartheta(\theta Lu)]_{su^4}-7[\theta_l(\theta Lu)]_{su^4},\\
0&\equiv a_\nu^2(asu)^2(bHu)^2(bsu)\\
&\quad+6\theta_l(\theta Lu)^2(\theta su)(Lsu),\\
0&\equiv\theta_\nu^3(\theta su)(Hsu)^2\\
&\quad\hbox{from }0\equiv\theta_l^2(\theta Lu)(\theta su)(Lsu)^2\\
&\quad(\hbox{reducer }a_\eta^2(Hsu)^2).
\end{aligned}
\end{equation}

\noindent 13th order.
\[
\text{A)}\ ((KLu)^3,s,u)^2,\ ((KLu)^3,s,u)^3,\qquad
\text{B)}\ (L_jsu)^2,
\qquad
\text{C)}\ ((aH^2u)^3,s,u)^2.
\]
\noindent 14th order.
\[
\text{A)}\ ((aLu)^2L_j,s,u)^2,
\qquad
\text{B)}\ ((\theta H^2u)^3,s,u)^2.
\]
\noindent 15th order.
\[
 ((KH^2u)^4,s,u)^2.
\]
\noindent\emph{Reductions:} Decomposition of $(K_l(KLu)^3,s,u)^2$, $(H_\vartheta(\theta H^2u)^3,s,u)$, $((K,H\cdot L,u)^5,s,u)$ according to \S\ 3. 2), c); that is, with simultaneous consideration of the decomposed forms $(K_l(KLu)^2,s,u)^3$, $(H_\vartheta(\theta H'u)^2,s,u)^2$, $((K,H\cdot L,u)^4,s,u)^2$.

\noindent\emph{Consequence:} $s^2$ does not enter into folding with individual forms of System III.

\begin{center}\textbf{\S\ 26.\quad System $s_{IV}$.}\end{center}
\noindent 1) \emph{Folding of $s$ with System I--III.}

\noindent\emph{Reductions:} According to the reductions of the last paragraphs $(a_\nu^2s_\nu(aHu)^2(asu)s_\eta$, etc.), the forms $(0,\lambda)$, $(1,\lambda)$, $(2,\lambda)$ do not permit foldings I and III simultaneously; hence, by \S\ 18, the forms $f,\D,N$ do not enter in products in folding to form System $s_{IV}$.

The form $j$ does not enter into folding, since, by the same reductions, the foldings contragredient to $j$,
\[
(K_\nu(k\nu x)^3,s,u),\qquad ((\vartheta^2\eta x)^4,s,u)\quad\hbox{etc.},
\]
correspond to the foldings cogredient to $j$:
\[
K_\nu(k\nu x)^2s_k,
\qquad (\vartheta^2\eta x)^3s_\vartheta\quad\hbox{etc.}
\]

\noindent 2) \emph{Folding of $s$ with System II--III}, that is, of $s$ with $H\cdot(1,\lambda)$, $H\cdot(2,\lambda)$, $H\cdot(3,\lambda)$.

\noindent\emph{Irreducible forms:}
\[
\begin{array}{rl}
\text{A)}&(a_\nu H,s,u)^4,
((aHu)^2H,s,u)^3,
((aHu)^2H,s,u)^4,\\
&((aLu)^2H,s,u)^4,
((aLu)^3H,s,u)^3,
((aH^2u)^3H,s,u)^4,\\[0.3em]
\text{B)}&(a_\nu^2H,s,u)^3,
(a_\eta(aHu)^2H,s,u)^3,
(a_l(aLu)^2H,s,u)^4,\\[0.3em]
\text{C)}&(\theta_\nu H,s,u)^4,
((\theta Hu)^2H,s,u)^3,
((\theta Hu)^2H,s,u)^4,\\
&((\theta Lu)^2H,s,u)^4,
((\theta Lu)^3H,s,u)^3,
((\theta H^2u)^3H,s,u)^4,\\[0.3em]
\text{D)}&(\theta_\eta(\theta Hu)^2H,s,u)^3,
(\theta_l(\theta Lu)^2H,s,u)^4,\\[0.3em]
\text{E)}&((KLu)^3H,s,u)^4,
((KH^2u)^4H,s,u)^4,\\[0.3em]
\text{F)}&((j\nu x)^2H,s,u)^3,
((j\nu x)^2H,s,u)^4,
((j\eta x)H,s,u)^4,\\
&(H_jH,s,u)^3,
(L_jH,s,u)^4.
\end{array}
\]
\noindent\emph{Reductions:} $\nu$ does not enter into folding, analogously to $j$.

B) and D). $(a_\nu^2H,s,u)^4$ and $(\theta_\nu^2H,s,u)^4$ decompose by formula (27.)c) and (29.)c), taking into account $(gHu)^2=-\frac13s\sigma$:
\[
(a_\nu^2H,s,u)^4=\frac13(asu)^4\sigma,
\qquad
(\theta_\nu^2H,s,u)^4=-\frac16(\theta su)^4\sigma;
\]
hence decomposition of $(a_\eta(aHu)H,s,u)^4$, $(\theta_\eta(\theta Hu)H,s,u)^4$.

$(\theta_\nu^2H,s,u)^3$, $(\theta_\nu^3H,s,u)^3$, and hence $(\theta_\eta^2(\theta Hu)H,s,u)^4$, decompose according to \S\ 25, forms of order 12, C).

\noindent 3) \emph{Folding of $s$ with System III--III}, that is, of $s$ with forms of System III:
\[
(3,\lambda)(3,\lambda),\quad (3,\lambda)(2,\lambda),\quad (3,\lambda)(1,\lambda),\quad
(2,\lambda)(2,\lambda),\quad (2,\lambda)(1,\lambda),\quad (1,\lambda)(1,\lambda).
\]
\noindent\emph{Irreducible forms:}
\[
\begin{array}{rl}
\text{A)}&(a_\nu^2a_\nu^2,s,u)^3,
(a_\nu^2a_\nu^2,s,u)^4,
(a^2a_\eta(aHu)^2,s,u)^3,\\[0.3em]
\text{B)}&(a_\nu H_j,s,u)^4,
((aHu)^2H_j,s,u)^3,
((aLu)^2H_j,s,u)^4,
((aLu)^3H_j,s,u)^2,
((aH^2u)^3H_j,s,u)^3.
\end{array}
\]
\noindent\emph{Reductions:}

Products II--III, for the same order in the symbols $\nu$ and $f$, are to be considered higher than products III--III.

The reductions arise partly from relations between products (syzygies), partly by replacing the products by the corresponding decomposed forms and reducing the folding of these forms with $s$. Products containing contravariants are always omitted, since they pass into products.

A) $f$ quadratic, symbols $\nu$ of arbitrary order. By \S\ 11, formulae (12.), and by analogous computation, one has:
\begin{align*}
1)\quad 0&=a_\nu b_{\nu_1}+\frac12f(aHu)^2-\frac14\theta H+\frac12u_xH_\vartheta-\frac14u_x^2H_\vartheta^2,\\
2)\quad 0&=a_\nu b_{\nu_1}^2+f a_\eta(aHu)^2-\theta_\eta-\frac12H_\vartheta,\\
3)\quad 0&=a_\nu^2b_{\nu_1}^2-H_\vartheta^2+2\theta_\eta^2.
\end{align*}
It follows that
\begin{align*}
4)\quad 0&=a_\nu^2b_{\nu_1}^2(j\nu_1x),\\
5)\quad L_\vartheta(\theta Lu)&=-a_\nu b_\eta(bHu)^2+\frac34a_\nu^2(bHu)^2+\frac34\theta_\nu^2H,\\
6)\quad L_\vartheta(\theta Lu)&=-6a_\nu b_\eta(bHu)^2+2a_\nu^2(bHu)^2-\frac12\theta_\nu^2H,\\
7)\quad 0&=a_\nu(bHu)^2+f(aLu)^3+\theta_\nu H,\\
8)\quad 0&=a_\nu(bLu)^3+\frac34(\theta Hu)^2H,\\
9)\quad 0&=(aHu)^2(bHu)^2+2(\theta Hu)^2H,\\
10)\quad 0&=a_\eta(aHu)^2b_\eta(bHu)^2,\\
11)\quad 0&\equiv [a_\nu^2b_\eta(bHu)]_{\widehat{su^4}}.
\end{align*}
The reductions of
\[
(H_\vartheta su)^4,
\qquad (L_\vartheta(\theta Lu),s,u)^3,
\qquad (L_\vartheta(\theta Lu),s,u)^4
\]
(formulae (31.) and (32.)) together with formulae 1) to 11) give the reduction of all foldings of $s$ with products that are quadratic in the symbols $f$ and arbitrary in $\nu$.

B) Products of two of the symbols $f,\theta,j$; $\nu$ in arbitrary order. By formulae (17.), (18.), (20.), and by direct computation, the relations are:
\begin{align*}
1)\quad 0&=a_\nu\theta_{\nu_1}+\frac14\theta(aHu)^2+\frac14f(\theta Hu)^2,\\
2)\quad 0&=\theta_\nu\theta_{\nu_1}+\frac12\theta(\theta Hu)^2.
\end{align*}
Therefore:
\begin{align*}
3)\quad 0&=3a_\nu(\theta Hu)^2+\theta(aLu)^3+2f(\theta Lu)^3,\\
4)\quad 0&=3\theta_\nu(aHu)^2+f(\theta Lu)^3+2\theta(aLu)^3,\\
5)\quad 0&=a_\nu(\theta Lu)^3,
&0&=\theta_\nu(aLu)^3,
&0&=(aHu)^2(\theta Hu)^2,\\
6)\quad 0&=a_\nu\theta_\nu^2+\theta_\nu a_\nu^2+f\theta_\eta(\theta Hu)^2+\theta a_\eta(aHu)^2,\\
7)\quad 0&=\theta_\nu\theta_\nu^2+\theta\theta_\eta(\theta Hu)^2+\frac16fH_j,\\
8)\quad 0&=a_\nu(j\nu x)^2+fH_j,\\
9)\quad 0&=(aHu)^2(j\nu x)^2-2a_\nu H_j,\\
10)\quad 0&=\theta_\nu(j\nu x)^2,
&0&=\theta H_j,\\
11)\quad 0&=a_\nu\theta_\nu^3-\frac34(j\eta x)^2,\\
12)\quad 0&=a_\nu^2\theta_\nu^2-\frac13(j\eta x)^2-\frac13(\D Hu)^2,\\
13)\quad 0&=a_\nu^2\theta_\nu^3-\frac12a_\sigma,\\
14)\quad 0&=\theta_\nu\theta_\nu^3,
&0&=a_\eta H_j.
\end{align*}
The formulae have been carried far enough for the products to become polarizable: that is, by folding only one new product arises, because a) the folding of the product into itself becomes reducible, and b) the remaining products arising by folding become reducible or lead to contravariants (cf. \S\ 3. 2), a) and b)).

Formulae 1) to 5) give the reduction of all foldings of $s$ with products arising from $a_\nu\theta_\nu$ by folding. Formulae 6) to 10) give the reduction of those products arising from $a_\nu\theta_\nu^2$ by folding. By \S\ 24 A), taking formulae 6) to 10) into account, the foldings of $s$ with products arising from $a_\nu^2\theta_\nu$ decompose.

One has:
\begin{align*}
0&\equiv a_\nu^2(asu)^2\theta_{\nu_1}(\theta su)^2,
&0&\equiv a_\nu^2(asu)\theta_{\nu_1}(\theta su)^3-K_\eta(KHu)^2(Ksu)^3,\\
0&\equiv a_\nu^2(asu)^2\theta_{\nu_1}(\theta su),
&0&\equiv a_\nu^2(asu)\theta_{\nu_1}(\theta su)^2,\\
0&\equiv a_\nu^2(asu)\theta_{\nu_1}(\theta su).
\end{align*}
The reducers
\[
 ((j\eta x)^2,s,u)^3\ [\hbox{from }(\widehat{\vartheta\eta su})^2(Hsu),\ \S\ 23.\ C3],
\]
\[
(H_j(j\eta x),s,u)^2\ [\S\ 24,\ B],
\quad ((\D Hu)^2,s,u)^2\ [\S\ 24,\ C],
\quad (a_\sigma su)^2\ [\S\ 24,\ E]
\]
together with formulae 11) to 14) give the reduction of all products arising from $a_\nu\theta_{\nu_1}^3$, $a_\nu^2\theta_{\nu_1}^2$, $a_\nu^2\theta_{\nu_1}^3$.

Our preceding computations can be summarized in the following \emph{conclusion}:

The relatively complete system modulo $(\R,t)$ consists of the following 331 forms and subsystems, which are also reproduced below in tabular order.

\clearpage
\begin{landscape}
\thispagestyle{plain}
\begin{center}\bfseries Table I (cf. p. 90)\end{center}
\vspace{-0.6em}
\begingroup
\setlength{\tabcolsep}{2pt}
\renewcommand{\arraystretch}{1.35}
\tiny
\begin{center}
\begin{adjustbox}{max width=0.98\linewidth,max totalheight=0.78\textheight,center}
\begin{tabular}{|c|c|c|c|c|c|c|c|c|c|c|c|c|c|c|c|c|c|c|c|c|c|c|}
\hline
\diagbox{$u$}{$x$} & 0 & 1 & 2 & 3 & 4 & 5 & 6 & 7 & 8 & 9 & 10 & 11 & 12 & 13 & 14 & 15 & 16 & 17 & 18 & 19 & 20 & 21 \\ \hline
0 & \mcell{i} &  &  &  & \mcell{f} &  & \mcell{\A} &  & \mcell{K_{\nu}(k\nu x)^3} &  & \mcell{K_{\eta}(k\eta x)^3} &  &  &  & \mcell{(\vartheta^3\eta x)^4} & \mcell{H_k(k\vartheta,\eta x)^4} &  & \mcell{\theta_l(\vartheta k,lx)^5} &  &  &  & \mcell{(\vartheta^2lx)^6} \\ \hline
1 &  & \mcell{u_x} &  &  &  & \mcell{\theta_\nu(\vartheta\nu x)^2} &  & \mcell{\theta_\eta(\vartheta\eta x)^2} & \mcell{N} & \mcell{(k\nu x)^3} & \mcell{\theta_\nu(\vartheta^2\nu x)^3} & \mcell{(k\eta x)^3} & \mcell{H_{\vartheta}(\vartheta^2\eta x)^3\\ \theta_\eta(\vartheta^2\eta x)^3} &  & \mcell{\theta_l(\vartheta^2lx)^4} &  & \mcell{(\vartheta k,\eta,x)^4} &  & \mcell{(\vartheta k,l,x)^5} &  &  &  \\ \hline
2 &  &  & \mcell{a_\nu^2} &  & \mcell{\theta\\ a_\eta^2} &  & \mcell{(\vartheta\nu x)^2} & \mcell{K_\nu(k\nu x)^2} & \mcell{(\vartheta\eta x)^2} & \mcell{H_k(k\eta x)^2\\K_\eta(k\eta x)^2} &  & \mcell{(\vartheta^2\nu x)^3\\ K_l(klx)^3} &  & \mcell{(\vartheta^2\eta x)^3} &  & \mcell{(\vartheta^2lx)^4} &  &  &  &  &  &  \\ \hline
3 &  & \mcell{\theta_\nu^3} &  & \mcell{a_\nu} & \mcell{\theta_\nu(\vartheta\nu x)} & \mcell{\A_\nu\\a_\eta} & \mcell{K\\H_{\vartheta}(\vartheta\eta x)\\\theta_\eta(\vartheta\eta x)} & \mcell{\A_\eta} & \mcell{(k\nu x)^2\\\theta_l(\vartheta lx)^2} &  & \mcell{(k\eta x)^2} &  & \mcell{(klx)^3} &  &  &  &  &  &  &  &  &  \\ \hline
4 & \mcell{\nu} &  & \mcell{H\\\theta_\nu^2} & \mcell{(j\nu x)^3\\a_\eta(aHu)} & \mcell{\theta_\eta^2} & \mcell{(\vartheta\nu x)\\a_l^2} &  & \mcell{N_\nu\\(\vartheta\eta x)} &  & \mcell{H_\eta\\N_\eta\\(\vartheta lx)^2} &  &  &  &  &  &  &  &  &  &  &  &  \\ \hline
5 &  & \mcell{a_\eta(aHu)^2} & \mcell{\theta_\eta^2(\theta Hu)} & \mcell{\theta_\nu} & \mcell{\frac{K_\nu^2}{(aHu)}} & \mcell{\theta_\eta} & \mcell{K_\eta^2\\(\A Hu)\\a_l} &  & \mcell{\A_l} &  &  &  &  &  &  &  &  &  &  &  &  &  \\ \hline
6 & \mcell{j\\\sigma} &  & \mcell{(j\nu x)^2\\(aHu)^2} & \mcell{L\\a_\nu^2(j\nu x)\\H_{\vartheta}(\theta Hu)\\\theta_\eta(\theta Hu)} & \mcell{\frac{K_\nu}{\theta_l^2}} &  &  & \mcell{K_\eta} &  &  &  &  &  &  &  &  &  &  &  &  &  &  \\ \hline
7 &  & \mcell{\theta_\eta(\theta Hu)^2} & \mcell{H_j(j\eta x)\\a_l(aLu)^2} &  & \mcell{a_\nu(j\nu x)\\(\theta Hu)} &  & \mcell{\A_\nu(j\nu x)\\\theta_l} & \mcell{K_l^2} &  &  &  &  &  &  &  &  &  &  &  &  &  &  \\ \hline
8 & \mcell{H_j^2\\a_l(aLu)^3} & \mcell{(\nu\sigma x)\\(j\nu x)} & \mcell{(\theta Hu)^2} & \mcell{K_\eta(KHu)^2\\(j\eta x)\\(aLu)^2} &  &  &  &  &  &  &  &  &  &  &  &  &  &  &  &  &  &  \\ \hline
9 &  & \mcell{H_j\\(aLu)^3} & \mcell{a_\eta(aHu)H_j\\L_{\vartheta}(\theta Lu)^2\\\theta_l(\theta Lu)^2} &  & \mcell{(KHu)} &  &  &  &  &  &  &  &  &  &  &  &  &  &  &  &  &  \\ \hline
10 & \mcell{\theta_l(\theta Lu)^3} & \mcell{L_j^2\\(j\sigma x)\\a_\eta(aH^2u)^3} &  & \mcell{H_j(aHu)\\(\theta Lu)^2} &  &  &  &  &  &  &  &  &  &  &  &  &  &  &  &  &  &  \\ \hline
11 &  & \mcell{(\theta Lu)^3} & \mcell{K_l(KLu)^3\\L_j\\(aH^3u)^3} &  &  &  &  &  &  &  &  &  &  &  &  &  &  &  &  &  &  &  \\ \hline
12 & \mcell{(aH^2u)^4} & \mcell{a_\eta(aLu)^2L_j\\H_{\vartheta}(\theta H^2u)^3\\\theta_\eta(\theta H^2u)^3} &  & \mcell{(KLu)^3} &  &  &  &  &  &  &  &  &  &  &  &  &  &  &  &  &  &  \\ \hline
13 & \mcell{(\nu\sigma j)} &  & \mcell{(aLu)^2L_j\\(\theta H^3u)^3} &  &  &  &  &  &  &  &  &  &  &  &  &  &  &  &  &  &  &  \\ \hline
14 & \mcell{(\theta H^2u)^4} & \mcell{K_\eta(KH^2u)^4\\(a_\eta H\!\cdot L,u)^4} &  &  &  &  &  &  &  &  &  &  &  &  &  &  &  &  &  &  &  &  \\ \hline
15 & \mcell{L_9(\theta,H\!\cdot L,u)^4} &  & \mcell{(KH^2u)^4} &  &  &  &  &  &  &  &  &  &  &  &  &  &  &  &  &  &  &  \\ \hline
16 & \mcell{(\theta,H\!\cdot L,u)^5} &  &  &  &  &  &  &  &  &  &  &  &  &  &  &  &  &  &  &  &  &  \\ \hline
17 & \mcell{K_\eta(K,H\!\cdot L,u)^5} &  &  &  &  &  &  &  &  &  &  &  &  &  &  &  &  &  &  &  &  &  \\ \hline
18 &  & \mcell{(K,H\!\cdot L,u)^5} &  &  &  &  &  &  &  &  &  &  &  &  &  &  &  &  &  &  &  &  \\ \hline
19 &  &  &  &  &  &  &  &  &  &  &  &  &  &  &  &  &  &  &  &  &  &  \\ \hline
20 &  &  &  &  &  &  &  &  &  &  &  &  &  &  &  &  &  &  &  &  &  &  \\ \hline
21 & \mcell{(KH^3u)^6} &  &  &  &  &  &  &  &  &  &  &  &  &  &  &  &  &  &  &  &  &  \\ \hline
\end{tabular}
\end{adjustbox}
\end{center}
\vspace{-0.2em}
\begin{center}\footnotesize (The numbers in the uppermost horizontal row denote the degree in the $x$; the numbers in the first vertical row denote the degree in the $u$.)\end{center}
\endgroup
\end{landscape}
\clearpage
\begin{landscape}
\thispagestyle{plain}
\begin{center}\bfseries Table II (cf. p. 90)\end{center}
\vspace{-0.6em}
\begingroup
\setlength{\tabcolsep}{2pt}
\renewcommand{\arraystretch}{1.35}
\tiny
\begin{center}
\begin{adjustbox}{max width=0.98\linewidth,max totalheight=0.78\textheight,center}
\begin{tabular}{|c|c|c|c|c|c|c|c|c|c|c|c|c|c|c|c|c|c|}
\hline
\diagbox{$u$}{$x$} & 0 & 1 & 2 & 3 & 4 & 5 & 6 & 7 & 8 & 9 & 10 & 11 & 12 & 13 & 14 & 15 & 16 \\ \hline
0 & \mcell{J} &  &  &  & \mcell{s} &  & \mcell{s_j^3} &  &  &  &  & \mcell{s_\nu} & \mcell{(k\nu x)^3s_\nu} & \mcell{(\vartheta^3\nu x)^2s_{\vartheta}\\(\vartheta^2\nu x)^2s_{\vartheta}} &  & \mcell{\theta_\eta(\vartheta^2\eta x)^2s_\eta} &  \\ \hline
1 &  &  &  &  &  &  & \mcell{(asu)} & \mcell{s_\eta} & \mcell{s_l^2\\(\A su)} & \mcell{s_\eta(Nsu)\\(\vartheta\nu x)^2s,u} & \mcell{((k\nu x)^3,s,u)_s\\(K_\nu(k\nu x)^3,s,u)} &  & \mcell{K_\eta(k\eta x)^3,s,u} &  & \mcell{(\vartheta^2\nu x)^2s_\nu} &  & \mcell{(\vartheta^2\eta x)^4,s,u} \\ \hline
2 &  &  &  &  & \mcell{(asu)^2} & \mcell{s_\theta(\theta su)} & \mcell{(\A su)^2/a_\nu s_\nu} & \mcell{(\vartheta\nu x)^2,s,u\\(\theta_l(\vartheta\nu x)^2,s,u)} &  & \mcell{\theta_\eta(\vartheta^2\eta x)^2,s,u} &  & \mcell{(k\eta x)^2s_\nu\\((k\nu x)^2,s,u)} &  & \mcell{(k\eta x)^3,s,u} &  &  &  \\ \hline
3 &  &  & \mcell{(asu)^3} & \mcell{s_\vartheta(\theta su)^2\\s_\nu} & \mcell{a_\nu s_\nu(asu)\\a_l^2(asu)} & \mcell{s_\eta} & \mcell{(\theta su)\\a_l^2su} &  & \mcell{(Nsu)^2\\(\vartheta\eta x)s_\nu\\((\vartheta\nu x)^2,s,u)} & \mcell{((k\nu x)^2,s,u)^2} & \mcell{(\vartheta^2\eta x)^2,s,u} &  &  &  &  &  &  \\ \hline
4 & \mcell{(asu)^4} & \mcell{s_\vartheta(\theta su)^3} & \mcell{a_\nu^2(asu)^2} & \mcell{(\theta_\nu^2su)} & \mcell{(\theta su)^2} & \mcell{s_k(Ksu)^3\\a_\nu(asu)\\s_\nu(asu)} & \mcell{((\vartheta\nu x)^2,s,u)^2} & \mcell{s_\nu(\A su)\\(\A su)\\(aHu)s_\eta\\(a_\nu su)} &  & \mcell{(\A_\eta su)} &  &  &  &  &  &  &  \\ \hline
5 &  &  & \mcell{(\theta su)^3} & \mcell{s_k(Ksu)^3,s_j\\s_\vartheta(\theta s^\prime u)^4\\s_\nu(asu)^2\\a_\nu(asu)^2} & \mcell{(Hsu)\\((\vartheta\nu x)^2,s,u)\\(\theta_\nu(\vartheta\nu x)^2,s,u)\\(\theta_\nu^2su)} & \mcell{((j\nu x)^2,s,u)\\(aHu)^2s_\nu\\(a_\eta su)^2} & \mcell{(Ksu)^2\\s_\eta\\(\theta_\eta^2su)} &  & \mcell{((k\nu x)^2,s,u)^2\\K_\nu s_\nu} &  &  &  &  &  &  &  &  \\ \hline
6 & \mcell{(\theta su)^4} & \mcell{s_\nu(asu)^3\\a_\nu(asu)^3} & \mcell{(Hsu)^2\\(\theta_\nu(\vartheta\nu x),s,u)^3\\(\theta_\nu^2su)^2} & \mcell{a_\eta s_\nu^2\\(a_\eta(aHu),s,u)^2\\(a_\eta(aH^2u),s,u)} & \mcell{(Ksu)^3} & \mcell{s_\eta(asu)\\((\vartheta\nu x),s,u)^3\\(\theta_\eta su)} & \mcell{((k\nu x)^2,s,u)^3} & \mcell{s_\nu(\A su)\\a_\nu(j\nu x)s_\nu\\(\theta H u)s_\eta\\(\theta_\eta su)} &  &  &  &  &  &  &  &  &  \\ \hline
7 & \mcell{\theta_\nu(\vartheta\nu x),s,u)^4} & \mcell{(a_\eta(aHu),s,u)^3} & \mcell{(Ksu)^4\\(\theta_\eta^2(\theta Hu),s,u)^2\\(a_\eta^2-a_l^2,s,u)^3} & \mcell{s_j(asu)^3\\s_k(K^2u)^3\\((j\nu x),s,u)^2\\(\theta_\eta su)^2} & \mcell{((j\nu x),s,u)\\((j\eta x),s,u)\\(aHu),s,u)\\(aH^2u),s,u)} & \mcell{((\vartheta\eta x),s,u)^3} & \mcell{(aLu)^3s_l/(asu)^3} &  &  &  &  &  &  &  &  &  &  \\ \hline
8 & \mcell{(a^2-a_l^2,s,u)^4} & \mcell{s_j(asu)^3\\s_k(K^2u)^3\\((j\nu x),s,u)^4\\(\theta_\eta su)^3} & \mcell{((j\nu x)^2,s,u)^2\\(aHu),s,u)^3\\((aHu)^2,s,u)^3} & \mcell{(Lsu)^2\\(a_l^2(j\nu x),s,u)^3\\H_\vartheta(\theta Hu),s,u)^3\\\theta_l(\theta Hu),s,u)^3} & \mcell{(asu)^3} & \mcell{(K,su)^3} &  & \mcell{(K_\vartheta su)^2} &  &  &  &  &  &  &  &  &  \\ \hline
9 & \mcell{K_j^2su^4\\((aHu),s,u)^4} & \mcell{(Lsu)^4\\(a_l^2(j\nu x),s,u)^4\\(\theta_\nu(\theta Hu),s,u)^2} & \mcell{(K_\nu^3u)^6\\(a_\eta(aLu)^2,s,u)^3\\(a_l^2H,s,u)^3} & \mcell{(K,su)^3} & \mcell{(a_\nu(j\nu x),s,u)^3\\(\theta Hu),s,u\\(\theta H u)^2,s,u} &  & \mcell{(\theta_lsu)^3} &  &  &  &  &  &  &  &  &  &  \\ \hline
10 &  &  & \mcell{(K,su)^4\\(a_l^2a_\eta(aHu)^2,s,u)^3} & \mcell{((j\eta x),s,u)^3\\(H_j,su)\\((aLu)^2,s,u)^2\\((aLu)^3,s,u)} &  &  &  &  &  &  &  &  &  &  &  &  &  \\ \hline
11 & \mcell{(a_\nu(j\nu x),s,u)^4\\((\theta Hu),s,u)^4} & \mcell{(K_l(KH u)^2,s,u)^3\\((j\eta x),s,u)^3\\((aLu)^2,s,u)^4\\(a_\eta H,s,u)^3} & \mcell{(H^2su)^3\\\theta_l(\theta Lu)^2,s,u)^3} &  & \mcell{((KHu)^2,s,u)^3} &  &  &  &  &  &  &  &  &  &  &  &  \\ \hline
12 &  & \mcell{(a_\eta(aHu)^2H,s,u)^3} & \mcell{((KH u)^3,s,u)^3} & \mcell{((aHu)H_j,s,u)^2\\((\theta Lu)^2,s,u)^3} &  &  &  &  &  &  &  &  &  &  &  &  &  \\ \hline
13 & \mcell{((KH u)^3,s,u)^4} & \mcell{((aHu)H_j,s,u)^3\\((\theta Lu)^2,s,u)^4\\(\nu\cdot H,s,u)^3} & \mcell{(L_jsu)^2\\((aH^3u)^3,s,u)^3\\((j\eta x)H,s,u)^3\\((aHu)^2H,s,u)^3} &  &  &  &  &  &  &  &  &  &  &  &  &  &  \\ \hline
14 & \mcell{((j\nu x)^2H,s,u)^4\\((aHu)^2H,s,u)^4} & \mcell{(\theta_\eta(\theta Hu)^2H,s,u)^3} &  & \mcell{((KLu)^3,s,u)^2} &  &  &  &  &  &  &  &  &  &  &  &  &  \\ \hline
15 & \mcell{(a_l(aLu)^3H,s,u)^4} & \mcell{((KLu)^3,s,u)^3} & \mcell{((aLu)^2L_j,s,u)^2\\((\theta Hu)^3H,s,u)^3\\((\theta Hu)^2H,s,u)^2} &  &  &  &  &  &  &  &  &  &  &  &  &  &  \\ \hline
16 & \mcell{((\theta Hu)^2H,s,u)^4\\(\sigma\cdot H_j,s,u)^4} & \mcell{((j\eta x)H,s,u)^4\\(H_jH,s,u)^3\\((aLu)^2H,s,u)^4\\((aLu)^3H,s,u)^3} &  &  &  &  &  &  &  &  &  &  &  &  &  &  &  \\ \hline
17 & \mcell{(\theta(\theta Lu)^2H,s,u)^4} &  & \mcell{((KH^3u)^3,s,u)^3} &  &  &  &  &  &  &  &  &  &  &  &  &  &  \\ \hline
18 &  & \mcell{((\theta Lu)^2H,s,u)^4\\((\theta Lu)^3H,s,u)^3\\(aHu)^2H,s,u)^4} &  &  &  &  &  &  &  &  &  &  &  &  &  &  &  \\ \hline
19 & \mcell{((aH^3u)^3H,s,u)^4\\(L_jH,s,u)^5} &  &  &  &  &  &  &  &  &  &  &  &  &  &  &  &  \\ \hline
20 &  & \mcell{((KLu)^3H,s,u)^5} & \mcell{((aLu)^3H,s,u)^5} &  &  &  &  &  &  &  &  &  &  &  &  &  &  \\ \hline
21 & \mcell{((\theta H^3u)^3H,s,u)^4\\((aLu)^2H_j,s,u)^4} &  &  &  &  &  &  &  &  &  &  &  &  &  &  &  &  \\ \hline
22 &  &  &  &  &  &  &  &  &  &  &  &  &  &  &  &  &  \\ \hline
23 & \mcell{((KH^3u)^4H,s,u)^4} & \mcell{((aH^3u)^4H,s,u)^3} &  &  &  &  &  &  &  &  &  &  &  &  &  &  &  \\ \hline
\end{tabular}
\end{adjustbox}
\end{center}
\vspace{-0.2em}
\begin{center}\footnotesize (The numbers in the uppermost horizontal row denote the degree in the $x$; the numbers in the first vertical row denote the degree in the $u$.)\end{center}
\endgroup
\end{landscape}



\begin{center}
{\Large\bfseries 3. On the Invariant Theory of Forms in $n$ Variables\footnote{Lecture delivered at the Salzburg meeting of natural scientists, 1909.}}\par
\vspace{1.2em}
Jahresbericht der DMV 19 (1910), pp. 101--104
\end{center}

\vspace{1em}
In the projective invariant theory of forms in $n$ variables the main problem has been solved: the finiteness of the system of forms has been proved. A series of further questions, however, have not been treated, namely those concerned with methods for investigating the connection among forms. I would like briefly to communicate the results that I have obtained with respect to such questions, but for the sake of clarity I shall first sketch the questions in the ternary domain, where they are known.

I have first in mind the two fundamental theorems of the symbolic method: the theorem that all invariant formations can be represented symbolically, and the second theorem that all relations existing among invariants are obtained by successive application of a small number of symbolic identities; thus the symbolic method gives all relations without leaving the domain of invariants.\footnote{Study: \emph{Methoden zur Theorie der ternären Formen}. II. \S\ 6. (Leipzig, Teubner, 1889.)} Built upon this are the processes for generating the forms, the symbolic folding process and the nonsymbolic differentiation processes, the theory of series expansions, and the like. In this connection I should also mention a theorem proved in my dissertation,\footnote{\emph{Über die Bildung des Formensystems der ternären biquadratischen Form}. \S\ 1. (Crelle's Journal Bd. 134. 1908.)} namely that all foldings can be generated by successive application of the two possible ``cogredient'' foldings; by this, for a symbolic product $s_x t_x u_\sigma u_\tau$, I mean the two foldings $(stu)$ and $(\sigma\tau x)$. On this theorem one can build the theory of reducers and of the reduction of systems of forms, in analogy with the binary domain, as I showed in my dissertation.

When we now pass to the $n$-ary domain, even the first theorem of the symbolic method holds only in a conditional measure. It is well known that already for quaternary forms there occur, besides the point and plane coordinates $x$ and $u$, line coordinates $p_{ik}$, the subdeterminants from two rows of point coordinates. Analogously, in the $n$-ary domain we have rows of variables $p_\rho$, whose individual elements are the subdeterminants from $\rho$ rows $x$, and rows of symbols $S^\rho$, whose elements behave like subdeterminants from $\rho$ rows $s$ cogredient to the $u$. The symbolic representation by determinants is known to hold only when the symbol rows $S^\rho$ are resolved into the rows $s$ and these are distributed among different determinants; only such aggregates of determinants as depend on the $S^\rho$ then have real meaning. But which aggregates of determinants accomplish this cannot be surveyed, and thereby all the other theorems mentioned in the ternary domain are lost.

I would now like to state a simple principle by which one obtains a symbolic representation explicit in the rows $S^\rho$, and thereby also recovers the remaining theorems. This is an application of the familiar theorem on corresponding matrices: ``To each row of variables $p_\rho$ there corresponds one-to-one another row $q_{n-\rho}$ whose individual elements are subdeterminants from $n-\rho$ rows $u$ connected with the $\rho$ rows $x$ by equations $(u\mid x)=0$.'' Correspondingly, to each symbol row $S^\rho$ there is assigned another $S_{n-\rho}$, behaving as if it were composed from $n-\rho$ rows cogredient to the $x$.

The theorem also permits a second formulation, suited for applications: ``To every matrix product $(v^{(1)}\cdots v^{(\rho)}\mid y^{(1)}\cdots y^{(\rho)})$,\footnote{That is, the sum over the products of corresponding determinants, formed from the matrix of the $\rho$ rows $v$ and that of the $\rho$ rows $y$.} where the rows of the second matrix are contragredient to those of the first, there is assigned a determinant; and conversely each determinant can be transformed into a matrix product with a prescribed number $\rho$ of rows.'' Thus determinant and matrix product appear as fully equivalent, and the first theorem of the symbolic method assumes the form: ``All invariant formations can be represented symbolically by matrix products.'' From two symbol rows $S^\sigma$, $T^\tau$ one can now, with the aid of rows of variables, very easily form a matrix product explicitly containing these symbol rows and hence representing an invariant formation of the form $(S^\sigma\mid p_\sigma)(T^\tau\mid p_\tau)$, namely the following:
\begin{equation}
 \pair{q_{n-\sigma-\lambda}S^\sigma}{T_{n-\tau}p_{\tau-\lambda}},
 \qquad (\lambda=1,2,\ldots,\tau\ \text{or}\ n-\sigma).
\tag{1}
\end{equation}

More important is the converse: that all invariant formations can be represented explicitly by matrix products, so that the above process is the extension of the binary and ternary folding process. To prove this converse it is necessary to transfer the second theorem of the symbolic method to the $n$-ary domain, namely to set up the totality of the independent indecomposable identities in which all rows $S^\rho$, $p_\rho$ are regarded as indecomposable. These identities arise from the known ones containing only rows $x$ and $u$,\footnote{E. Pascal in \emph{Memorie d. R. Acc. d. Lincei}. (4) 5 (1888), p. 375.} by replacing the product theorem for determinants by a system of product theorems for matrices, and by contracting different matrix products into a single one by means of precisely these product theorems. One obtains in this way two dualistic formulas, of which only one is given:
\begin{equation}
\begin{aligned}
 \pair{S_1^{\rho_1}S_2^{\rho_2}\cdots S_k^{\rho_k}}{x p_{\rho-1}}
 &=\sum_{i=1}^{k}\eps\,\pair{S_1^{\rho_1}\cdots S_{i-1}^{\rho_{i-1}}S_{i+1}^{\rho_{i+1}}\cdots S_k^{\rho_k}q_{n-\rho_i+1}}{S_{n-\rho_i}x},\\
 \rho&=\sum_{i=1}^{k}\rho_i<n,\qquad \eps=\pm1.
\end{aligned}
\tag{2}
\end{equation}

The only specialization contained in these formulas, namely that a row $x$ enters, cannot be avoided; for completely general symbol and variable rows we arrive at a ``decomposition identity,'' an identity in which a row $p_\rho$ is decomposed into the individual rows $x$. The simplest special case of the decomposition identity was already used by Clebsch in his ``fundamental problem of invariant theory'' to derive the elementary covariants in the theory of series expansions. With the aid of these identities, which mediate the transition from the non-explicit determinant expressions to the matrix product, one can in fact show that the desired explicit symbolic representation is obtained with the matrix representation.

For the folding process, however, a theorem entirely analogous to that in the ternary domain holds, and the reduction theorems can likewise be built upon it analogously. If in (1) one denotes the number $\lambda$ as the defect of the folding, and calls foldings for which $\sigma=\tau$ cogredient, then all foldings can be generated by successive application of the cogredient foldings of defect 1. The proof of this theorem rests essentially on the fact that the most general forms can be replaced by ``normal forms,'' that is, by forms that vanish identically under the application of all invariant differential operations. In the quaternary domain this latter fact was proved by Mertens;\footnote{Über invariante Gebilde quaternärer Formen. Wiener Berichte. Bd. 98. 1889.} our knowledge of the most general differential operations--which, as is well known, arise from the folding processes by replacing the symbol rows by differential symbols--allows us, using the decomposition identity, to transfer Mertens's proof almost verbatim to the $n$-ary domain.

\clearpage
\begin{center}
{\Large\bfseries 4. On the Invariant Theory of Forms in $n$ Variables}\par
\vspace{1.2em}
Journal für die reine und angewandte Mathematik 139 (1911), pp. 118--154
\end{center}

\section*{Introduction.}
In the projective invariant theory of forms in $n$ variables, the questions adjoining the main problem--the proof of finiteness of the system of forms--have scarcely been treated when one goes beyond the binary and ternary domains. These are the questions concerning the mutual dependence of the forms and the methods for controlling this connection among forms. These methods rest essentially--as has been completely proved in the binary and ternary domains--on two theorems designated by Study as the ``fundamental theorems of the symbolic method,'' namely the theorem on the symbolic representability of the forms and the theorem on symbolic identities.\footnote{E. Study, \emph{Methoden zur Theorie der ternären Formen} (Leipzig, Teubner, 1889). II. \S\ 6. For the occurrence of these fundamental theorems also in the invariant theory of special transformation groups, cf. Study, \emph{Das Apollonische Problem} (Math. Ann. Bd. 49 [1897]) and \emph{Zur Differentialgeometrie der analytischen Kurven} (Transactions of the American Math. Society, Bd. 1 [1909]).} The latter theorem says that all relations existing among integral invariant formations are obtained by successive application of a small number of symbolic identities; thus the symbolic method, and only it, gives knowledge of the connection among forms without leaving the domain of invariants.

Gordan, in his Erlangen program,\footnote{P. Gordan, \emph{Über das Formensystem binärer Formen}. S. 46. (Leipzig, Teubner, 1875).} already pointed to the reason for the difficulty in passing to $n$ variables: it lies in the fact that the first theorem of the symbolic method is no longer valid in its general form. In the $n$-ary domain there occur, as is well known, variable rows of higher level $p_\rho$\footnote{For the notation see \S\ 1.} and analogously symbol rows of higher level $S^\rho$; one arrives at these rows both when, as originally with Clebsch,\footnote{A. Clebsch, über eine Fundamentalaufgabe der Invariantentheorie. (Abh. der Ges. d. Wiss. zu Göttingen, Bd. XVII, 1872).} one starts from forms with arbitrarily many rows $p_\rho$, and when one starts, following F. Deruyts and A. Capelli,\footnote{Cf. A. Capelli, \emph{Lezioni sulla teoria delle forme algebriche} (Napoli 1902), \S\ 22--24, and the literature cited there.} from forms containing only arbitrarily many cogredient rows $x$. A symbolic representation that would contain the rows $p_\rho$, $S^\rho$ explicitly\footnote{More precise definition of ``explicit representation'' in \S\ 2.} does not exist; it is necessary to resolve the $p_\rho$, $S^\rho$ into the individual rows $x$ and the $s$ contragredient to them, and to distribute these among different determinants. It is indeed possible, by means of differential conditions (cf. formula (19.)), to verify that a given aggregate of determinants has the property of depending only on the rows $p_\rho$, $S^\rho$; but in this way one obtains no overview of the totality of invariant formations and of the processes for their generation.\footnote{A computationally very valuable development of the symbolism due to Study (communicated by F. Engel, Ber. d. K. Sächs. Ges. d. Wiss., math.-phys. Kl., 1900, p. 126, and for the quaternary domain by E. Study, \emph{Geometrie der Dynamen}, p. 135) is also not an explicit representation in our sense. It would, for example, hardly lead to the nonsymbolic differentiation processes for generating the forms.}

We shall now show in \S\S\ 2 and 6 how, using the ``theorem on corresponding matrices,''\footnote{Apart from Grassmann's \emph{Ausdehnungslehre} of 1862, no. 112, first in A. Brill, über zwei Berührungsprobleme, \S\ 2 (Math. Ann. Bd. 4. 1871).} the first fundamental theorem is recovered in its generality by replacing representation by determinants with representation by ``matrix products.'' In \S\ 2 it is shown that every matrix product explicitly containing the rows $S^\rho,p_\rho$ is an invariant formation of the ground form; the converse given in \S\ 6, that every invariant formation can be represented explicitly by matrix products, succeeds only through the transfer of the second fundamental theorem (\S\S\ 3--5).

This theorem is known for forms that contain only variable rows $x$.\footnote{E. Pascal, Giorn. di mat. 26 (1888), pp. 33, 102; Rendic. d. Accad. d. Linc. (4) 4 (1888), p. 119; ib. Mem. (4) 5 (1888), p. 375.} The general problem, to set up the totality of the indecomposable identities in which all rows $S^\rho,p_\rho$ are regarded as indecomposable, was formulated by Study;\footnote{``Methoden \dots'', p. 204, note 17).} at the same time he sees in the number of identities that occur a measure of the difficulty of the methods in the $n$-ary domain. Since it is now possible to combine the totality of the identities in a single formula (36.), this difficulty is at least pressed down to a minimum. In applications, however, certain distinguished special cases of (36.) will often occur, such as (20.), (21.), characterized by the fact that they admit explicit representation without substitution of new rows; or formula (31.), designated as the ``decomposition identity,'' since a row $p_\rho$ appears to be decomposed into the rows $y$.

Upon the two fundamental theorems the further theory can now easily be built. The first theorem gives directly the processes for generating the forms, the symbolic folding process and the nonsymbolic differentiation processes,\footnote{In the paper ``Invariante Gebilde von Nullsystemen'' (Sitzungsber. d. Ak. d. Wiss. Wien, Bd. 97, Abt. IIa, 1888), Mertens arrived by another path, not directly generalizable, at the quaternary folding process. The alternating invariant of 6 rows $S^2$ occurring there differs from the one arising with us by a single application of substitution (11.) by an aggregate of even invariants. For a special case the quaternary folding process is also found in P. Gordan, \emph{Das simultane System von zwei quadratischen quaternären Formen} (Math. Ann. Bd. 56. 1903).} while the decomposition identity eliminates, among these processes, all equivalent ones (\S\ 6). Knowledge of the differentiation processes further yields the transfer of the concept of ``normal form'' to the $n$-ary domain and, by means of a proof procedure given by Mertens\footnote{F. Mertens, Über invariante Gebilde quaternärer Formen. (Sitzber. d. Ak. d. Wiss. Wien. Bd. 98. Abt. IIa. 1889.)} in the quaternary domain, the theorem that a system of invariant formations can be replaced by an equivalent system of normal forms (\S\ 7). Under this latter assumption I showed in my dissertation\footnote{Über die Bildung des Formensystems der ternären biquadratischen Form. (Dieses Journal, Bd. 134. 1908.)} that the ternary folding process can be generated by successive application of two fundamental foldings. On the basis of this theorem and the theory of series expansions, I then developed, by introducing the concept of the ``form row,'' the principal reduction theorems for systems of forms. Entirely analogously, in the $n$-ary domain the folding process can be generated from $n-1$ fundamental foldings (\S\ 8), and the reduction theorems can be set up (\S\ 9).

\paragraph{Afterword.} On the occasion of the Salzburg communication, Prof. E. Müller of Vienna drew my attention to the fact that the calculation with matrix products underlying this paper is, in principle, identical with the calculus of Grassmann's theory of extension.\footnote{Hermann Graßmanns gesammelte mathematische und physikalische Werke. Ersten Bandes zweiter Teil: Die Ausdehnungslehre von 1862. In Gemeinschaft mit H. Graßmann d. Jüngeren herausgegeben v. Fr. Engel (Leipzig, Teubner 1896).} Indeed, Grassmann's ``combinatorial product'' $[S^{\rho_1}S^{\rho_2}\cdots S^{\rho_i}]$ can be regarded as a matrix given by these rows (cf. \S\ 2), specifically the ``progressive product'' as the matrix of the rows themselves, the ``regressive'' as the matrix of the assigned rows $S_{n-\rho_1},S_{n-\rho_2},\ldots,S_{n-\rho_i}$, while the matrix corresponding to the ``mixed product'' arises when several rows $S$ are combined by substitution (9.) into new rows $P,Q$. Our ``assigned row'' corresponds to Grassmann's ``complement,'' our ``matrix product'' to a ``mixed product of level number 0.''

Under this correspondence, the development given in the \emph{Ausdehnungslehre} of 1862, no. 113, becomes identical with a formula dualistic to our formula (32.). In the special case $\rho=1$, (32.) goes over into (17.), the dualistic formula into (16.). From these special cases of Grassmann's development, Müller\footnote{E. Müller, Beiträge zur Graßmannschen Ausdehnungslehre. 1. Mitteilung: Einige allgemeine Sätze. (Sitzungsber. d. Ak. d. Wiss. Wien; Bd. 118. Abt. IIa. Juli 1909), Satz 16 und 18.} has recently derived the identities (20.), (21.) mentioned above.

In contrast to the Grassmann--Müller derivation, our derivation at the same time gives the proof of completeness of the identities, which appears to be the essential point for the questions considered here, and it proceeds from (20.), (21.) further to the most general indecomposable identities (36.), which do not seem to have been noticed before now.

\section*{\S\ 1. Notations and definitions. Summary of known results.}
We denote, as usual, by the variable row $x$ the row of elements $x_1,x_2,\ldots,x_n$, and analogously by the variable row $p_\rho$ ($\rho=1,2,\ldots,n-1$) the row of the $\binom n\rho$ elements $p_{i_1i_2\ldots i_\rho}$, where $p_{i_1i_2\ldots i_\rho}$ denotes the $\binom n\rho$ determinants formed from the matrix of the $\rho$ rows $x^{(1)},x^{(2)},\ldots,x^{(\rho)}$,
\[
\left|\begin{array}{cccc}
 x_{i_1}^{(1)}&x_{i_2}^{(1)}&\cdots&x_{i_\rho}^{(1)}\\
 x_{i_1}^{(2)}&x_{i_2}^{(2)}&\cdots&x_{i_\rho}^{(2)}\\
 \vdots&\vdots&&\vdots\\
 x_{i_1}^{(\rho)}&x_{i_2}^{(\rho)}&\cdots&x_{i_\rho}^{(\rho)}
\end{array}\right|,
\qquad (i_1<i_2<\cdots<i_\rho=1,2,\ldots,n).
\]
According to the theorem on corresponding matrices,\footnote{Cf. for example Gordan--Kerschensteiner, \emph{Vorlesungen über Invariantentheorie}, Bd. I, p. 99.} to every row $p_\rho$ there is assigned one-to-one a second row $q_{n-\rho}$ by the relation, valid for each element of the row,
\begin{equation}
 p_{i_1i_2\ldots i_\rho}=(-1)^{\frac{\rho(\rho+1)}{2}+i_1+i_2+\cdots+i_\rho}q_{j_1j_2\ldots j_{n-\rho}},
\tag{1}
\end{equation}
where the $i$ and $j$ are each in good order and together form a complete permutation of the numbers 1 through $n$. The row $q_{n-\rho}$ consists of the $\binom n\rho$ determinants of degree $n-\rho$, $q_{j_1j_2\ldots j_{n-\rho}}$, formed from the matrix of $n-\rho$ rows $u^{(1)},u^{(2)},\ldots,u^{(n-\rho)}$ contragredient to the $x$. These rows satisfy the $\rho(n-\rho)$ condition equations
\begin{equation}
 (u^{(k)}\mid x^{(l)})=\sum_{\alpha=1}^{n}u_\alpha^{(k)}x_\alpha^{(l)}=0;
 \qquad (k=1,2,\ldots,n-\rho;\ l=1,2,\ldots,\rho).
\tag{2}
\end{equation}
We shall assume them normalized (with Clebsch\footnote{a. a. O. \S\ 5.}) so that the proportionality factor otherwise entering in (1.) is set equal to one.

According to Clebsch\footnote{a. a. O. \S\ 12.} (and likewise according to the later investigations of F. Deruyts and A. Capelli, cf. the Introduction), the most general forms can be reduced to those that contain at most one of each of the $n-1$ rows $p_\rho$, and hence can be represented symbolically as follows:
\begin{equation}
 (S^1\mid p_1)^{m_1}(S^2\mid p_2)^{m_2}\cdots(S^{n-1}\mid p_{n-1})^{m_{n-1}},
\tag{3}
\end{equation}
where, in analogy with (2.), we have set
\[
 (S^\rho\mid p_\rho)=\sum S^{i_1i_2\ldots i_\rho}p_{i_1i_2\ldots i_\rho}.
\]
In consequence of the nonlinear identical relations existing between the elements of a row $p_\rho$, the symbol rows $S^\rho$ can be normalized so that they satisfy the very same identities, hence so that they behave like the determinants of a $\rho$-rowed matrix whose rows $s^{(1)},s^{(2)},\ldots,s^{(\rho)}$ are contragredient to the $x$.\footnote{Clebsch, a. a. O. \S\ 4.} Thus to each row $S^\rho$ there can be assigned one-to-one another $S_{n-\rho}$ through the relation
\begin{equation}
 S_{i_1i_2\ldots i_{n-\rho}}=(-1)^{\frac{(n-\rho)(n-\rho+1)}2+i_1+i_2+\cdots+i_{n-\rho}}S^{j_1j_2\ldots j_\rho},
\tag{4}
\end{equation}
where again the $i$ and $j$ are each in good order and together exactly exhaust the numbers 1 through $n$. The elements $S_{i_1i_2\ldots i_{n-\rho}}$ of the row $S_{n-\rho}$ behave like the determinants formed from $(n-\rho)$ rows $\sigma^{(1)},\sigma^{(2)},\ldots,\sigma^{(n-\rho)}$ cogredient to the $x$, and the rows $s$ and $\sigma$ are linked by the $\rho(n-\rho)$ conditions
\begin{equation}
 (s^{(k)}\mid \sigma^{(l)})=0,\qquad(k=1,2,\ldots,\rho;\ l=1,2,\ldots,n-\rho).
\tag{5}
\end{equation}
In consequence of the relations (1.) and (4.), each factor of (3.) is capable of the equivalent fourfold symbolic representation
\begin{equation}
 (S^\rho\mid p_\rho)=(S^\rho q_{n-\rho})=(S_{n-\rho}p_\rho)=(-1)^{\rho(n-\rho)}(q_{n-\rho}\mid S_{n-\rho}),
\tag{6}
\end{equation}
where, as always in what follows, $(S^\rho q_{n-\rho})$ and $(S_{n-\rho}p_\rho)$ are abbreviations for the determinants $(s^{(1)}\cdots s^{(\rho)}u^{(1)}\cdots u^{(n-\rho)})$ and $(\sigma^{(1)}\cdots\sigma^{(n-\rho)}x^{(1)}\cdots x^{(\rho)})$; by Laplace expansion these prove to depend only on the rows $S^\rho,q_{n-\rho}$ and respectively $S_{n-\rho},p_\rho$, as the above notation is meant to indicate. The expressions $(S^\rho\mid p_\rho)$ and $(q_{n-\rho}\mid S_{n-\rho})$ mean, according to the preceding discussion, the products of the matrices of the rows $s$ and $x$, respectively of the rows $u$ and $\sigma$, where by matrix product is meant the sum over the products of corresponding determinants, that is, the value of the determinant of the product matrix.

The characteristic number $\rho(n-\rho)$ belonging to a symbol row (respectively variable row) will be called, following Clebsch, the weight of the row;\footnote{a. a. O. \S\ 11.} the number $\rho$, which gives the number of rows $s$, the upper weight index; the number $n-\rho$, which gives the number of rows $\sigma$, the lower weight index. It follows from the relations (4.) that a symbol row can occur in one and the same relation once with upper and once with lower weight index; the same holds for the occurrence of the corresponding rows $p_\rho$ and $q_{n-\rho}$. It remains to note that we generally denote variable rows of higher level whose associated matrix consists of rows $x$ by $p$, those whose matrix consists of rows $u$ by $q$; symbol rows cogredient to the $x$ by small Greek letters $\sigma,\tau,\alpha,\beta$; symbol rows cogredient to the $u$ by small Latin letters $s,t,a,b$; all other symbol rows by capital Latin letters with weight index: $S^\rho$, $T^\tau$, $A^\rho$ or $S_{n-\sigma}$, $T_{n-\tau}$, $A_{n-\rho}$. Corresponding to an upper or lower weight index, we also write the individual elements of a row with upper or lower indices, as for instance in (4.).

\section*{\S\ 2. Representation by matrix products.}
In what follows we always understand by ``matrix product'' the sum over the products of corresponding determinants of two matrices with the same number of rows, where the rows of the second matrix are contragredient to those of the first. The rows of each matrix may be: 1. individually given; 2. combined, as in (6.), into a row $S^\rho$; 3. united into different rows $S^\rho,S^\sigma\ldots$. The matrix product is to be denoted by the symbol $(\ \mid\ )$ used in (6.). The representation (6.) is a special application of a second formulation of the theorem on corresponding matrices: ``To every $\rho$-rowed matrix product $(v^{(1)}v^{(2)}\cdots v^{(\rho)}\mid y^{(1)}\cdots y^{(\rho)})$ there is assigned a determinant; and conversely to every determinant there can be assigned a matrix product with prescribed number $\rho$ of rows ($\rho=1,2,\ldots,n-1$).''\footnote{For $\rho>n$ the matrix product is known to vanish; for $\rho=n$ it goes into the product of two determinants.} Indeed, one has only to combine the rows $y$ into a row $p_\rho$ (respectively the rows $v$ into a row $q_\rho$) and to carry out the substitutions given by (1.) in order to pass from a given matrix product to the determinant and conversely. The value of the determinant is, however, not given by its individual elements, but first results from Laplace expansion. Since determinant and matrix product therefore prove to be equivalent, the theorem on the symbolic representability of the forms (the first fundamental theorem) is expressed also as follows:

\emph{Theorem I.} ``All invariant formations can be represented symbolically by matrix products, and conversely all matrix products are invariant formations.''\footnote{Invariant formations of prescribed ground forms are, of course, only those aggregates of matrix products that depend on the symbol rows $S^\rho\ldots$.}

Representation by matrix products allows us to overcome the essential difficulty of the non-explicit representation caused by the determinant representation.\footnote{Cf. the Introduction.} By ``explicit representation'' we mean a representation into which only the rows $S^\rho,T^\tau,\ldots$ themselves enter, or rows $R^\rho$ uniquely derivable from these, but in which the individual component rows $s,t,\ldots$ no longer occur. We shall obtain in \S\ 6, for all invariant formations that depend only on the symbol rows $S^\rho,T^\tau,\ldots$, a representation explicit in these symbol rows, and thereby the transfer to the $n$-ary domain of the generating processes, the folding process and the differentiation processes. The possibility of explicit representation is evident from the fact that only the simplest matrix products can be assigned precisely to the determinants occurring in the determinant representation, so that all other matrix products are combinations of different determinants into a single expression. For if one resolves the symbol and variable rows into the individual rows $s,u,x$, then, according to the first fundamental theorem in its usual form, an aggregate of determinants must necessarily arise. The proof that, conversely, all determinant aggregates that represent invariant formations of prescribed ground forms can be combined into explicit matrix products follows by using the second fundamental theorem (\S\ 6).

In order to obtain from two symbol rows $S^\sigma,T^\tau$, with the help of variable rows, an invariant formation of the ground form $(S^\sigma\mid p_\sigma)(T^\tau\mid p_\tau)$ that explicitly contains all rows (the folding process), we need only form a matrix product of the third kind according to the definition given at the beginning of the paragraph:
\begin{equation}
 \pair{q_{n-\sigma-\lambda}S^\sigma}{T_{n-\tau}p_{\tau-\lambda}},
 \qquad (\lambda=1,2,\ldots,\tau,n-\sigma).
\tag{7}
\end{equation}
Expanded, (7.) gives:
\begin{equation}
\begin{aligned}
&\pair{v^{(1)}\cdots v^{(n-\sigma-\lambda)}s^{(1)}\cdots s^{(\sigma)}}
       {\alpha^{(1)}\cdots\alpha^{(n-\tau)}y^{(1)}\cdots y^{(\tau-\lambda)}} \\
&\qquad =\sum_{i,k,l}\eps\,q_{k_1\ldots k_{n-\sigma-\lambda}}
 S^{k_{n-\sigma-\lambda+1}\ldots k_{n-\lambda}}
 T_{l_1\ldots l_{n-\tau}}
 p^{l_{n-\tau+1}\ldots l_{n-\lambda}},
\end{aligned}
\tag{8}
\end{equation}
that is, an expression depending only on the rows $S,T,q,p$. Here $\eps=\pm1$\footnote{This notation is to be retained throughout.} and the sum is to be understood in such a way that the indices of each row $S\ldots$ are in good order, and that for each combination $i_1,i_2,\ldots,i_{n-\lambda}$ the $k$ as well as the $l$ individually run through all then possible permutations of the numbers $i$. The number $\lambda$ entering into (7.) will be called the defect of the folding in the case $\sigma>\tau$; the reason for this latter restriction appears in \S\ 6.\footnote{The number $\lambda$ always runs up to the smaller of the two last indicated numbers.}

The determinants of the two matrices entering in (7.) can be regarded as elements of new rows $Q^{n-\lambda},P_{n-\lambda}$ whose associated rows are $Q_\lambda,P^\lambda$. According to (8.), these rows $Q,P$ can be expressed through the original $S$ and $q$, respectively $T$ and $p$. We indicate this by the equations
\begin{equation}
\begin{aligned}
 q_{n-\sigma-\lambda}S^\sigma&=Q^{n-\lambda}\sim Q_\lambda, &\quad\text{or also}\quad Q_\lambda&=[q_{n-\sigma-\lambda}S^\sigma]_\lambda,\\
 T_{n-\tau}p_{\tau-\lambda}&=P_{n-\lambda}\sim P^\lambda, &\quad\text{or also}\quad P^\lambda&=[T_{n-\tau}p_{\tau-\lambda}]^\lambda.
\end{aligned}
\tag{9}
\end{equation}
Taking (4.) into account, one then has, analogously to (6.), the equivalent fourfold symbolic representation for the matrix product (7.):
\begin{equation}
\pair{q_{n-\sigma-\lambda}S^\sigma}{T_{n-\tau}p_{\tau-\lambda}}
=(q_{n-\sigma-\lambda}S^\sigma P^\lambda)=(Q_\lambda T_{n-\tau}p_{\tau-\lambda})=(-1)^{\lambda(n-\lambda)}\pair{P^\lambda}{Q_\lambda}.
\tag{10}
\end{equation}
A further substitution of new rows can be performed on (7.) by setting
\begin{equation}
 \pair{q_{n-\sigma-\lambda}S^\sigma}{T_{n-\tau}p_{\tau-\lambda}}
 =\pair{R^{\sigma+\lambda}}{p_{\sigma+\lambda}}(R^{\tau-\lambda}p_{\tau-\lambda}).
\tag{11}
\end{equation}
Here too (8.) shows that the products of the elements of the rows $R$ can be expressed uniquely through the rows $S$ and $T$. The substitutions (9.) and (11.) will be used especially when a further folding with other rows is performed with (7.).

The matrix representation enables us to state invariantly the quadratic relations between the elements of a row $p_\rho$ (from which, as is known, all the others follow), hence, according to \S\ 1, the relations linear in the coefficients of the ground form that the rows $S^\rho$ satisfy. These are the identities
\begin{equation}
 \pair{q_{n-\rho-\lambda}S^\rho}{S_{n-\rho}p_{\rho-\lambda}}=0,
 \qquad (\lambda=1,2,\ldots,\rho,n-\rho),
\tag{12}
\end{equation}
where the $p$ and $q$ are to be regarded as arbitrary quantities. We shall see in \S\ 5 that the identities with odd defect number $\lambda$ are consequences of those with higher defect. The relations (12.) are a direct consequence of the relations (5.); for one has
\[
 \pair{q_{n-\rho-\lambda}S^\rho}{S_{n-\rho}p_{\rho-\lambda}}
 =(v^{(1)}\cdots v^{(n-\rho-\lambda)}s^{(1)}\cdots s^{(\rho)}\mid \sigma^{(1)}\cdots\sigma^{(n-\rho)}y^{(1)}\cdots y^{(\rho-\lambda)}),
\]
and the application of the product theorem gives, in consequence of (5.), a vanishing $(n-\lambda)$-rowed determinant. The identities (12.) say that in order to find all types of foldings it suffices to fold the product of forms linear in each variable row; or also, by \S\ 6, that the normalized form $(S^\rho\mid p_\rho)^m$ possesses no invariant formation linear in the coefficients.\footnote{Cf. an addendum by Study in Grassmann's works, first volume, second part, p. 510 (condition that a quantity $S^\rho$ is simple).}

The conditions (12.) are also sufficient to characterize the rows $p_\rho$. For the property that the individual elements of $p_\rho$ are determinants of a matrix is an invariant property, and therefore it must be expressible invariantly; but by Theorem VI, the conditions (12.) give the greatest possible invariant restriction for the $p_\rho$, namely that the linear form formed from the row $p_\rho$ as coefficients possesses no invariant formations at all.

\section*{\S\ 3. The symbolic identities.}
We now have to transfer the second fundamental theorem of the symbolic method, namely to set up the totality of the irreducible indecomposable identities in which all rows $S^\rho,p_\rho$ are regarded as indecomposable. We also stipulate the following. Corresponding to the meaning of symbol and variable rows, those identities are to be regarded as composite that arise when one specializes the variable row $p_\rho$ to $p_{\rho-1}y$ and applies one of the identities of the system to the resulting expressions; while, on the other hand, identities containing $k$ symbol rows are to count as indecomposable even if they arise from identities with fewer symbol rows by the above process. The rows $S^\rho,p_\rho$ need not necessarily enter the identities explicitly; for their interpretation as indecomposable quantities it suffices to show that the matrix products that occur depend only on these rows. We shall show, however, that in this case, partly by applying substitution (11.), an explicit representation can always be obtained. The general identities are obtained, by means of the principle of matrix representation, from the special identities given by Pascal, which contain only rows $x$ and $u$ and constitute a direct transfer of those given by Study for the ternary domain:\footnote{For formulation and literature of the problem, cf. the Introduction.}
\begin{align}
\sum_{i=1}^n(-1)^{i+1}\pair{a^{(i)}}{x}(a^{(1)}\cdots a^{(i-1)}a^{(i+1)}\cdots a^{(n)}u)&=(-1)^{n+1}\pair{u}{x}(a^{(1)}\cdots a^{(n)}),\tag{13}\\
\sum\pm\pair{a^{(1)}}{x^{(1)}}\pair{a^{(2)}}{x^{(2)}}\cdots\pair{a^{(n)}}{x^{(n)}}&=(a^{(1)}\cdots a^{(n)})(x^{(1)}\cdots x^{(n)}),\tag{14}\\
\sum_{i=1}^n(-1)^{i+1}\pair{u}{a^{(i)}}(a^{(1)}\cdots a^{(i-1)}a^{(i+1)}\cdots a^{(n)}x)&=(-1)^{n+1}\pair{u}{x}(a^{(1)}\cdots a^{(n)}).\tag{15}
\end{align}
Two further identities arise from (13.) and (15.) through the substitutions $(a^{(i)}\mid x)=(a^{(i)}q_{n-1})$ and respectively $(u\mid a^{(i)})=(p_{n-1}a^{(i)})$, and so by our interpretation they are not to be regarded as essentially different. (14.) represents the product theorem for determinants; we shall show how (13.) and (14.) (and, dualistically, (15.) and (14.)) can be replaced by the system of suitably formed product theorems for matrices. The product theorem, formed once for $\rho$-rowed matrices and once for $(\rho-1)$-rowed matrices, reads
\[
\pair{a^{(1)}a^{(2)}\cdots a^{(\rho)}}{xp_{\rho-1}}
=(a^{(1)}a^{(2)}\cdots a^{(\rho)}\mid xy^{(2)}\cdots y^{(\rho)})
=\sum\pm\pair{a^{(1)}}{x}\pair{a^{(2)}}{y^{(2)}}\cdots\pair{a^{(\rho)}}{y^{(\rho)}},
\]
\[
\sum\pm\pair{a^{(2)}}{y^{(2)}}\cdots\pair{a^{(\rho)}}{y^{(\rho)}}
=(a^{(2)}\cdots a^{(\rho)}\mid y^{(2)}\cdots y^{(\rho)})
=(a^{(2)}\cdots a^{(\rho)}p_{\rho-1})=(a^{(2)}\cdots a^{(\rho)}q_{n-\rho+1});
\]
and from this follows the system of product theorems, where $\rho=2,3,\ldots,n$:
\begin{equation}
\begin{aligned}
\pair{a^{(1)}a^{(2)}\cdots a^{(\rho)}}{xp_{\rho-1}}
 &=\sum_{i=1}^{\rho}(-1)^{i+1}\pair{a^{(i)}}{x}(a^{(1)}\cdots a^{(i-1)}a^{(i+1)}\cdots a^{(\rho)}\mid p_{\rho-1})\\
 &=\sum_{i=1}^{\rho}(-1)^{i+1}\pair{a^{(i)}}{x}(a^{(1)}\cdots a^{(i-1)}a^{(i+1)}\cdots a^{(\rho)}q_{n-\rho+1}).
\end{aligned}
\tag{16}
\end{equation}
For $\rho=n$, (16.) goes over into (13.); if we further specialize $p_{n-1}=y^{(2)}p_{n-2}$, $p_{n-2}=y^{(3)}p_{n-3}$, etc., and apply the identities (16.) to the expressions $(a^{(1)}\cdots a^{(i-1)}a^{(i+1)}\cdots a^{(n)}\mid y^{(2)}p_{n-2})$, $(a^{(1)}\cdots a^{(i-1)}a^{(i+1)}\cdots a^{(h-1)}a^{(h+1)}\cdots a^{(n)}\mid y^{(3)}p_{n-3})$, etc., then we pass from (13.) to (14.). Thus the identity (14.) is, by our definition, composite from the identities (16.); equivalently, the system (16.) is equivalent to the identities (13.) and (14.). The system (16.) satisfies one of the conditions required for the general identities: it contains the row $p_{\rho-1}$ as an indecomposable quantity and at the same time in explicit form. It is the only system satisfying this condition and only this condition. For we can form no further determinants or matrix products from the rows $a,x,p$; and between these no relations independent of (16.) can exist, since otherwise, by eliminating $(a^{(1)}\cdots a^{(\rho)}\mid xp_{\rho-1})$, there would arise a relation between $\rho$ ($\rho<n$) expressions $(a^{(i)}\mid x)(a^{(1)}\cdots a^{(i-1)}a^{(i+1)}a^{(\rho)}q_{n-\rho+1})$, whereas by (13.) at least $n+1$ such expressions must enter into a relation. From analogous considerations one obtains the system equivalent to (15.) and (14.):
\begin{equation}
\begin{aligned}
\pair{u q_{\rho-1}}{a^{(1)}\cdots a^{(\rho)}}
 &=\sum_{i=1}^{\rho}(-1)^{i+1}\pair{u}{a^{(i)}}(q_{\rho-1}\mid a^{(1)}\cdots a^{(i-1)}a^{(i+1)}\cdots a^{(\rho)})\\
 &=\sum_{i=1}^{\rho}(-1)^{i+1}\pair{u}{a^{(i)}}(p_{n-\rho+1}a^{(1)}\cdots a^{(i-1)}a^{(i+1)}\cdots a^{(\rho)}).
\end{aligned}
\tag{17}
\end{equation}
In order to pass from (16.) to identities between arbitrary rows $A^{\rho_1},B^{\rho_2},\ldots,x,p$, we observe that these identities must become identical with (16.) as soon as we resolve
$A^{\rho_1}=a^{(1)}a^{(2)}\cdots a^{(\rho_1)}$, $B^{\rho_2}=b^{(1)}b^{(2)}\cdots b^{(\rho_2)}$, etc., and that they can also be combined into final expressions only by the operations determined by (16.) as soon as the individual expressions are of the type entering into (16.). Thus we obtain, if we also set
\begin{equation}
 \rho=\rho_1+\rho_2+\cdots+\rho_k\le n,
\tag{18}
\end{equation}
\begin{align*}
&\pair{a^{(1)}\cdots a^{(\rho_1)}b^{(1)}\cdots b^{(\rho_2)}\cdots k^{(1)}\cdots k^{(\rho_k)}}{xp_{\rho-1}}
  =\pair{A^{\rho_1}B^{\rho_2}\cdots K^{\rho_k}}{xp_{\rho-1}} \\
&\quad=\sum_{i=1}^{\rho_1}(-1)^{i+1}\pair{a^{(i)}}{x}(a^{(1)}\cdots a^{(i-1)}a^{(i+1)}\cdots a^{(\rho_1)}B^{\rho_2}\cdots K^{\rho_k}q_{n-\rho+1})\\
&\qquad+(-1)^{\rho_1\rho_2}\sum_{i=1}^{\rho_2}(-1)^{i+1}\pair{b^{(i)}}{x}(b^{(1)}\cdots b^{(i-1)}b^{(i+1)}\cdots b^{(\rho_2)}A^{\rho_1}\cdots K^{\rho_k}q_{n-\rho+1})\\
&\qquad+\cdots\\
&\qquad+(-1)^{(\rho_1+\cdots+\rho_{k-1})\rho_k}\sum_{i=1}^{\rho_k}(-1)^{i+1}\pair{k^{(i)}}{x}(k^{(1)}\cdots k^{(i-1)}k^{(i+1)}\cdots k^{(\rho_k)}A^{\rho_1}\cdots K^{\rho_{k-1}}q_{n-\rho+1}).
\end{align*}
Each individual sum (but not already a part of the sum) does in fact contain the rows $A^{\rho_1},B^{\rho_2},\ldots$ as indecomposable quantities; for it satisfies the necessary and sufficient conditions\footnote{Capelli, a. a. O. \S\ 23, gives sufficient conditions: $\bigl(a^{(k)}\mid \frac{\partial}{\partial a^{(i)}}\bigr)=0$, $(k>i,\ i=1,2,\ldots,\rho_1-1)$. These most simply imply the necessary and sufficient conditions by application of the Poisson bracket process after Wellstein, Von den Differentialgleichungen der projektiven Invarianten, Math. Ann. 67 (1909) \S\ 3.}
\begin{equation}
 \left(a^{(i+1)}\middle|\frac{\partial}{\partial a^{(i)}}\right)=0;\ \ldots\ ;
 \left(k^{(j+1)}\middle|\frac{\partial}{\partial k^{(j)}}\right)=0.
 \quad (i=1,2,\ldots,\rho_1-1;\ j=1,2,\ldots,\rho_k-1)
\tag{19}
\end{equation}
We show that it also contains all rows explicitly. For if we set $B^{\rho_2}\cdots K^{\rho_k}q_{n-\rho+1}=q_{n-\rho_1+1}$, then by (16.) and (10.) we obtain
\begin{align*}
\sum(-1)^{i+1}\pair{a^{(i)}}{x}(a^{(1)}\cdots a^{(i-1)}a^{(i+1)}\cdots a^{(\rho_1)}q_{n-\rho+1})
&=\pair{A^{\rho_1}}{xp_{\rho_1-1}}\\
&=(A_{n-\rho_1}xp_{\rho_1-1})=(-1)^{(\rho_1+1)(n+1)}\pair{q_{n-\rho_1+1}}{A_{n-\rho_1}x}\\
&=(-1)^{(\rho_1+1)(n+1)}\pair{B^{\rho_2}\cdots K^{\rho_k}q_{n-\rho+1}}{A_{n-\rho_1}x}.
\end{align*}
If we compute the remaining sums analogously and subsequently set, since the rows are already characterized as distinct by the weight indices, $B^{\rho_2}=A^{\rho_2}$, $C^{\rho_3}=A^{\rho_3},\ldots$, then we obtain the desired identities:
\begin{equation}
\begin{aligned}
\pair{A^{\rho_1}A^{\rho_2}\cdots A^{\rho_k}}{xp_{\rho-1}}
&=\sum_{i=1}^{k}(-1)^{\delta_i}\pair{A^{\rho_1}\cdots A^{\rho_{i-1}}A^{\rho_{i+1}}\cdots A^{\rho_k}q_{n-\rho+1}}{A_{n-\rho_i}x},\\
\delta_i&=(\rho_1+\rho_2+\cdots+\rho_{i-1})\rho_i+(\rho_i+1)(n+1),
\end{aligned}
\tag{20}
\end{equation}
and from (17.) the dualistic identities:
\begin{equation}
\begin{aligned}
\pair{u q_{\sigma-1}}{A_{\sigma_1}A_{\sigma_2}\cdots A_{\sigma_k}}
&=\sum_{i=1}^{k}(-1)^{\eps_i}\pair{u A^{n-\sigma_i}}{p_{n-\sigma+1}A_{\sigma_1}\cdots A_{\sigma_{i-1}}A_{\sigma_{i+1}}\cdots A_{\sigma_k}},\\
\eps_i&=(\sigma_1+\sigma_2+\cdots+\sigma_{i-1})\sigma_i+(\sigma_i+1)(n+\sigma).
\end{aligned}
\tag{21}
\end{equation}
\emph{Theorem II.} With (20.) and (21.) the totality of the indecomposable identities between the rows $A^{\rho_i},x,p$, respectively $A_{\sigma_i},u,q$, is exhausted, and at the same time the totality of those indecomposable identities that admit explicit representation without carrying out substitution (11.).

The first part of the theorem follows from the considerations leading to (20.) and (21.); the second part from the impossibility of an identity between two symbol rows:
\[
\pair{q_\rho A^\sigma}{B_{\rho+\sigma-\tau}p_\tau}
=\eps_1\pair{q_\rho B^{n-\rho-\sigma+\tau}}{A_{n-\sigma}p_\tau}
+\eps_2\pair{q_\rho q_{n-\tau}}{B_{\rho+\sigma-\tau}A_{n-\sigma}},
\]
where $\rho\le\tau\le\sigma$ and none of the rows has weight $1\cdot(n-1)$. (Other assumptions about $\rho,\tau,\sigma$ can be reduced to this one by interchanging the row designations.) This impossibility follows, for example, by expanding (8.) under the specialization $q_\rho=lM^{\rho-1}$, $q_{n-\tau}=lN^{n-\tau-1}$, if one sets $l_1=1$, $M^{2,3,\ldots,\rho}=1$, $A^{\rho+1,\ldots,\rho+\sigma}=1$, all other elements of the rows $l,M,A$ equal to zero, while $N$ and $B$ are arbitrary. Since identities between arbitrary rows, by resolving a row $p_\tau$ into $y^{(1)}y^{(2)}\cdots y^{(\tau)}$, pass into identities composite from (20.), the second part of the assertion is proved as soon as (20.) can be generated from identities with only two symbol rows. In fact, from
\begin{equation}
\pair{A^{\rho_1}A^{\rho_2}}{xp_{\rho-1}}
=\eps\pair{A^{\rho_1}q_{n-\rho+1}}{A_{n-\rho_2}x}
+\pair{A^{\rho_1}}{xp_{\rho_1-1}},
\quad \text{where }p_{\rho_1-1}\sim A^{\rho_2}q_{n-\rho+1},
\tag{20a}
\end{equation}
by specializing $A^{\rho_1}=B^{\sigma_1}B^{\sigma_2}$, that is, by
\[
\pair{B^{\sigma_1}B^{\sigma_2}}{xp_{\rho_1-1}}
=\eps\pair{B^{\sigma_1}q_{n-\rho_1+1}}{B_{n-\sigma_2}x}
+\pair{B^{\sigma_1}}{xp_{\sigma_1-1}},
\quad \text{where } p_{\sigma_1-1}\sim B^{\sigma_2}q_{n-\rho_1+1},
\]
one obtains an identity (20.) in three symbol rows; and by further specializations of $B^{\sigma_1}=C^{\tau_1}C^{\tau_2}$, etc., the general identities (20.). The nonexistence of the explicit representation of an identity between two symbol rows and two arbitrary variable rows thus implies the nonexistence in the case of arbitrarily many symbol rows, provided that among the variable rows no row $x$ or $u$ is contained.

\clearpage
\section*{\S\ 4. The decomposition identities.}
We now consider the most general case of identities among arbitrary symbol and variable rows in whose final expressions, without performing substitution (11.), a row $p_\tau$ occurs resolved into its individual rows $y^{(1)}\cdots y^{(\tau)}$; we call these identities ``decomposition identities.'' Following the remark at the end of the preceding paragraph, we first determine them for the case of only two symbol rows, that is, we expand the expression $\pair{q_\rho A^\sigma}{B_{\rho+\sigma-\tau}p_\tau}$. We set:
\begin{equation}
\begin{aligned}
q_{\rho-\alpha}^{(i_1\ldots i_\alpha)}&\sim p_{n-\rho}\,y^{(i_1)}\cdots y^{(i_\alpha)},\qquad
p_{\tau-\alpha}^{(i_1\ldots i_\alpha)}=y^{(i_{\alpha+1})}\cdots y^{(i_\tau)},\qquad
p_\tau=y^{(1)}\cdots y^{(\tau)}.
\end{aligned}
\tag{22}
\end{equation}
With respect to the numbers $\rho,\sigma,\tau$ four cases must be distinguished:
\[
\begin{array}{llll}
1)&\rho+\sigma\le n,& \tau\le\rho,& \tau\le\sigma;\\[2pt]
2)&\rho+\sigma\le n,& \tau\ge\rho,& \tau\le\sigma;\\[2pt]
3)&\rho+\sigma\le n,& \tau\le\rho,& \tau>\sigma;\\[2pt]
4)&\rho+\sigma\le n,& \tau\ge\rho,& \tau\ge\sigma.
\end{array}
\]
We shall single out the first case, and then briefly characterize the others from it. By (20.), (10.) and (9.) we obtain, when we decompose the row $y^{(1)}y^{(2)}\cdots y^{(\tau)}B_{\rho+\sigma-\tau}$ into $y^{(1)}$ and $y^{(2)}\cdots y^{(\tau)}B_{\rho+\sigma-\tau}$:
\begin{equation}
\begin{aligned}
\pair{q_\rho A^\sigma}{y^{(1)}y^{(2)}\cdots y^{(\tau)}B_{\rho+\sigma-\tau}}
&=\pair{q_\rho^{(1)}A^\sigma}{y^{(2)}\cdots y^{(\tau)}B_{\rho+\sigma-\tau}}\\
&\quad+(-1)^\rho\pair{q_\rho[A_{n-\sigma}y^{(1)}]^{\sigma-1}}{y^{(2)}\cdots y^{(\tau)}B_{\rho+\sigma-\tau}}.
\end{aligned}
\tag{23}
\end{equation}
The last term in (23.) now has exactly the form of the original expression--only $\sigma$ is replaced by $\sigma-1$, and $\tau$ by $\tau-1$--and therefore again admits equation (23.). Thus, since $\tau\le\sigma$, we can apply operation (23.) $\tau$ times and, taking account of (10.), obtain the relation:
\begin{equation}
\begin{aligned}
\pair{q_\rho A^\sigma}{p_\tau B_{\rho+\sigma-\tau}}
&=(-1)^{\rho\sigma+(\sigma+\tau)(n+1)}
  \pair{q_\rho B^{n-\rho-\sigma+\tau}}{A_{n-\sigma}p_\tau}\\
&\quad+\sum_{i_1=1}^{\tau}(-1)^{\rho(i_1-1)}
 \pair{q_{\rho-1}^{(i_1)}[A_{n-\sigma}y^{(1)}\cdots y^{(i_1-1)}]^{\sigma-i_1+1}}
 {y^{(i_1+1)}\cdots y^{(\tau)}B_{\rho+\sigma-\tau}}.
\end{aligned}
\tag{24}
\end{equation}
Every term under the summation sign is again of the type of the original expression; $\rho$ is replaced by $\rho-1$, $\sigma$ by $\sigma-i_1+1$, and $\tau$ by $\tau-i_1$, so that condition 1) is even more strongly satisfied. Since $\tau\le\rho$, we can therefore apply operation (24.) $\tau$ times and obtain the desired decomposition identity:
\begin{equation}
\begin{aligned}
\pair{q_\rho A^\sigma}{p_\tau B_{\rho+\sigma-\tau}}
&=\eps_{i_0}\pair{q_\rho B^{n-\rho-\sigma+\tau}}{A_{n-\sigma}p_\tau}\\
&\quad+\sum_{i_1=1}^{\tau}\eps_{i_1}
 \pair{q_{\rho-1}^{(i_1)}B^{n-\rho-\sigma+\tau}}{A_{n-\sigma}p_{\tau-1}^{(i_1)}}\\
&\quad+\sum_{i_2=i_1+1}^{\tau}\eps_{i_2}
 \pair{q_{\rho-2}^{(i_1i_2)}B^{n-\rho-\sigma+\tau}}{A_{n-\sigma}p_{\tau-2}^{(i_1i_2)}}+\cdots\\
&\quad+\eps_{i_\tau}\pair{q_{\rho-\tau}^{(i_1\ldots i_\tau)}B^{n-\rho-\sigma+\tau}}{A_{n-\sigma}}\\
&=\sum_{\alpha=0}^{\tau}\sum_{i_\alpha=i_{\alpha-1}+1}^{\tau}\eps_{i_\alpha}
 \pair{q_{\rho-\alpha}^{(i_1\ldots i_\alpha)}B^{n-\rho-\sigma+\tau}}{A_{n-\sigma}p_{\tau-\alpha}^{(i_1\ldots i_\alpha)}}.
\end{aligned}
\tag{25}
\end{equation}
For the sign factor $\eps$ one obtains from (24.):
\begin{equation}
\begin{aligned}
\eps_{i_\alpha}&=(-1)^{\delta_{i_\alpha}},\qquad
\delta_{i_\alpha}=\sum_{\beta=1}^{\alpha}(\rho-\beta+1)(i_{\beta-1}+i_\beta+1)\\
&\quad +(\rho-\alpha)(\sigma+i_\alpha+\alpha)+(\sigma+\tau+\alpha)(n+1),
\end{aligned}
\tag{26}
\end{equation}
where $i_0=0$ is to be put. The summation sign in (25.) is to be understood as follows: to each combination $i_1i_2\ldots i_{\alpha-1}$, where $i_1<i_2<\cdots<i_{\alpha-1}$, all values $i_\alpha$ for which $i_{\alpha-1}<i_\alpha\le\tau$ are adjoined. Taking (26.) into account, one easily verifies that the individual sums in (25.) satisfy the differential equations analogous to (19.):
\begin{equation}
 \left(\frac{\partial}{\partial y^{(i)}}\middle|y^{(i+1)}\right)=0,
 \qquad (i=1,2,\ldots,\tau-1),
\tag{27}
\end{equation}
and hence depend only on the row $p_\tau$. Thus the decomposition identity is an indecomposable identity among the rows $A,B,q,p$; we shall return to its explicit representation in the next paragraph.

In case 2), operation (23.) can likewise be applied $\tau$ times, but operation (24.) stops after the $\rho$-th application. Thus in the final expression of (25.) the first summation sign must be replaced by $\sum_{\alpha=0}^{\rho}$. In cases 3) and 4), operation (23.) can be carried out only $\sigma$ times; in place of (24.) one then has the relation:
\begin{equation}
\begin{aligned}
\pair{q_\rho A^\sigma}{p_\tau B_{\rho+\sigma-\tau}}
&=(-1)^{\rho\sigma}
 \pair{q_\rho^{(\sigma+1\ldots\tau)}B^{n-\rho+\tau-\sigma}}{A_{n-\sigma}p_{\tau-\sigma}^{(\sigma+1\ldots\tau)}}\\
&\quad+\sum_{i_1=1}^{\sigma}(-1)^{\rho(i_1-1)}
 \pair{q_{\rho-1}^{(i_1)}[A_{n-\sigma}y^{(1)}\cdots y^{(i_1-1)}]^{\sigma-i_1+1}}
 {y^{(i_1+1)}\cdots y^{(\tau)}B_{\rho+\sigma-\tau}};
\end{aligned}
\tag{28}
\end{equation}
repeating (28.) $\tau-\sigma$ times, where the summation sign becomes $\sum_{i_\beta=i_{\beta-1}+1}^{\sigma+\beta-1}$--as the $(\sigma-i_1+1)$-fold application of (23.) to the terms under the summation sign in (28.) shows--gives all terms of the individual sum in (25.) corresponding to the index $\alpha=\tau-\sigma$, with the corresponding sign. Then cases 3) and 4) pass over into cases 1) and 2), for the numbers corresponding to $\sigma,\tau$ become $\sigma'=\sigma-i_{\tau-\sigma}+\tau-\sigma$, $\tau'=\tau-i_{\tau-\sigma}=\sigma'$; and as the procedure continues, $\tau'<\sigma'$. Thus in the final expression of (25.) the first summation sign must be replaced by $\sum_{\alpha=\tau-\sigma}^{\tau,\rho}$.

Analogously, from (21.) one obtains the dual decomposition identity, in which a row $q_\rho$ is resolved into the individual rows $v^{(1)},\ldots,v^{(\rho)}$. If we put
\begin{equation}
 \dot p_{\tau-\beta}^{(i_1\ldots i_\beta)}\sim v^{(i_1)}\cdots v^{(i_\beta)}q_{n-\tau};\qquad
 \dot q_{\rho-\beta}^{(i_1\ldots i_\beta)}=v^{(i_{\beta+1})}\cdots v^{(i_\rho)};\qquad
 q_\rho=v^{(1)}\cdots v^{(\rho)},
\tag{29}
\end{equation}
then
\begin{equation}
 \pair{q_\rho A^\sigma}{p_\tau B_{\rho+\sigma-\tau}}
 =\sum_{\beta=\max(0,\tau-\sigma)}^{\min(\tau,\rho)}
   \sum_{i_\beta=i_{\beta-1}+1}^{\rho}\eps
\pair{\dot q_{\rho-\beta}^{(i_1\ldots i_\beta)}B^{n-\rho-\sigma+\tau}}{A_{n-\sigma}\dot p_{\tau-\beta}^{(i_1\ldots i_\beta)}} ,
\tag{30}
\end{equation}
where the summation index $\beta$ runs from the larger of the lower values to the smaller of the upper values. All terms of an individual sum in (25.) and (30.) are polars of a single form, produced by the polar processes
\[
\left(\frac{\partial}{\partial p_{\tau-\alpha}}\middle|p_{\tau-\alpha}^{(i_1\ldots i_\alpha)}\right),\qquad
\left(q_{\rho-\alpha}^{(i_1\ldots i_\alpha)}\middle|\frac{\partial}{\partial q_{\rho-\alpha}}\right).
\]
To express this, we denote by $\Pi$ an aggregate of such polar operations, multiplied by constants, and obtain, in place of (25.) and (30.), the relation
\begin{equation}
 \pair{q_\rho A^\sigma}{p_\tau B_{\rho+\sigma-\tau}}
 =\sum_{\alpha=\max(0,\tau-\sigma)}^{\min(\tau,\rho)}
 \Pi\pair{q_{\rho-\alpha}B^{n-\rho-\sigma+\tau}}{A_{n-\sigma}p_{\tau-\alpha}}.
\tag{31}
\end{equation}

\section*{\S\ 5. The decomposition identities in explicit form.}
If in (25), in case 4), we specialize $\tau$ to $\sigma+\rho$, we obtain
\begin{equation}
 \pair{q_\rho A^\sigma}{y^{(1)}\cdots y^{(\rho+\sigma)}}
 =\sum_{1\le i_1<\cdots<i_\rho\le \rho+\sigma}
 \eps_{i_\rho}\pair{q_\rho}{y^{(i_1)}\cdots y^{(i_\rho)}}
 \pair{A^\sigma}{y^{(i_{\rho+1})}\cdots y^{(i_{\rho+\sigma})}},
\tag{32}
\end{equation}
where $\eps_{i_\rho}=(-1)^{\delta_{i_\rho}}$ and $\delta_{i_\rho}=\rho(\rho+1)/2+i_1+\cdots+i_\rho$. This is a more general form of the product theorem for matrices than (16.) and (17.), and follows by Laplace expansion of the determinant of the product matrix
\[
 \sum\pm(v^{(1)}y^{(1)})\cdots(v^{(\rho)}y^{(\rho)})(a^{(1)}y^{(\rho+1)})\cdots(a^{(\sigma)}y^{(\rho+\sigma)}).
\]
Thus the more general product theorem is composed out of the identities (20.) and (21.), respectively, and the selection of the special form (16.), (17.) is thereby explained.

Relation (32.) now gives the means for explicitly representing the decomposition identities. For this purpose we regard the forms occurring in (31.) as ground forms; that is, we perform substitution (11.), which changes nothing in the explicit representation in the rows $A,B$, since by (8.) the newly introduced rows $R$ can be expressed through the rows $A,B$ themselves and not through their component rows $a^{(1)}a^{(2)}\ldots$, $b^{(1)}b^{(2)}\ldots$. Thus let
\begin{equation}
 \pair{q_{\rho-\alpha}B^{n-\rho-\sigma+\tau}}{A_{n-\sigma}p_{\tau-\alpha}}
 =\pair{R_{\rho-\alpha}p_{n-\rho+\alpha}}{R^{\tau-\alpha}p_{\tau-\alpha}}.
\tag{33}
\end{equation}
Then, for the individual sum of (25.) corresponding to the index $\alpha$, we get
\begin{equation}
\begin{aligned}
&\sum \eps_{i_\alpha}\pair{R_{\rho-\alpha}p_{n-\rho}y^{(i_1)}\cdots y^{(i_\alpha)}}{R^{\tau-\alpha}y^{(i_{\alpha+1})}\cdots y^{(i_\tau)}}\\
&\quad=\sum \eps_{i_\alpha}\pair{[R_{\rho-\alpha}p_{n-\rho}]^\alpha y^{(i_1)}\cdots y^{(i_\alpha)}}{R^{\tau-\alpha}y^{(i_{\alpha+1})}\cdots y^{(i_\tau)}}\\
&\quad=\eps\pair{[R_{\rho-\alpha}p_{n-\rho}]^\alpha R^{\tau-\alpha}}{p_\tau}
 =\eps\pair{R^{\tau-\alpha}q_{n-\tau}}{R_{\rho-\alpha}p_{n-\rho}}.
\end{aligned}
\tag{34}
\end{equation}
Thus an explicit final expression is in fact obtained, and the symmetric form in which $q_\rho$ and $p_\tau$ occur shows that (30.) leads to the same final expression, which is therefore independent of the original decomposition. There is a difference between (25.) and (30.) only as long as the rows $y$ and $v$ have independent meaning, so that the decomposition identity passes, by our definition, into a decomposable identity. To arrive at identities with more symbol rows, one need only, according to \S\ 3 (end), resolve a row $q_\rho$ into $q_{\rho_1}C^{\sigma_1}$, or $p_\tau$ into $p_{\tau_1}C_{\tau-\tau_1}$. The resulting expressions
\begin{equation}
 \pair{q_{\rho_1}C^{\sigma_1}}{[R^{\tau-\alpha}q_{n-\tau}]_\lambda R_{\rho+\sigma_1-\alpha-\lambda}},
 \qquad
 \pair{[R_{\rho-\alpha}p_{n-\rho}]^\lambda R^{\tau-\alpha}}{p_{\tau_1}C_{\tau-\tau_1}}
\tag{35}
\end{equation}
are exactly of the type that occurs for two symbol rows, and therefore again admit explicit representation after a substitution corresponding to (33.). Repetition therefore gives explicit representation for arbitrarily many rows. It may also be noted that the first and last sums in (25.) and (30.), respectively, give the explicit representation without applying (33), solely by formula (32), or else reduce altogether to a single term containing all rows explicitly. The relation obtained from (32.) by resolving $q_\rho$ into $q_{\rho_1}C^{\sigma_1}$ and so forth can be regarded as the product theorem in its most general form.

In summary we shall show:

\emph{Theorem III:} With the identities
\begin{equation}
 \pair{q_\rho A^\sigma}{p_\tau B_{\rho+\sigma-\tau}}
 =\sum_{\alpha=\max(0,\tau-\sigma)}^{\min(\tau,\rho)}
 \eps\pair{R^{\tau-\alpha}q_{n-\tau}}{R_{\rho-\alpha}p_{n-\rho}},
\tag{36}
\end{equation}
where the rows $R$ are determined by (33), together with those obtained from them by resolving $q_\rho,p_\tau$ into $q_{\rho_1}C^{\sigma_1}$, etc., the totality of indecomposable identities is exhausted.

Indeed, the resulting identities are the most general of their kind, since the individual rows are no longer subject to any restriction as to weight or number, as was still the case in (20.) and (21.). Nor do any further identities among these rows exist; for when a row $p_\tau$ is resolved into $y^{(1)}\cdots y^{(\tau)}$, the most general identities are composed out of (20.), hence by \S\ 3 (end) out of (20a.); and the composition discussed following (16.) is precisely the one carried out in \S\ 4. The transition to arbitrarily many rows is supplied, as the discussion of (20a.) shows, by applying always the same operation to the expressions that arise from resolving $q_\rho$ into $q_{\rho_1}C^{\sigma_1}$, and so on. It also follows that the direct transition from (20.) instead of from (20a.) gives nothing new. This is easily verified by calculation, by specializing once a row $q_\rho$ in (25.), and another time carrying out the transition from (20.) by the method of \S\ 4. In both cases the same sums arise, which pass by application of the general product theorem (see above) into the identities obtained from (36.). The identities (20.), (21.) correspond to the special cases $\tau=1$, $\rho=1$; only two individual sums then occur, and hence, in accordance with an earlier remark, explicit representation without the substitution (33.), agreeing with what was found before.

At the end let us briefly mention two special cases of the decomposition identity. Put in (31): $\tau=\sigma$, $\rho=n-\sigma$, and denote the row $p_\tau$ by $p'_\sigma\sim q'_{n-\sigma}$. Then
\begin{equation}
\begin{aligned}
\pair{A^\sigma}{p_\sigma}\pair{B^\sigma}{p'_\sigma}
-\pair{A^\sigma}{p'_\sigma}\pair{B^\sigma}{p_\sigma}
=\sum_{\alpha=1}^{\min(\sigma,n-\sigma)}
\Pi\pair{q_{n-\sigma-\alpha}B^\sigma}{A_{n-\sigma}p_{\sigma-\alpha}}.
\end{aligned}
\tag{37}
\end{equation}
This is Clebsch's formula\footnote{Clebsch, loc. cit., pp. 25--32. Clebsch proves that the individual terms on the right-hand sides depend only on the rows $A,B$; the explicit representation would be obtained by applying (32.) to his expressions $\Phi_h$.} for expanding a form with two cogredient rows $p_\sigma,p'_\sigma$ into elementary covariants. The elementary covariants belonging to a form $\pair{A^\sigma}{p_\sigma}^m\pair{B^\sigma}{p'_\sigma}^n$ are then the following:
\begin{equation}
 \left\{\sum_{\alpha=1}^{\min(\sigma,n-\sigma)}
 \eps\pair{q_{n-\sigma-\alpha}B^\sigma}{A_{n-\sigma}p_{\sigma-\alpha}}
 \right\}^{\rho}
 \pair{A^\sigma}{p_\sigma}^{m-\rho}\pair{B^\sigma}{p'_\sigma}^{n-\rho},
 \qquad (\rho=0,1,\ldots,m,n).
\tag{38}
\end{equation}
They arise, as we shall briefly say, from the ground form by ``cogredient folding of defect $\rho$'' (cf. \S\ 2).

Second, put $\tau=\sigma-\lambda$, $\rho=n-\sigma-\lambda$, $B^\sigma=A^\sigma$. Then
\begin{equation}
\begin{aligned}
\pair{q_{n-\sigma-\lambda}A^\sigma}{A_{n-\sigma}p_{\sigma-\lambda}}
&=(-1)^\lambda\pair{q_{n-\sigma-\lambda}A^\sigma}{A_{n-\sigma}p_{\sigma-\lambda}}\\
&\quad+(-1)^{(n-\sigma)(\sigma-\lambda)}
 \sum_{\alpha=1}^{\min(\sigma-\lambda,n-\sigma-\lambda)}
 \Pi\pair{q_{n-\sigma-\lambda-\alpha}A^\sigma}{A_{n-\sigma}p_{\sigma-\lambda}}.
\end{aligned}
\tag{39}
\end{equation}
Thus, as noted in \S\ 2, the conditions (12.) characterizing the symbol rows for odd defect $\lambda$ are consequences of the conditions of higher defect. For a linear form $\pair{A^\sigma}{p_\sigma}=\pair{B^\sigma}{p_\sigma}$, (39.) implies that its irreducible invariant formations that are quadratic in the coefficients reduce to ones of even defect; this agrees with the known fact that a linear form in the binary and ternary domains has no invariant formations.

It should be noted that the general identities were originally derived under the assumption that the individual rows satisfy conditions (12.). Since, however, these individual rows enter the identities only linearly, the latter also hold for the more general rows that do not satisfy conditions (12.).

\section*{\S\ 6. Explicit representation of the invariant formations.\\Folding and differentiation processes.}
The results of the preceding paragraphs allow us to complete the theorem stated in \S\ 2 on symbolic representability (the first fundamental theorem) by adding explicit representability; namely, we show:

\emph{Theorem IV:} ``All invariant formations of arbitrary ground forms can be represented explicitly by matrix products; that is, aside from variable rows $p_\rho$, only the rows $S^\sigma,\ldots,T^\tau$ of the original forms, or the rows $P,Q,R$ derived from them by substitution (9.) or (11.), enter the final expressions.''

For the proof, suppose that an invariant formation depends, besides variable rows, on the rows $S^\sigma,\ldots,T^\tau$; and suppose that these rows have been resolved sufficiently far into individual rows $S^\sigma=S^{\sigma_1}\cdots S^{\sigma_i},\ldots,T^\tau=T^{\tau_1}\cdots T^{\tau_k}$ to make representation by determinants possible. Let the rows be ordered in a sequence $S^\sigma,\ldots,T^\tau$ that is arbitrary in itself but fixed once and for all. We pick out an aggregate of determinants that depends on the first total row $S^\sigma=S^{\sigma_1}\cdots S^{\sigma_i}$, not merely on individual component rows. This dependence, however, is by \S\ 3--\S\ 5 a consequence of the general identities. If we now resolve a variable row $p_\rho$ into the rows $y,v$, or one of the later rows $T$ into the individual rows $t$, respectively $\tau$, then the most general identities are composed out of (20.) and (21.). But these latter identities always lead to an expression containing the row $S^\sigma=S^{\sigma_1}\cdots S^{\sigma_i}$ explicitly, namely the left-hand side of (20.) and (21.); for by (19.) one verifies that only the complete sum on the right--not a part of it corresponding to a term of (20a.)--has the property of depending only on $S^\sigma$. As the procedure continues, all rows that have once entered explicitly remain explicit; application of the identities can at most require the combining of several rows into a single one, i.e. substitution (9.). If finally only a single row $T^\tau$ remains to be brought to explicit form, two cases are possible. There may still occur a free variable row, that is, one not connected with other rows by substitution (9.); resolving it into component rows $y$ or $v$ completes the procedure and gives, as in (31.), polars of one or more forms containing all rows explicitly. If no free variable row remains, then from the manner in which they arose we recognize, in the totality of the expressions that contain the individual rows $t$ or $\tau$ but depend on the row $T^\tau$, the terms of a decomposition identity; by \S\ 5, however, these always admit explicit representation in all rows after applying substitution (11.). This proves the theorem in full generality. It should also be noted that when only two symbol rows occur, substitution (11.) never appears. For, since $\sigma,n-\sigma,\tau,n-\tau<n$, no variable rows can enter only when the two symbol rows are united in a single determinant and thus occur explicitly; but as soon as variable rows enter, (34.) and (33.) show that substitution (11.) becomes unnecessary.

It follows from what has just been said that, in order to generate all invariant formations, it suffices to unite the symbol rows, after adjoining variable rows, into all possible matrix products. The most general matrix product is now of the form
\begin{equation}
 \pair{P^{\mu_1}\cdots P^{\mu_i}}{Q_{\nu_1}\cdots Q_{\nu_k}},\qquad \sum\mu=\sum\nu,
\tag{40}
\end{equation}
where the $P$ and $Q$ may be rows of the original forms, or may have arisen from the original rows by a sequence of substitutions (9.) or (11.), or finally may be variable rows. But the expressions (40.) can be generated by successively uniting two symbol rows at a time into a matrix product, with the help of variable rows; for from
\[
 \pair{q_{n-\mu-\lambda}P^\mu}{Q_{\nu_1}p_{n-\mu-\nu_1-\lambda}}=\pair{R_\lambda Q_{\nu_1}}{p_{n-\mu-\nu_1-\lambda}}
\]
one obtains, by uniting with the row $Q_{\nu_2}$, an expression
\[
 \pair{[R_\lambda Q_{\nu_1}]^{n-\lambda-\nu_1}}{Q_{\nu_2}p_{n-\mu-\nu_1-\nu_2-\lambda}}
 =\pair{q_{n-\mu-\lambda}P^\mu}{Q_{\nu_1}Q_{\nu_2}p_{n-\mu-\nu_1-\nu_2-\lambda}},
\]
and hence, by continuing the procedure, an expression (40.). If $P$ and $Q$ are not rows of the original forms, then (9.) and (11.), in conjunction with what was just shown, imply that they too arose through a sequence of unions of two symbol rows at a time. Thus:

\emph{Theorem V:} The process for generating all invariant formations, the folding process, consists in successively uniting two symbol rows at a time into a matrix product, with the help of variable rows.

From two symbol rows $S^\sigma,T^\tau$, where $\sigma\ge\tau$, the following matrix products can be formed:
\begin{align}
&\pair{q_{n-\sigma-\lambda}S^\sigma}{T_{n-\tau}p_{\tau-\lambda}}, &&(\lambda=1,2,\ldots,\min(\tau,n-\sigma)),\tag{41}\\
&\pair{q_{n-\tau-\mu}T^\tau}{S_{n-\sigma}p_{\sigma-\mu}}, &&(\mu=\sigma-\tau+\lambda),\tag{42}\\
&\pair{q_{n-\tau-\nu}T^\tau}{S_{n-\sigma}p_{\sigma-\nu}}, &&(\nu=1,2,\ldots,\sigma-\tau).
\tag{43}
\end{align}
From (31.), however, the following values are obtained for (42.) and (43.):
\begin{align}
&\pair{q_{n-\sigma-\lambda}S^\sigma}{T_{n-\tau}p_{\tau-\lambda}}
 +\sum_{\alpha}\Pi\pair{q_{n-\sigma-\lambda-\alpha}S^\sigma}{T_{n-\tau}p_{\tau-\lambda-\alpha}},\tag{42a}\\
&\Pi(q_{n-\sigma}S^\sigma)(T_{n-\tau}p_\tau)
 +\sum_{\beta}\Pi\pair{q_{n-\sigma-\beta}S^\sigma}{T_{n-\tau}p_{\tau-\beta}}.
\tag{43a}
\end{align}
Thus the forms (42.) can be expressed linearly by (41.) and polars of (41.), and the forms (43.) by polars of the ground form and of the forms (41.). Likewise (31.) gives the linear dependence of the forms (41.) on (42.) and polars of (42.). Hence, according to Clebsch's definition,\footnote{Clebsch, loc. cit., \S\ 7.} the forms (42.) are equivalent to the forms (41.), while (43.) leads back to the ground form. The generating process is therefore given with equal right by (41.) or (42.); we shall define the process (41.), which is simpler with respect to the number $\lambda$, as the folding process. The number $\lambda$, the defect of the folding (cf. \S\ 2), gives the number of rows missing from the determinant product, namely in that matrix product which, for the given folding, contains the maximal number of rows. Since $\lambda$ runs only up to the smaller of the two values $\tau,n-\sigma$, where $\sigma\ge\tau$, we have
\begin{equation}
 \lambda\le\frac n2,\ n\text{ even};\qquad
 \lambda\le\frac{n-1}{2},\ n\text{ odd}.
\tag{44}
\end{equation}
The reduction of (42.) and (43.) to (41.) remains valid even when, through further folding, the rows $p,q$ have passed into symbol rows; for by (36.) one has
\[
 \pair{A^{n-\tau-\varkappa}T^\tau}{S_{n-\sigma}B_{\sigma-\varkappa}}
 =\sum_{\alpha=\max(0,\varkappa-\sigma+\tau)}^{\min(\tau,n-\sigma)}
 \eps\pair{R^{\tau-\alpha}B^{n-\sigma+\varkappa}}{A_{\tau+\varkappa}R_{n-\sigma-\alpha}},
\]
where the rows $R$ have arisen from $S$ and $T$ according to (33.). Thus these forms containing symbol rows only are obtained by folding $\pair{q_\tau A^{n-\tau-\varkappa}}{B_{\sigma-\varkappa}p_{n-\sigma}}$ or $\pair{q_{\tau-\alpha}B^{n-\sigma+\varkappa}}{A_{\tau+\varkappa}p_{n-\sigma-\alpha}}$--according as $\sigma-\tau\gtreqless 2\varkappa$--with the ground form and the forms (41.).

From the symbolic folding process one obtains, in the usual way, the invariant differentiation processes if one merely replaces the symbol rows $S^\sigma,T_{n-\tau}$ by the rows of differentiation symbols $\frac{\partial}{\partial p_\sigma}$, $\frac{\partial}{\partial q_{n-\tau}}$; as is well known, the product $\frac{\partial}{\partial p_{i_1\ldots i_\sigma}}\frac{\partial}{\partial q_{j_1\ldots j_{n-\tau}}}$ then has the real meaning $\frac{\partial^2}{\partial p_{i_1\ldots i_\sigma}\partial q_{j_1\ldots j_{n-\tau}}}$. Since the rows $\frac{\partial}{\partial p_\sigma}$ and $\frac{\partial}{\partial q_{n-\tau}}$ are cogredient with the rows $S^\sigma,T_{n-\tau}$, they satisfy the same identities as these, in particular also the ones existing between (42.) and (42a.), and between (43.) and (43a.). In summary we obtain:

\emph{Theorem VI:} All invariant formations of arbitrary ground forms arise from these by repeated application of the symbolic folding process (41.), or also of the invariant differentiation processes
\begin{equation}
 \left(q_{n-\sigma-\lambda}\frac{\partial}{\partial p_\sigma}\middle|\frac{\partial}{\partial q_{n-\tau}}p_{\tau-\lambda}\right),
 \qquad(\sigma\ge\tau,\ \lambda=1,2,\ldots,\min(\tau,n-\sigma)).
\tag{45}
\end{equation}
It remains to note that in every folding (41.) the total weight of the form, that is, the sum of the weights of the individual variable rows (cf. \S\ 1), decreases by the positive number $2(\lambda^2+\lambda(\sigma-\tau))$. This implies, independently of the general finiteness proof, the finiteness of the form row, that is, of the totality of invariant formations of a given ground form that are linear in the coefficients.\footnote{Cf. Clebsch, loc. cit., \S\ 11 and 12: finiteness of the system of elementary covariants.}

It should further be noted that Theorem VI still holds for forms that do not satisfy conditions (12.). For each factor in (3), $\pair{S^\rho}{p_\rho}^{m_\rho}$, can be replaced by a product of $m_\rho$ linear factors $\pair{A^\rho}{p_\rho}\pair{B^\rho}{p_\rho}\cdots\pair{M^\rho}{p_\rho}$; the invariant formations thereby pass into those of a system of linear forms, which need only afterwards be set equal to one another. For every single linear form, however, conditions (12.)--which, when differential symbols are introduced, contain only second-order differential quotients--are always satisfied.

\section*{\S\ 7. Expansion of the general forms according to normal forms.}
In what follows (\S\ 8) we shall derive a principal property of the folding process (41.), namely its composability from fundamental foldings. For this we need the transfer to the $n$-ary domain of an important theorem given by Mertens\footnote{F. Mertens, \emph{\"Uber invariante Gebilde quatern\"arer Formen} (see the Introduction), cited as ``Mertens.''} in the quaternary domain. This theorem says that any finite system of forms can be replaced by an equivalent system of normal forms. By a ``normal form'' we mean--generalizing the binary and ternary definition\footnote{Study, loc. cit., p. 54.}--a form all of whose invariant formations linear in the coefficients vanish identically, with the exception of the ground form itself. Thus by Theorem VI (\S\ 6) a normal form $\varphi$ satisfies the system of differential equations
\begin{equation}
 \left(q_{n-\sigma-\lambda}\frac{\partial}{\partial p_\sigma}\middle|\frac{\partial}{\partial q_{n-\tau}}p_{\tau-\lambda}\right)\varphi=0,
 \qquad(\sigma\ge\tau,\ \lambda=1,2,\ldots,\min(\tau,n-\sigma)),
\tag{46}
\end{equation}
where the rows $q_{n-\sigma-\lambda},p_{\tau-\lambda}$ are to be regarded as arbitrary quantities. In the quaternary domain--taking (39.) into account for $\sigma=\tau=2$--these differential equations agree completely with the ten differential equations given by Mertens without introduction of the arbitrary rows $p,q$ (loc. cit., formula 28). Their knowledge, together with Theorem VI, allows us to transfer his proof almost word for word to $n$ variables. The only thing to add is the verification, obtained by the decomposition identity (30.), that the constructed forms are normal forms; in the quaternary domain this verification still follows without further transformation. It is therefore enough to give a short exposition of the proof, and for simplicity in the case relevant below, namely that of a single ground form $f$ and its form row (cf. \S\ 6, end); for invariant formations of arbitrary degree in the coefficients of arbitrarily many ground forms, the proof remains the same.

The basic idea of the proof is to replace, by introducing the transformed coefficients, a form with $n-1$ rows $p_\rho$ by forms with $n$ rows $\xi$ cogredient to the $x$, namely the rows of the transformation quantities. From these forms the sought normal forms are obtained directly by invariant differentiation processes. The expansion of the general forms according to normal forms is mediated by a series expansion for forms with $\rho$ ($\rho<n$) cogredient rows $\xi$; invariant differentiation processes lead back to the forms in the original rows $p_\rho$.

Let $\xi^{(1)},\ldots,\xi^{(n)}$ be the rows of transformation quantities, so that
\begin{equation}
\begin{aligned}
x_i&=\xi_i^{(1)}x'_1+\xi_i^{(2)}x'_2+\cdots+\xi_i^{(n)}x'_n,\\
p_{i_1\ldots i_\rho}&=\sum_j(\xi_{i_1}^{(j_1)}\cdots \xi_{i_\rho}^{(j_\rho)})p'_{j_1\ldots j_\rho}.
\end{aligned}
\tag{47}
\end{equation}
The transformed coefficients of the forms $f,\varphi,\ldots$ will be denoted by $(f),(\varphi),\ldots$, and let
\[
 f=\pair{S^1}{p_1}^{m_1}\pair{S^2}{p_2}^{m_2}\cdots\pair{S^{n-1}}{p_{n-1}}^{m_{n-1}}.
\]
Then
\begin{equation}
\begin{aligned}
(f)&=\pair{S^1}{\xi^{(1)}}^{m_{11}}\cdots\pair{S^1}{\xi^{(n)}}^{m_{1n}}
\pair{S^2}{\xi^{(1)}\xi^{(2)}}^{m_{21}}\cdots
\pair{S^{n-1}}{\xi^{(2)}\cdots\xi^{(n)}}^{m_{n-1,n}}\\
&\quad\cdot\pair{s^{(1)}}{\xi^{(1)}}^{k_1}\pair{s^{(2)}}{\xi^{(2)}}^{k_2}\cdots\pair{s^{(n)}}{\xi^{(n)}}^{k_n},
\end{aligned}
\tag{48}
\end{equation}
where $m_{11}+\cdots+m_{1n}=m_1$, and so on. From each coefficient $(f)$ for which
\begin{equation}
 k_1\ge k_2\ge k_3\ge\cdots\ge k_n,
\tag{49}
\end{equation}
the differentiation process exhausting exactly the rows $\xi$,
\begin{equation}
\left(\frac{\partial}{\partial\xi^{(1)}}\middle|p_1\right)^{k_1-k_2}
\left(\frac{\partial}{\partial\xi^{(1)}}\frac{\partial}{\partial\xi^{(2)}}\middle|p_2\right)^{k_2-k_3}
\cdots
\left(\frac{\partial}{\partial\xi^{(1)}}\cdots\frac{\partial}{\partial\xi^{(n-1)}}\middle|p_{n-1}\right)^{k_{n-1}-k_n}
\left(\frac{\partial}{\partial\xi^{(1)}}\cdots\frac{\partial}{\partial\xi^{(n)}}\right)^{k_n}
\tag{50}
\end{equation}
gives the form
\begin{equation}
 \varphi=\bigl(s^{(1)}\mid p_1\bigr)^{k_1-k_2}
 \bigl(s^{(1)}s^{(2)}\mid p_2\bigr)^{k_2-k_3}\cdots
 \bigl(s^{(1)}\cdots s^{(n-1)}\mid p_{n-1}\bigr)^{k_{n-1}-k_n}
 \bigl(s^{(1)}s^{(2)}\cdots s^{(n)}\bigr)^{k_n}.
\tag{51}
\end{equation}
The forms $\varphi$ are the sought normal forms. Indeed, they have the following properties defining normal forms.

1. They are invariant formations of the ground form $f$ that are linear in the coefficients. This follows from the fact that they arise from $f$ by the invariant differentiation processes $\left(\frac{\partial}{\partial p_\rho}\middle|\xi^{(i_1)}\xi^{(i_2)}\cdots\xi^{(i_\rho)}\right)$ and (50.). Thus (by \S\ 6, end) they can be expressed linearly by polars of the finite form row of $f$. These polars are obtained by regarding the variables $p_\rho$ entering $\varphi$ as coefficients of a form decomposed into linear forms with the variable rows $p_\rho$:
\begin{equation}
 H=(q_{n-1}\mid p_{n-1})^{k_1-k_2}(q_{n-2}\mid p_{n-2})^{k_2-k_3}\cdots(q_1\mid p_1)^{k_{n-1}-k_n}.
\tag{52}
\end{equation}
The simultaneous invariants of the form rows $f$ and $H$ give all relevant polars, since the latter must be of the same dimension in the individual rows $p_\rho$ as $\varphi$ itself.

2. They satisfy the differential equations (46.). If the differential process (45.) is applied to $\varphi$, then
\[
 \varphi'=\bigl(q_{n-\sigma-\lambda}s^{(1)}\cdots s^{(\sigma)}\mid [s^{(1)}\cdots s^{(\tau)}]_{n-\tau}p_{\tau-\lambda}\bigr),\qquad (\sigma>\tau).
\]
From the decomposition identity (30.) it follows, if the rows $q_\sigma=s^{(1)}\cdots s^{(\sigma)}$, $p_{n-\tau}=[s^{(1)}\cdots s^{(\tau)}]_{n-\tau}$ formed out of the rows $s$ are identified with the $q_\rho,p_\tau$ occurring there, that
\[
 \varphi'=\sum_{\beta=\sigma-\tau+\lambda}^{n-\tau,\sigma}\sum_{i_\beta}\eps\,
 \pair{q_{\sigma-\beta}^{(i_1\ldots i_\beta)}q_{n-\tau+\lambda}}{p_{\sigma+\lambda}p_{n-\tau-\beta}^{(i_1\ldots i_\beta)}}.
\]
Here
\[
 p_{n-\tau-\beta}^{(i_1\ldots i_\beta)}\sim s^{(1)}\cdots s^{(\tau)}s^{(i_1)}\cdots s^{(i_\beta)},
\]
where $i_1,
\ldots,i_\beta$ denote any $\beta$ distinct numbers chosen from the numbers 1 to $\sigma$. Since $\beta>\sigma-\tau$, it follows that among these numbers $i_1,
\ldots,i_\beta$ there must be some between 1 and $\tau$; hence $p_{n-\tau-\beta}^{(i_1\ldots i_\beta)}$ vanishes identically, and therefore every individual matrix product, and with it $\varphi'$, vanishes.

For the equivalence of the system of normal forms with the form row $f$, it remains only to show that all forms of the form row can be expanded in polars of the normal forms in the sense fixed under 1. Let $\Pi$ be an aggregate of polar operations
\begin{equation}
 \left(\frac{\partial}{\partial\xi^{(i)}}\middle|\xi^{(k)}\right),
 \qquad (i\ne k;\ i=1,2,\ldots,\rho;\ k=1,2,\ldots,\rho),
\tag{53}
\end{equation}
and let $\Pi^\sigma$ be an aggregate of the special operations (53.) for which $i=\sigma$ and $k<\sigma$. For a form $F$ of the $\rho$ rows $\xi^{(1)},\ldots,\xi^{(\rho)}$ that has degree $k_\rho$ in the row $\xi^{(\rho)}$, the expansion\footnote{F. Mertens, \emph{\"Uber eine Formel der Determinantentheorie}. Sitzungsberichte der Akademie der Wissenschaften, Vienna, Math.-Naturw. Kl., vol. 91, 2nd section (1885), formula 20).} holds:
\begin{equation}
 F=\sum_{\alpha=0}^{k_\rho}\Pi F_\alpha,
 \qquad(\rho\le n),
\tag{54}
\end{equation}
where $F_\alpha$ has degree $\alpha$ in the last row $\xi^{(\rho)}$ and arises from $F$ by the invariant differential operations
\begin{equation}
 \left(\frac{\partial}{\partial\xi^{(1)}}\cdots\frac{\partial}{\partial\xi^{(\rho)}}\middle|\xi^{(1)}\cdots\xi^{(\rho)}\right)^\alpha
\tag{55}
\end{equation}
and $\Pi^\rho$. If, in constructing $F_\alpha$, one replaces process (55.) by
\begin{equation}
 \left(\frac{\partial}{\partial\xi^{(1)}}\cdots\frac{\partial}{\partial\xi^{(\rho)}}\middle|p_\rho\right)^\alpha,
\tag{56}
\end{equation}
then a form $F_\alpha(p_\rho)$ with only $\rho-1$ rows $\xi^{(1)},\ldots,\xi^{(\rho-1)}$ is obtained, to which (54.) can again be applied. Continuing the procedure leads to forms $G(p_\rho,p_{\rho-1},\ldots,p_1)$. In order to obtain from these the expansion of $F$ according to (54.), one replaces the individual rows $p_\rho$ by $\xi^{(1)},\ldots,\xi^{(\rho)}$ and applies to the resulting forms the polar process $\Pi$ (53.); this gives (cf. (48.)) a sum of transformed coefficients $(G)$. For $\rho=n$, the process (55.) remains at the first step. If $c$ denotes numerical constants and we put, for brevity,
\begin{equation}
 (\xi^{(1)}\xi^{(2)}\cdots\xi^{(n)})\sim\Delta;
 \qquad
 \left(\frac{\partial}{\partial\xi^{(1)}}\frac{\partial}{\partial\xi^{(2)}}\cdots\frac{\partial}{\partial\xi^{(n)}}\right)=\nabla,
\tag{57}
\end{equation}
then, in the case $\rho=n$,
\begin{equation}
 F=\sum c\Delta^\alpha(G),
\tag{58}
\end{equation}
whereas for $\rho<n$ only $\alpha=0$, and hence $\sum c(G)$, appears on the right.

Applying (58.) to a transformed coefficient $(f)$ (48.), and using the commutativity of the polar processes $\Pi^\sigma$ with all operations (56.) for which $\rho\ge\sigma$, as well as the fact that polars of transformed coefficients again give transformed coefficients, one obtains
\[
 G=\left(\frac{\partial}{\partial\xi^{(1)}}\middle|p_1\right)^{\beta_1}
 \left(\frac{\partial}{\partial\xi^{(1)}}\frac{\partial}{\partial\xi^{(2)}}\middle|p_2\right)^{\beta_2}
 \cdots
 \left(\frac{\partial}{\partial\xi^{(1)}}\cdots\frac{\partial}{\partial\xi^{(n-1)}}\middle|p_{n-1}\right)^{\beta_{n-1}}
 \nabla^{\beta_n}\Pi(f)=\sum c\varphi,
\]
and from (58.) it follows that
\begin{equation}
 (f)=\sum c\Delta^\alpha(\varphi)\qquad\text{(Mertens, formula 23)).}
\tag{59}
\end{equation}
To pass from (59.) to the expansion of $f$ itself, we again regard the variables $p_\rho$ as coefficients of a form $H$ (52.) with the degrees $m_1,
\ldots,m_{n-1}$ corresponding to $f$, and put $m=m_1+\cdots+m_{n-1}$. Then, by the definition of invariants,
\[
 f\cdot\Delta^m=\sum(f)(H)=\sum c\Delta^\alpha(\varphi)(H),
\]
and application of the process $\nabla^m$ that exactly exhausts the rows $\xi$ gives
\begin{equation}
 f=\sum c\Phi,
\tag{60}
\end{equation}
where the forms $\Phi$ are derived from $\varphi$ and $H$ by invariant differentiation processes with respect to the variables, and hence are simultaneous invariants of the form rows $\varphi$ and $H$.\footnote{Instead of the direct conclusion, Mertens, \S\ 7, introduces a series of further differential processes that essentially serve to form the form row $H$; this is done by our Theorem VI.} But since $\varphi$ is a normal form, the form row $\varphi$ reduces to the form $\varphi$ itself; the form row $H$ contains all possible invariant variable combinations. Thus the simultaneous invariants of $\varphi$ and $H$ are the polars of $\varphi$ in the sense indicated under 1.; in other words, (60.) gives the expansion of $f$ in polars of the normal forms. It remains to verify the same expansion for an arbitrary form $\theta$ of the form row $f$. Let $\Gamma$ be the normal forms derived from $(\theta)$ by (50.), so that
\begin{equation}
 (\theta)=\sum c\Delta^\beta(\Gamma).
\tag{61}
\end{equation}
The forms $\theta$ are characterized by the relation, valid for every transformed coefficient,
\[
 (\theta)\Delta^\sigma=\sum c(f)\qquad\text{(Mertens, formula 9)).}
\]
But every form $F$ of the $n$ rows $\xi$ which contains the factor $\Delta^\sigma$ also contains the second $\nabla^\sigma\Pi F$ (Mertens, formula 20)); hence
\[
 (\theta)=\sum c\nabla^\sigma(f),\qquad\text{therefore}\qquad \Gamma=\sum c\varphi,
\]
and from (61.) we obtain the relation corresponding to (59.),
\[
 (\theta)=\sum c\Delta^\beta(\varphi),
\]
which implies for $\theta$ the expansion (60.) in polars of the normal forms $\varphi$. Thus:

\emph{Theorem VII.} Every form row $f$ can be replaced by an equivalent system of normal forms.

\section*{\S\ 8. Reduction of the folding process to fundamental foldings.}
Following Theorem VII, we now carry out the reduction of the folding process to fundamental foldings mentioned at the beginning of \S\ 7. We call (cf. \S\ 5) every folding of two cogredient rows $S^\rho,T^\rho$ a ``cogredient folding'', and specifically the cogredient folding of defect 1 of two rows $S^\rho,T^\rho$ a ``fundamental folding of type $\rho$''. We then prove, completing Theorem VI:

\emph{Theorem VIII:} The general folding process (41.) can be composed from the $(n-1)$ fundamental foldings of type $\rho$, where $\rho=1,2,\ldots,n-1$. In other words, to generate all invariant formations it suffices to apply successively the $n-1$ fundamental foldings.

In the binary domain the folding process (41.) coincides directly with the only possible fundamental folding; in the ternary domain I proved Theorem VIII in \S\ 1 of my dissertation (cf. the Introduction). There, because of (44.), the general cogredient foldings are still identical with the fundamental foldings, and it is noteworthy that Theorem VIII for $n$ variables gives the reduction to fundamental foldings, not merely the much weaker reduction to cogredient foldings. The proof is modeled in principle on the ternary proof: the most general foldings are constructed from the fundamental foldings, using the symbolic identities. We shall generate:
\begin{enumerate}
\item arbitrary foldings of defect 1 by cogredient foldings of defect 1, that is, by fundamental foldings (the analogue of the ternary case),
\item arbitrary foldings of defect $\lambda$ by cogredient foldings of defect $\lambda$ and arbitrary foldings of defect 1,
\item cogredient foldings of defect $\lambda$ by cogredient foldings of defect $\lambda-1$ and arbitrary foldings of defect 1.
\end{enumerate}
As is essential, we assume by Theorem VII that the two forms $S$ and $T$ to be folded are normal forms. Thus, by (46.), the relations hold:
\begin{equation}
\begin{aligned}
\pair{q_{n-\sigma_1-\lambda}S^{\sigma_1}}{S_{n-\sigma_2}p_{\sigma_2-\lambda}}&=0,
&& (\sigma_1\ge\sigma_2,\ \lambda=1,2,\ldots,\sigma_2,n-\sigma_1),\\
\pair{q_{n-\tau_1-\lambda}T^{\tau_1}}{T_{n-\tau_2}p_{\tau_2-\lambda}}&=0,
&& (\tau_1\ge\tau_2,\ \lambda=1,2,\ldots,\tau_2,n-\tau_1).
\end{aligned}
\tag{62}
\end{equation}
It is further sufficient, by \S\ 2 (end), to suppose that the forms $S$ and $T$ are linear in all variable rows; that is, the constructed forms belong to the form row
\[
 f=\pair{S^1}{p_1}\pair{S^2}{p_2}\cdots\pair{S^{n-1}}{p_{n-1}}
   \pair{T^1}{p_1}\pair{T^2}{p_2}\cdots\pair{T^{n-1}}{p_{n-1}}.
\]
1) Generation of $\pair{q_{n-\sigma-1}S^\sigma}{T_{n-\tau}p_{\tau-1}}$, where $\sigma>\tau$, from fundamental foldings of type $\tau+\alpha$, $\alpha=0,1,\ldots,\sigma-\tau$. From the fundamental folding of type $\tau$
\[
 \pair{q_{n-\tau-1}S^\tau}{T_{n-\tau}p_{\tau-1}}
\]
we obtain, by the substitution $w\sim T_{n-\tau}p_{\tau-1}$, a form of the form row $f$ with three rows of upper weight index (cf. \S\ 1) $\tau+1$:
\begin{equation}
 \pair{wS^\tau}{p_{\tau+1}}\pair{S^{\tau+1}}{p_{\tau+1}}\pair{T^{\tau+1}}{p_{\tau+1}}.
\tag{63}
\end{equation}
A fundamental folding of type $\tau+1$ applied to (63.) gives
\[
 \pair{q_{n-\tau-2}wS^\tau}{S_{n-\tau-1}p_\tau};
\]
from (31.), by the substitution $q_{n-\tau-2}w=q'_{n-\tau-1}$, this has the value
\[
 \eps\pair{q_{n-\tau-2}wS^{\tau+1}}{S^\tau p_\tau}
 +\sum_{\alpha=1}^{\tau,n-\tau-1}\Pi\pair{q'_{n-\tau-1-\alpha}S^{\tau+1}}{S_{n-\tau}p_{\tau-\alpha}}.
\]
The expression under the summation sign vanishes identically by (62.); thus we have reached a form
\[
 \pair{wS^{\tau+1}}{p_{\tau+2}}\pair{S^{\tau+2}}{p_{\tau+2}}\pair{T^{\tau+2}}{p_{\tau+2}},
\]
which is obtained from (63.) by changing $\tau$ into $\tau+1$. Repeating the process $\sigma-\tau$ times therefore leads to the desired folding:
\[
 \pair{wS^\sigma}{p_{\sigma+1}}=\eps\pair{q_{n-\sigma-1}S^\sigma}{T_{n-\tau}p_{\tau-1}}.
\]
We indicate the process just carried out by the formula
\begin{equation}
 T^\tau S^\tau,
\quad S^{\tau+1},\quad S^{\tau+2},\ldots,S^\sigma;
\tag{64}
\end{equation}
and see that only the normal-form property of $S$, not that of $T$, is used. A process dual to (64.) can be performed, using only the normal-form property of $T$, in the reverse order, namely:
\begin{equation}
 S^\sigma T^\sigma,
\quad T^{\sigma-1},\quad T^{\sigma-2},\ldots,T^\tau,
\tag{65}
\end{equation}
according to the relations:
\begin{align*}
\pair{q_{n-\sigma-1}S^\sigma}{T_{1-\sigma}p_{\sigma-1}}&=\eps\pair{T_{n-\sigma}y}{p_{\sigma-1}},\quad y\sim q_{n-\sigma-1}S^\sigma,\\
\pair{q_{n-\sigma}T^{\sigma-1}}{T_{n-\sigma}yp_{\sigma-2}}&=\eps\pair{T_{n-\sigma+1}y}{p_{\sigma-2}},\\
&\vdots\\
\pair{q_{n-\tau-1}T^\tau}{T_{n-\tau-1}yp_{\tau-1}}&=\eps\pair{q_{n-\sigma-1}S^\sigma}{T_{n-\tau}p_{\tau-1}}.
\end{align*}
We should also point out the connection with the decrease of total weight given in \S\ 6 (end): for foldings of defect 1 this decrease has the value $2(1+(\sigma-\tau))$, and for fundamental foldings the value $2\cdot1$. Therefore, if composition out of fundamental foldings is possible at all, exactly $\sigma-\tau+1$ such foldings are necessarily required, as was carried out above.

2) Generation of general foldings from cogredient and fundamental foldings. The weight decrease for a general folding (41) is $2(\lambda^2+\lambda(\sigma-\tau))=2[\lambda^2+(\sigma-\tau)(1+(\lambda-1))]$. This gives the possibility of decomposing it into a cogredient folding of defect $\lambda$ (weight decrease $2\lambda^2$) and $\sigma-\tau$ foldings of defect 1 with difference of weight indices $\lambda-1$ (weight decrease $2\{1+(\lambda-1)\}$). In the case $\lambda=1$ this decomposition agrees with the one given under 1); we show that it can indeed be realized.

From the cogredient folding of defect $\lambda$
\[
 \pair{q_{n-\tau-\lambda}S^\tau}{T_{n-\tau}p_{\tau-\lambda}}
\]
we obtain, by the substitution $R^\lambda\sim T_{n-\tau}p_{\tau-\lambda}$, a form with the two rows of upper weight indices $\tau+\lambda$ and $\tau+1$:
\begin{equation}
 \pair{R^\lambda S^\tau}{p_{\tau+\lambda}}\pair{S^{\tau+1}}{p_{\tau+1}}.
\tag{66}
\end{equation}
The row with the lower weight index $\tau+1$ belongs to a normal form; hence (66.) admits a folding of defect 1 composed out of $\lambda$ fundamental foldings, in the order (65):
\[
 (R^\lambda S^\tau)S^{\tau+\lambda},\quad S^{\tau+\lambda-1},\quad S^{\tau+\lambda-2},\ldots,S^{\tau+1}.
\]
Taking account of (31.) and (62.), this gives
\[
 \pair{q_{n-\tau-\lambda-1}R^\lambda S^\tau}{S_{n-\tau-1}p_\tau}=\eps\pair{R^\lambda S^{\tau+1}}{p_{\tau+\lambda+1}},
\]
that is, a form obtained from (66.) by replacing $\tau$ by $\tau+1$, just as in 1) for (63.). Repeating the process $\sigma-\tau$ times leads to the desired folding:
\[
 \pair{R^\lambda S^\sigma}{p_{\sigma+\lambda}}=\eps\pair{q_{n-\sigma-\lambda}S^\sigma}{T_{n-\tau}p_{\tau-\lambda}}.
\]
In the whole process only the normal-form property of $S$ was used. Viewing $T$ as a normal form gives the dual process, with complete reversal of the order, according to the formulas
\begin{align*}
\pair{q_{n-\sigma-\lambda}S^\sigma}{T_{n-\sigma}p_{\sigma-\lambda}}&=\eps\pair{T_{n-\sigma}R_\lambda}{p_{\sigma-\lambda}},\quad R_\lambda\sim q_{n-\sigma-\lambda}S^\sigma,\\
\pair{q_{n-\sigma}T^{\sigma-1}}{T_{n-\sigma}R_\lambda p_{\sigma-\lambda-1}}&=\eps\pair{T_{n-\sigma+1}R_\lambda}{p_{\sigma-\lambda-1}},\\
&\vdots\\
\pair{q_{n-\tau-1}T^\tau}{T_{n-\tau-1}R_\lambda p_{\tau-\lambda}}&=\eps\pair{q_{n-\sigma-\lambda}S^\sigma}{T_{n-\tau}p_{\tau-\lambda}}.
\end{align*}
3) Generation of cogredient foldings of defect $\lambda$ from cogredient foldings of defect $\lambda-1$ and fundamental foldings. The weight decrease for cogredient foldings of defect $\lambda$ is $2\lambda^2=2((\lambda-1)^2+2\lambda-1)$, giving the possibility of a decomposition into a cogredient folding of defect $\lambda-1$ and $2\lambda-1$ fundamental foldings. This is most easily found by first putting $n=2\lambda$. The only possible cogredient folding of defect $\lambda$ is $\pair{S^\lambda}{T_\lambda}=\pair{S^\lambda}{T_{n/2}}$; it contains one row $u$ and one row $x$ fewer than the folding of defect $\lambda-1$ of the same rows, $\pair{uS^\lambda}{T_\lambda x}$. Conversely, we compose the folding of defect $\lambda$ out of fundamental foldings which effect the disappearance of a row $u$ and a row $x$ (weight decrease $2(n-1)=2(2\lambda-1)$) and at the same time generate a new symbol row $R^\lambda$, together with a cogredient folding of defect $\lambda-1$. The simplest folding that does this is the following one of defect 1:
\[
 \pair{sT^\lambda}{\sigma p_\lambda}=\eps\pair{S^{2\lambda-1}}{R_{\lambda-1}p_\lambda}=\eps\pair{R^\lambda}{p_\lambda},
\]
where
\[
 R^{\lambda+1}=sT^\lambda,
 \quad s=S^1,
 \quad \sigma=S_1,
 \quad R^\lambda\sim\sigma R_{\lambda-1}.
\]
By (65.) and (64.) it is composed out of the $2\lambda-1$ fundamental foldings
\begin{equation}
 T^\lambda S^\lambda,S^{\lambda-1},\ldots,S^1;
 \qquad R^{\lambda+1}S^{\lambda+1},S^{\lambda+2},\ldots,S^{2\lambda-1}.
\tag{67}
\end{equation}
A folding of defect $\lambda-1$, taking account of (21.) and (62.), gives
\[
 \pair{uS^\lambda}{\sigma R_{\lambda-1}x}=\eps\pair{R^{\lambda+1}}{S_\lambda x}
 =\eps\pair{sT^\lambda}{S^\lambda x}=\eps\pair{S^\lambda}{T_\lambda}.
\]
The general case, for two rows $S^\rho,T^\rho$ and arbitrary $n$, is obtained directly from the special one. In place of (67.) one has the following $2\lambda-1$ fundamental foldings:
\begin{equation}
 T^\rho S^\rho,S^{\rho-1},\ldots,S^{\rho-\lambda+1};
 \qquad R^{\rho+1}S^{\rho+1},S^{\rho+2},\ldots,S^{\rho+\lambda-1}.
\tag{68}
\end{equation}
The first half of the foldings gives the form
\[
 \pair{q_{n-\rho-1}T^\rho}{S_{n-\rho+\lambda-1}p_{\rho-\lambda}},
\]
from which, by the substitutions
\[
 S_{n-\rho+\lambda-1}p_{\rho-\lambda}\sim w,
 \qquad T^\rho w=R^{\rho+1},
\]
and application of the second half of the foldings, one obtains
\[
 \pair{q_{n-\rho-\lambda}S^{\rho+\lambda-1}}{R_{n-\rho-1}p_\rho}.
\]
Substituting
\[
 q_{n-\rho-\lambda}S^{\rho+\lambda-1}\sim y,
\]
one obtains, by a cogredient folding of defect $\lambda-1$,
\begin{equation}
 \pair{q_{n-\rho-\lambda+1}S^\rho}{yR_{n-\rho-1}p_{\rho-\lambda+1}}.
\tag{69}
\end{equation}
We show that this is the desired folding of defect $\lambda$. Substitute
\[
 R_{n-\rho-1}p_{\rho-\lambda+1}=R_{n-\lambda}
\]
and observe that resolving $y$ gives
\[
 \pair{q_{n-\rho-\lambda+1}R^\lambda}{S_{n-\rho}y}
 =\eps\pair{q_{n-\rho-\lambda}S^{\rho+\lambda-1}}{S_{n-\rho}p'_{\rho-1}}=0.
\]
Thus by (20.) the value of (69.) is
\[
 \eps\pair{S^\rho R^\lambda}{p_{\rho+\lambda-1}y}
 =\eps\pair{q_{n-\rho-\lambda}S^{\rho+\lambda-1}}{[S^\rho R^\lambda]_{n-\rho-\lambda}p_{\rho+\lambda-1}}.
\]
This folding is of type (43.), hence equivalent to
\[
 \pair{q_{n-\rho-\lambda}S^\rho}{R_{n-\rho-1}p_{\rho-\lambda+1}}
 =\eps\pair{wT^\rho}{R_\lambda p_{\rho-\lambda+1}},
 \qquad R_\lambda\sim q_{n-\rho-\lambda}S^\rho.
\]
Here $R_\lambda$ is already of the desired final form; the expression corresponds to the form $\pair{sT^\lambda}{S_\lambda x}$ in the special case. Now, observing that resolving $w$ and $R_\lambda$ gives
\[
 \pair{wR^{n-\lambda}}{T_{n-\rho}p_{\rho-\lambda+1}}
 =\eps\pair{q'_{n-1}R^{n-\lambda}}{S_{n-\rho+\lambda-1}p_{\rho-\lambda}}
 =\eps\pair{q''_{n-\sigma-1}S^\rho}{S_{n-\rho+\lambda-1}p_{\rho-\lambda}}=0,
\]
we get by (21.) the value
\[
 \pair{wq_{n-\rho+\lambda-1}}{T_{n-\rho}R_\lambda}
 =\eps\pair{q_{n-\rho+\lambda-1}[T_{n-\rho}R_\lambda]^{n-\lambda}}{S_{n-\rho+\lambda-1}p_{\rho-\lambda}}
 \quad\text{equivalent to}\quad
 \pair{q_{n-\rho-\lambda}S^\rho}{T_{n-\rho}p_{\rho-\lambda}}.
\]
In the whole process only the normal-form property of $S$, not that of $T$, was used. Thus the cogredient folding of defect $\lambda-1$ used to generate (69.) can be replaced by $2\lambda-3$ fundamental foldings and a cogredient folding of defect $\lambda-2$, which in turn admits a decomposition into fundamental foldings, and so on. Analogously, the generation is obtained by using the normal-form property of $T$ and applying the dualistic process, that is, the fundamental foldings
\[
 S^\rho T^\rho,\quad T^{\rho-1},\ldots,T^{\rho-\lambda+1};
 \qquad R^{\rho-1}T^{\rho-1},\quad T^{\rho-2},\ldots,T^{\rho-\lambda+1},
\]
and a folding dual to (69.). The folding
\[
\pair{q_{n-\rho-\lambda}T^\rho}{S_{n-\rho}p_{\rho-\lambda}},
\]
equivalent to the preceding one, is then obtained.

Adding the decomposition obtained under 2), we have indeed obtained a generation of the most general foldings from fundamental foldings. It remains to show that this generation is unique, apart from permutation of the order, i.e. that to each folding there corresponds one and only one number $\alpha_\rho$ of fundamental foldings of type $\rho$, where $\rho=1,2,\ldots,n-1$.

Let $m_\rho$ be the degrees in the variable rows $p_\rho$ of the initial form $f$, and $l_\rho$ the degrees of the form produced by folding. Each fundamental folding of type $\rho$ transforms the degrees $m_{\rho-1},m_\rho,m_{\rho+1}$ respectively into $m_{\rho-1}+1,m_\rho-2,m_{\rho+1}+1$. Thus the $\alpha$ satisfy the system of equations
\begin{equation}
 m_\rho-l_\rho=-\alpha_{\rho-1}+2\alpha_\rho-\alpha_{\rho+1},
 \qquad (\rho=1,2,\ldots,n-1),
\tag{70}
\end{equation}
where $\alpha_0$ and $\alpha_n=0$ are to be put, and whose determinant has value $n$.\footnote{If $\Delta_n$ denotes the $n$-rowed determinant, then the recurrence formula holds: $\Delta_n=2\Delta_{n-1}-\Delta_{n-2}$; that is, $\Delta_n-\Delta_{n-1}=\Delta_{n-1}-\Delta_{n-2}=\cdots=\Delta_2-\Delta_1=1$, since $\Delta_2=3$, $\Delta_1=2$.} The $\alpha$ are thereby uniquely determined, and indeed, by Theorem VIII, as positive integers; this gives conditions for the differences $m_\rho-l_\rho$.

The results of this paragraph hold by virtue of relations (62.); however, if the relations (62.) do not hold, the final expressions are by no means equivalent to the initial expressions. The definition of equivalence given in \S\ 6 also admits the formulation:\footnote{Cf. the consideration at the end of the next paragraph.} ``Two forms are equivalent if and only if they differ by forms of higher folding, i.e. by forms that exhibit a greater weight decrease.'' This greater weight decrease always occurs when the two rows $q_{n-\sigma-i},p_{\tau-\lambda}$ entering the folding are really variable rows, as happens in (41.), (42.), in the foldings appearing in this paragraph as equivalent, and also in the equivalent forms of Theorem VII. But it need not occur when these rows are derived from symbol rows by substitution, as was the case for the forms of this paragraph which could be expressed in terms of one another by virtue of (62.). One easily sees that the vanishing forms even have, in part, a higher total weight.

\section*{\S\ 9. Form rows. Reduction theorems.}
We now define the form row introduced in \S\ 6 (end) somewhat more sharply as ``the totality of invariant formations of a given initial form that are linear in the coefficients, i.e. the totality of forms that arise from the initial form by folding in itself, arranged according to higher forms.''

To explain the higher forms, suppose by Theorem VII that the initial form is given as a product of two normal forms $S\cdot T$. By a higher form we mean forms with higher folding in itself, and among forms with the same folding in itself, specially normalized ones. More precisely: suppose the form $A$ has arisen by $\alpha_\rho$ fundamental foldings of type $\rho$, and the form $B$ by $\beta_\rho$ such foldings. Then $B$ is higher than $A$:
\begin{enumerate}
\item if, in the system of inequalities $\beta_\rho\ge\alpha_\rho$, $\rho=1,2,\ldots,n-1$, the inequality sign holds for at least one value of $\rho$;\footnote{If in the system of inequalities the sign $>$ holds in part and the sign $<$ in part, then the forms are not comparable, since a form arising from another by folding always has a positive value system of fundamental foldings, hence in our case positive or vanishing values $\beta_\rho-\alpha_\rho$.}
\item if the equality sign holds for all values $\rho$, but fewer symbol rows have been combined into a single matrix product. This normalization proves necessary because of the reductions in \S\ 8. Accordingly, for example, $\pair{S^{n-1}}{T_{n-1}}$ is higher than $\pair{S^{n-1}}{[S^1T^\rho]_{n-\rho-1}p_\rho}$;
\item if the equality sign holds for all values $\rho$ and the number of symbol rows combined into one matrix product is the same, then $B$ is made higher than $A$ by a normalization arbitrary in itself but fixed once and for all. This normalization is always possible because of the finite number of forms belonging to a value system $\alpha$; it can be achieved, for instance, by privileging the one form $S$ (larger number of symbol rows $S$, larger sum of the upper weight indices of the rows $S$, and so on).
\end{enumerate}
It is useful to think of the form row as arranged as an $(n-1)$-dimensional manifold of lattice points, where the $n-1$ directions correspond to the $n-1$ fundamental foldings. To every form there then corresponds one and only one value system $\alpha$, i.e. a positive integral lattice point, while the different forms belonging to a value system $\alpha$ are uniquely determined in their order by the above normalization. An arbitrary form of the form row gives the initial form of a new form row, contained as a part of the total form row; indeed it is contained as the part consisting of the forms obtained from the new initial form by proceeding in the $n-1$ positive directions.

By virtue of Theorem VIII and the general theory of series expansions, exactly the reduction theorems of the binary and ternary domains (cf. dissertation, \S\ 3) hold for form rows. We shall briefly derive the principal theorem. We first state the definition:

``In a given system of forms, suppose a certain finite number of forms has been combined into a module. We call a form reducible if it can be expressed by forms that have invariants as factors, or by `higher forms'; that is, by higher forms of the total form row to which the form belongs, or by forms that contain the symbols of the module in higher order. A reducible form row is called a reducer.'' Then:

\emph{Theorem IX:} If the initial form of a form row is reducible by having arisen through folding with a reducer, then the entire form row arising from this initial form is reducible.

For the proof, observe that a form $h$ arising by folding a form $f$ with a second form $g$ can, because of (45.), be regarded as a form $f$ containing several cogredient variable rows $p_\rho,p'_\rho$. Such forms can, by the general theory of series expansion, be represented as polars of the elementary covariants of $f$, by applying the series expansion successively to each pair of cogredient variable rows $p_\rho,p'_\rho$. By (38.), the elementary covariants that occur are the forms generated by cogredient folding. Since, however, the order in which the series expansion is applied is still arbitrary, we may in particular choose an order corresponding to the generation of the general foldings from fundamental foldings. The resulting system of elementary covariants is then identical with the form row $f$; the forms that arise by cogredient folding of higher defect are arranged among those that arise by fundamental foldings. But one also obtains polars of the form row $f$ for any order of the series expansion, to which a second system of elementary covariants may correspond. Since the form row exhausts the totality of invariant formations, every form of the second system can be expressed linearly by polars of the form row, and indeed by polars of those forms that show the same or a greater weight decrease. The latter follows because (cf. \S\ 7) the polars can be regarded as simultaneous invariants of a form row $H$ (52.), with the degree numbers corresponding to the form $h$, and of the form row $f$. Every folding of $H$ in itself corresponds to a weight decrease which, after substitution (11.) is performed, is transferred to the symbol rows and hence to the forms of $f$.\footnote{This consideration also underlies the second formulation of the equivalence concept (\S\ 8, end). Further explanations of the different systems of elementary covariants in the ternary domain are found in Study, loc. cit., II, \S\ 7, Theorems 7) and 8). By our Theorem VII the same results also hold for $n$ variables.}

An arbitrary form of the form row $h$ has arisen by folding $h$ in itself; that is, on the one hand by a higher folding of $g$ with $f$, and on the other by folding $f$ or $g$ in itself. In both cases the series expansion leads to polars of the form row $f$, just as was shown above for the initial form $h$. If in particular $f$ is a reducer, then an expansion in polars of the reducer results; hence the form row $h$ becomes reducible.

Applications of the theorem to the reduction of complete systems of forms follow analogously to the binary and ternary domains.

\end{document}
